\documentclass[aps,preprint]{revtex4}%
\usepackage{amsfonts}
\usepackage{amsmath}
\usepackage{amssymb}
\usepackage{graphicx}%
\setcounter{MaxMatrixCols}{30}
%TCIDATA{OutputFilter=latex2.dll}
%TCIDATA{Version=4.10.0.2363}
%TCIDATA{CSTFile=revtex4.cst}
%TCIDATA{Created=Friday, August 29, 2003 10:08:38}
%TCIDATA{LastRevised=Thursday, August 19, 2004 08:49:20}
%TCIDATA{<META NAME="GraphicsSave" CONTENT="32">}
%TCIDATA{<META NAME="DocumentShell" CONTENT="Articles\SW\REVTeX 4">}
%TCIDATA{Language=American English}
\newtheorem{theorem}{Theorem}
\newtheorem{acknowledgement}[theorem]{Acknowledgement}
\newtheorem{algorithm}[theorem]{Algorithm}
\newtheorem{axiom}[theorem]{Axiom}
\newtheorem{claim}[theorem]{Claim}
\newtheorem{conclusion}[theorem]{Conclusion}
\newtheorem{condition}[theorem]{Condition}
\newtheorem{conjecture}[theorem]{Conjecture}
\newtheorem{corollary}[theorem]{Corollary}
\newtheorem{criterion}[theorem]{Criterion}
\newtheorem{definition}[theorem]{Definition}
\newtheorem{example}[theorem]{Example}
\newtheorem{exercise}[theorem]{Exercise}
\newtheorem{lemma}[theorem]{Lemma}
\newtheorem{notation}[theorem]{Notation}
\newtheorem{problem}[theorem]{Problem}
\newtheorem{proposition}[theorem]{Proposition}
\newtheorem{remark}[theorem]{Remark}
\newtheorem{solution}[theorem]{Solution}
\newtheorem{summary}[theorem]{Summary}
\newenvironment{proof}[1][Proof]{\noindent\textbf{#1.} }{\ \rule{0.5em}{0.5em}}
\begin{document}
\preprint{ }
\title{Correlated One Particle Wave Functions}
\author{B. Weiner}
\affiliation{Department of Physics, Pennsylvania State University, DuBois PA 15801}
\author{J. V. Ortiz}
\affiliation{Department of Chemistry, Kansas State University, Manhattan, KS 66506-3701}
\pacs{}

\begin{abstract}


\end{abstract}
\date{08/11/2004}
\startpage{1}
\maketitle
\tableofcontents


\section{Introduction \label{intro}}

Independent Particle States (IPS's) have proved a mainstay of molecular
multi-particle quantum mechanics both in physical modelling and numerical
applications for many decades. The major attraction of these states lies in
the conceptually concrete and chemically understandable pictures of molecular
structure and properties that they provide. Apart from an initial choice of
$s$ atomic orbitals and possibly a decision about the molecular point group
and spin symmetry, one can obtain an unbiased description of the electronic
states of a molecule by finding an IPS\ that minimizes the energy via the Self
Consistent Field (SCF)\ equations. These equations provide a set of one
electron states $\left\{  \left.  \left\vert \varphi_{j}\right\rangle
\right\vert 1\leq j\leq2s\right\}  $, that are eigenstates of a self
consistent Fock operator corresponding to the eigenvalues $\left\{  \left.
\varepsilon_{j}\right\vert 1\leq j\leq2s\right\}  $ and are unique up to
unitary transformations that mix degenerate eigenstates. If the IPS\ describes
$M$ electrons then in general it is determined by the one-electron eigenstates
corresponding to the $M$ lowest eigenvalues $\left\{  \left.  \varepsilon
_{j}\right\vert 1\leq j\leq M\right\}  $. The wave functions of each of these
one electron states can then be analyzed in terms of a basis of atomic
orbitals, producing the familiar picture of one electron Probability
Distributions (PD's) labelled by one electron energies, that has been
invaluable in the understanding of chemical processes. The physical meaning of
the energies $\left\{  \left.  \varepsilon_{j}\right\vert 1\leq j\leq
M\right\}  $ can be deduced by considering the electron propagator
\cite{ortiz}, or by directly examining the form of the Fock operator matrix
elements, which show that they correspond to ionization potentials.

It is often the case, however, that the IPS\ model is not sufficiently
accurate to predict or explain the finer details of chemical reaction
mechanisms or spectroscopic observations. In order to achieve this one must
take into account electron-electron correlation. Traditionally this is
accomplished by perturbational or configurational interaction methods that
increase numerical accuracy at the cost of a clear, tangible, chemical model
of the molecular system. The clarity is lost as pictorial models need be based
on one particle properties, but in general one particle properties do not
determine correlated multi-electron states. In particular one cannot use
chemical intuition to rearrange the positions of atoms and bonds and be
certain that the new conformation actually corresponds to a multi-electron
state in contrast to the case of states determined by one particle properties.

However, we have described \cite{Weiner & Ortiz 2003} a class of correlated
states that generalize the IPS model, explicitly contain a description of
correlation effects, and are the \emph{most general} states that produce a
\emph{one particle} theory of electrons i.e. multi-particle states determined
by one particle properties. These are the Antisymmetrized Geminal Power (AGP)
states for an even number, $2N$, of electrons and the Generalized
Antisymmetrized Geminal Power (GAGP) ones for an odd number, $2N+1$. In the
IPS\ case the $2N$-electron state is completely determined by the occupation
numbers of the First Order Reduced Density Operator (FORDO) and a set of one
particle \emph{Natural }General Spin Orbitals (NGSOs). The occupation numbers
in this case are equal to one or zero. In the AGP\ case the $2N$-electron
state is completely determined by a set of one particle \emph{Canonical
}General Spin Orbitals (CGSOs) and occupation numbers of the FORDO, but now
the occupation numbers can be any value in the range $\left[  0,1\right]  ,$
though they must be at least doubly degenerate. This theory \emph{explicitly
contains} correlation and generalizes the one-particle theory based on
uncorrelated Independent Particle States (IPS), which it contains as a special
case. In this paper we further show that one can associate an ionization
potential with each CGSO and the orbitals with the largest occupation have the
largest ionization potential. The AGP\ model produces a simple pictorial
interpretation in terms of occupation numbers, orbitals and ionization
potentials, which totally \emph{determine} the $2N$-electron state, and thus
can be used as an effective and rigorous chemical modelling tool.

In contrast to IPS's when using AGP's it is important to note the distinction between

\begin{itemize}
\item AGP-NGSO's, which together with a set of doubly degenerate occupation
numbers, determine an AGP-FORDO, that correspond to a \emph{set }of
AGP\ states and

\item AGP-CGSO's, which together with a set of doubly degenerate occupation
numbers, determine a \emph{unique} AGP state.
\end{itemize}

The CGSO's are the NGSO's of the FORDO of a given AGP state, but the NGSO's of
the FORDO\ of this AGP state are \emph{not }in general the CSGO's of this
AGP\ state. In this paper when we are explicitly referring to CGSO's we will
use the symbol $\mathbf{\psi}$ and when we refer to NGSO's we will use the
symbol $\mathbf{\varphi.}$

The difference between NGSO's and CGSO's allows us to partition the manifold
of $2N$-electron AGP states into equivalence classes, $\left[  D^{1}\right]
_{2N}$, indexed by FORDO's $D^{1}$ that are at least doubly degenerate, where
each element of a class can be parameterized by a vector of angles
\begin{equation}
\mathbf{\theta\,}=\left(  0,\theta_{2},\cdots,\theta_{s}\right)
,\quad0<\theta_{j}\leq2\pi;\quad2\leq j\leq s,
\end{equation}
and a representative FORDO $D^{1}.$ As the occupation numbers $\left\{
\left.  N_{j}\right\vert \quad1\leq j\leq s\right\}  $ of the $2N$-AGP\ state
$\left\vert g^{N}\right\rangle $ are in $1-1$ correspondence \cite{Colemann}
with the occupation numbers $\left\{  n_{j}\quad1\leq j\leq s\right\}  $ of
the $2$-electron geminal state $\left\vert g\right\rangle $ each
representative $D^{1}$ can, in turn, be parameterized by a vector of geminal
occupation numbers%
\begin{equation}
\mathbf{n}=\left(  n_{1},\cdots,n_{s}\right)  ,\quad0\leq n_{j}\leq
1;\quad1\leq j\leq s;\quad%
%TCIMACRO{\tsum \limits_{1\leq j\leq s}}%
%BeginExpansion
{\textstyle\sum\limits_{1\leq j\leq s}}
%EndExpansion
n_{j}=1
\end{equation}
and a representative set of NGSO's $\left\vert \mathbf{\varphi}\right\rangle
=\left\{  \left.  \left\vert \varphi_{j}\right\rangle \right\vert 1\leq
j\leq2s\right\}  $.

The energy functional $\mathcal{E}\left(  g^{N}\right)  $ over the manifold of
AGP\ states can thus be formulated as a functional $\mathcal{E}\left(
\mathbf{n},\mathbf{\theta,\varphi}\right)  $ specifically describing the phase
and FORDO $D^{1}\left(  \mathbf{n},\mathbf{\varphi}\right)  $ dependency. This
leads to a Density Matrix Functional (DMF), $\mathcal{F}$, defined on the set
of equivalence classes $\left\{  \left[  D^{1}\right]  \right\}  $ by\
\begin{equation}
\mathcal{F}\left(  D^{1}\right)  =\min_{\mathbf{\theta}}\left\{
\mathcal{E}\left(  \mathbf{n},\mathbf{\theta,\varphi}\right)  \right\}  ,
\end{equation}
which can be utilized in Density Matrix Functional Theory (DMFT) as a concrete
example of a DMF.

If we concentrate on the occupation and CGSO dependency of the energy we
denote the functional by $\mathcal{E}\left(  \mathbf{n},\mathbf{\psi}\right)
$, where the geminal occupation numbers $\mathbf{n}$ must satisfy the
constraints
\begin{subequations}
\begin{align}
n_{j}  &  \geq0;\quad1\leq j\leq s,\\
\sum_{1\leq j\leq s}n_{j}  &  =1,
\end{align}
and the CGSO's the orthonormality constraints%
\end{subequations}
\begin{equation}
\left\langle \left.  \psi_{j}\right\vert \psi_{k}\right\rangle =\delta
_{jk};\quad1\leq j,k\leq2s.
\end{equation}
Variation of the energy with respect to (wrt) orbital changes described by a
Taylors series expansion involve functional derivatives of $\mathcal{E}\left(
\mathbf{n},\mathbf{\psi}\right)  $ with respect to $\left\{  \left.  \psi
_{j}\right\vert 1\leq j\leq2s\right\}  .$ We show that these derivatives can
be expressed in terms of a generalized Fock operator $F\left(  \mathbf{n}%
,\mathbf{\psi}\right)  $, which is \emph{only self adjoint} at stationary
points $\mathbf{\psi}_{s}$ of the orbital variation. Further if $\mathbf{n}%
_{s}$ is a stationary point of the occupation number variation then the
diagonal elements of the matrix of $F\left(  \mathbf{n}_{s},\mathbf{\psi}%
_{s}\right)  $ in the CGSO basis $\left\{  \left\vert \left\vert \psi
_{j}\right\rangle \right\vert 1\leq j\leq2s\right\}  $ are ionization
potentials associated with these orbitals weighted by their occupation numbers
$\left\{  N_{j}\right\}  $.

Unlike the IPS case the Fock operator is in general \emph{not} diagonal in the
basis $\left\{  \left.  \left\vert \psi_{j}\right\rangle \right\vert 1\leq
j\leq2s\right\}  $. It is, however, possible to diagonalize $F\left(
\mathbf{n}_{s},\mathbf{\psi}_{s}\right)  $ in terms of eigenvectors $\left\{
\left.  \left\vert \chi_{j}\right\rangle \right\vert 1\leq j\leq2s\right\}  $
and eigenvalues $\left\{  \left.  \varepsilon_{j}\right\vert 1\leq
j\leq2s\right\}  $ such that $\left\{  \varepsilon_{j}\right\}  $ are
ionization potentials associated with the orbitals $\left\{  \left\vert
\chi_{j}\right\rangle \right\}  $, but the geminal $\left\vert g\left(
\mathbf{n}_{s},\mathbf{\psi}_{s}\right)  \right\rangle ,$ that determines the
energy optimized AGP\ $\left\vert g\left(  \mathbf{n}_{s},\mathbf{\psi}%
_{s}\right)  ^{N}\right\rangle ,$ will \emph{not }be in a canonical form when
expressed in terms of the basis $\left\{  \left.  \left\vert \chi
_{j}\right\rangle \right\vert 1\leq j\leq2s\right\}  .$

One can avoid the use of constrained domains for the energy functional by
utilizing the unitary group of one electron Hilbert space $\mathcal{H}^{1},$
parameterized by $\left(  \mathbf{\lambda},\mathbf{\theta}\right)  ,$ to vary
the orbitals and a commutative non unitary group of positive transformations,
parameterized by $\mathbf{\eta},$ to vary the occupation numbers. This allows
us to express the energy as a function $\mathcal{E}\left(  \mathbf{\eta
},\mathbf{\lambda},\mathbf{\theta}\right)  $ where there are no constraints on
$\left(  \mathbf{\eta},\mathbf{\lambda},\mathbf{\theta}\right)  .$ Moreover
the unitary transformations produce NGSO's $\mathbf{\varphi=}$
$\mathbf{\varphi}\left(  \mathbf{\lambda}\right)  $ that belong to different
equivalence classes $\left[  D^{1}\right]  _{2N}$ and $\mathbf{\theta}$
produces transformations within an equivalence class by altering the phases of
$\mathbf{\varphi.}$

The one electron model that we obtain describes molecules by weighted PD's,
each determined by an occupation number and a phase optimized GSO associated
with an ionization potential, in an analogous fashion to the IPS case, but
also contains correlation effects. Pictorial manipulations of these PD's
(taking into account relaxation of the ionization energies and optimal phases)
allows one to analyze chemical processes and structures in a rigorously
meaningful fashion, as one is assured that the pictures produced corresponds
to multi electron states.

The structure of this article is as follows: in section \ref{AGPstate} we
describe the canonical form of a geminal, and the AGP\ state, in terms of
geminal occupation numbers $\left\{  \left.  n_{j}\right\vert 1\leq j\leq
s\right\}  $ and CGSOs $\left\{  \left.  \left\vert \psi_{j}\right\rangle
\right\vert 1\leq j\leq2s\right\}  .$ Together with the AGP-FORDO occupation
numbers $\left\{  \left.  N_{j}\right\vert 1\leq j\leq s\right\}  ,$ which are
in $1-1$ correspondence with the geminal occupation numbers, these CGSOs
determine both the AGP\ state and the AGP-FORDO. In section \ref{transform} we
examine the $1-1$ transformations between the AGP $\left\{  \left.
N_{j}\right\vert 1\leq j\leq s\right\}  $ and geminal $\left\{  \left.
n_{j}\right\vert 1\leq j\leq s\right\}  $ occupation numbers. In section
\ref{Energy} we discuss the energy expressions in terms of the occupation
numbers $\left\{  \left.  n_{j}\right\vert 1\leq j\leq s\right\}  $, CGSOs
$\left\{  \left.  \left\vert \psi_{j}\right\rangle \right\vert 1\leq
j\leq2s\right\}  $,NGSOs $\left\{  \left.  \left\vert \varphi_{j}\right\rangle
\right\vert 1\leq j\leq2s\right\}  $ and phases $\left\{  \left.  \theta
_{j}\right\vert 1\leq j\leq s\right\}  .$ In section \ref{Evariation} we
describe the necessary conditions needed to minimize the energy over a
manifold of AGP\ states. In section \ref{Fock} we derive a generalized Fock
operator that produces the derivative of the energy functional with respect to
NGSO's for fixed FORDO occupation numbers. In section \ref{aufbau} we analyze
the total energy in terms of the eigenvalues of this Fock operator and show
that an aufbau principle similar to the IPS\ case is effective in constructing
multiparticle states. In \ref{procedure} we form a combined procedure that
produces an energy optimized AGP state and we conclude in section
\ref{Discussion} with a review of the material presented and a discussion of
its implications.

\section{The $2N$-electron AGP state\label{AGPstate}}

In this section we review the canonical form of a geminal and AGP state and
their associated first and second order reduced density operators. We also
characterize equivalence classes of AGP states that have the same FORDO.

\subsection{Canonical Form}

The $2N$-electron AGP state, $\left\vert g^{N}\right\rangle $, is the $N^{th}$
antisymmetric power of the $2$-electron state described by the geminal
\begin{equation}
\left\vert g\right\rangle =\sum_{1\leq j\leq s}n_{j}^{\frac{1}{2}}\left\vert
\psi_{j}\psi_{j+s}\right\rangle , \label{geminal}%
\end{equation}
where $\left\{  \left.  n_{j}^{\frac{1}{2}}\right\vert 1\leq j\leq s\right\}
$ are the positive square roots of $\left\{  \left.  n_{j}\right\vert 1\leq
j\leq s\right\}  ,$ and is given by \cite{weiner&ortiz2004}
\begin{equation}
\left\vert g^{N}\right\rangle =\left(  S_{N}\right)  ^{-\frac{1}{2}}%
\sum_{1\leq j_{1}<\cdots<j_{N}\leq s}n_{j_{1}}^{\frac{1}{2}}\cdots n_{j_{s}%
}^{\frac{1}{2}}\left\vert \psi_{j_{1}}\psi_{j_{1}+s}\cdots\psi_{j_{N}}%
\psi_{j_{N}+s}\right\rangle . \label{agp}%
\end{equation}
The normalization factor, $\left(  S_{N}\right)  ^{-\frac{1}{2}}$, is given in
terms of the elementary symmetric function
\begin{equation}
S_{N}=\sum_{1\leq j_{1}<\cdots<j_{N}\leq s}n_{j_{1}}\cdots n_{j_{N}},
\end{equation}
and the geminals $\left\vert g\right\rangle $ represent normalized two
electron states so that
\begin{equation}
\sum_{1\leq j\leq s}n_{j}=1.
\end{equation}
It is evident from the expansion eq. $\left(  \ref{geminal}\right)  $ that the
geminal and geminal power state are at the least invariant to the group
\begin{equation}
SU\left(  2\right)  ^{s}=\overset{s\text{-factors}}{\overbrace{\,SU\left(
2\right)  \times\cdots\times SU\left(  2\right)  }},
\end{equation}
that is formed by $2\times2$ \emph{special} unitary transformations of the
paired CGSO's $\left\{  \left.  \left\vert \psi_{j}\right\rangle ,\left\vert
\psi_{j+s}\right\rangle \right\vert 1\leq j\leq s\right\}  ,$ while if
degeneracies exist amongst the $\left\{  \left.  n_{j}\right\vert 1\leq j\leq
s\right\}  $ this invariance group is larger \cite{Weiner2004}. The canonical
CSGO's are unique up to transformations belonging this invariance group.

The CGSO's $\left\{  \left.  \left\vert \psi_{j}\right\rangle \right\vert
1\leq j\leq2s\right\}  $ can be expanded in terms of a spin orbitals basis
$\left\{  \left.  \left\vert \varsigma_{j\alpha}\right\rangle ,\left\vert
\varsigma_{j\beta}\right\rangle \right\vert 1\leq j\leq s\right\}  $ of one
electron state space $\mathcal{H}^{1}$ as
\begin{equation}
\left\vert \psi_{j}\right\rangle =\sum_{1\leq j\leq s}\left\{  a_{j}\left\vert
\varsigma_{j\alpha}\right\rangle +b_{j}\left\vert \varsigma_{j\beta
}\right\rangle \right\}  ,
\end{equation}
leading to one electron GSO wavefunctions defined on the spectrum of position
and spin operators as
\begin{equation}
\psi_{j}\left(  \mathbf{r,}\zeta\right)  =\left\langle \left.  \mathbf{r,}%
\zeta\right\vert \psi_{j}\right\rangle .
\end{equation}
It should be noted that if expansion eq. $\left(  \ref{geminal}\right)  $ is
used to define the canonical expansion it is usually \emph{not possible} to
find a set of CSGO's such that
\begin{equation}
\psi_{j+s}\left(  \mathbf{r,}\zeta\right)  =\psi_{j}\left(  \mathbf{r,}%
\zeta\right)  ^{\ast}.
\end{equation}


\subsection{First Order Reduced Density Operators\label{FORDO}}

The states $\left\vert g\right\rangle $ and $\left\vert g^{N}\right\rangle $
define FORDO's given by \cite{Coleman}
\begin{equation}
D^{1}\left(  g\right)  =\sum_{1\leq j\leq s}n_{j}\left\{  \left\vert \psi
_{j}\right\rangle \left\langle \psi_{j}\right\vert +\left\vert \psi
_{j+s}\right\rangle \left\langle \psi_{j+s}\right\vert \right\}  \label{d1g}%
\end{equation}
and
\begin{equation}
D^{1}\left(  g^{N}\right)  =\sum_{1\leq j\leq s}N_{j}\left\{  \left\vert
\psi_{j}\right\rangle \left\langle \psi_{j}\right\vert +\left\vert \psi
_{j+s}\right\rangle \left\langle \psi_{j+s}\right\vert \right\}  \label{dng}%
\end{equation}
respectively. It is clear from the expansions eqs. $\left(  \ref{d1g}\right)
$ and $\left(  \ref{dng}\right)  $ that the CGSO's $\left\{  \left.
\left\vert \psi_{j}\right\rangle \right\vert 1\leq j\leq r\right\}  $ are
simultaneously the Natural General Spin Orbitals (NGSO's) of $D^{1}\left(
g\right)  $ and $D^{1}\left(  g^{N}\right)  $. However it is very important to
note that the NGSO's of $D^{1}\left(  g\right)  $ and $D^{1}\left(
g^{N}\right)  $ \emph{are not }in general the CGSO's of $\left\vert
g\right\rangle $. This is a ramification of

\begin{enumerate}
\item the general invariance of FORDO's to the multiplication of their NGSO's
by complex phase factors
\begin{equation}
D^{1}=\sum_{1\leq j\leq r}\nu_{j}\left\{  \left\vert \chi_{j}\right\rangle
\left\langle \chi_{j}\right\vert \right\}  =\sum_{1\leq j\leq r}\nu
_{j}\left\{  e^{i\theta_{j}}\left\vert \chi_{j}\right\rangle \left\langle
\chi_{j}\right\vert e^{-i\theta_{j}}\right\}  ,
\end{equation}
where $\left\{  \left.  \nu_{j}\right\vert 1\leq j\leq s\right\}  $ are
occupation numbers and

\item the observation that geminals \emph{are not }invariant to such phase
changes of their CGSO's.
\end{enumerate}

One can form equivalence classes $\left[  D^{1}\left(  g\left(  \mathbf{0}%
\right)  \right)  \right]  $ of geminals $\left\{  \left\vert g\left(
\mathbf{\theta}\right)  \right\rangle \right\}  $ that have the same FORDO
$D^{1}\left(  g\left(  \mathbf{0}\right)  \right)  $ and equivalence classes
$\left[  D^{1}\left(  g^{N}\left(  \mathbf{0}\right)  \right)  \right]  $ of
$2N$-electron AGP states $\left\{  \left\vert g^{N}\left(  \mathbf{\theta
}\right)  \right\rangle \right\}  $, that have the same FORDO $D^{1}\left(
g^{N}\left(  \mathbf{0}\right)  \right)  $. For a fixed $N$ these classes
correspond to each other in a bijective manner \cite{coleman} and can be
parameterized by the angles $\mathbf{\theta}$, where
\begin{align}
\left\vert g\left(  \mathbf{\theta}\right)  \right\rangle  &  =\sum_{1\leq
j\leq s}n_{j}^{\frac{1}{2}}e^{i\theta_{j}}\left\vert \varphi_{j}\varphi
_{j+s}\right\rangle =\sum_{1\leq j\leq s}n_{j}^{\frac{1}{2}}\left\vert
e^{i\alpha_{j}}\varphi_{j}e^{i\alpha_{j+s}}\varphi_{j+s}\right\rangle
;\nonumber\\
&  \qquad\qquad\qquad\qquad\qquad\qquad\qquad\theta_{1}=0,\,0<\theta_{j}%
\leq2\pi,\,2\leq j\leq s,
\end{align}
where $\mathbf{\theta=}\left(  0,\theta_{2},\cdots,\theta_{s}\right)  $ and
$\alpha_{j}+\alpha_{j+s}=\theta_{j}$. Choosing $\theta_{1}=0$ eliminates the
overall phase factor indeterminacy of $\left\vert g\left(  \mathbf{\theta
}\right)  \right\rangle $. The individual phase factors $\left\{  \left.
\alpha_{j}\right\vert 1\leq j\leq s\right\}  $ are not unique, though their
sums $\left\{  \left.  \theta_{j}\right\vert 1\leq j\leq s\right\}  $ are,
allowing one to obtain a unique representative set $\left\{  \left.
\alpha_{j}\right\vert 1\leq j\leq s\right\}  $ by defining
\begin{equation}
\alpha_{j}=\frac{\theta_{j}}{2},\alpha_{j+s}=\frac{\theta_{j}}{2};\quad1\leq
j\leq s.
\end{equation}
Thus one has the correspondence
\begin{equation}
\left\{  \left\vert g\left(  \mathbf{\theta}\right)  \right\rangle
\longleftrightarrow\left\vert g^{N}\left(  \mathbf{\theta}\right)
\right\rangle ;\quad\theta_{1}=0;\,0<\theta_{j}\leq2\pi,2\leq j\leq s\right\}
\longleftrightarrow\left[  D^{1}\left(  g^{N}\left(  \mathbf{0}\right)
\right)  \right]  .
\end{equation}
Each of the elements of the class $\left[  D^{1}\left(  g^{N}\left(
\mathbf{0}\right)  \right)  \right]  $ can also be expressed in canonical form
as
\begin{equation}
\sum_{1\leq j\leq s}n_{j}^{\frac{1}{2}}\left\vert \chi_{j}\chi_{j+s}%
\right\rangle ,
\end{equation}
where $\left\vert \chi_{j}\right\rangle =\left\vert e^{i\varphi_{j}}%
\varphi_{j}\right\rangle $.

\subsection{Second Order Reduced Density Operators\label{SORDO}}

The AGP\ states $\left\vert g^{N}\left(  \mathbf{\theta}\right)  \right\rangle
$ are clearly distinct as they have different Second Order Reduced Density
Operators (SORDO's), $D^{2}\left(  g^{N}\left(  \mathbf{\theta}\right)
\right)  $, given by \cite{Coleman}
\begin{align}
&  D^{2}\left(  g^{N}\left(  \mathbf{\theta}\right)  \right)  =\frac{1}{S_{N}%
}\left\{  \sum_{1\leq j\leq s}n_{j}\left\vert \varphi_{j}\varphi
_{j+s}\right\rangle \left\langle \varphi_{j}\varphi_{j+s}\right\vert \right.
+\nonumber\\
&  \left.  \sum_{1\leq j,k\leq s}n_{j}^{\frac{1}{2}}n_{k}^{\frac{1}{2}%
}e^{i\left(  \theta_{j}-\theta_{k}\right)  }S_{N-1}\left(  \jmath k\right)
\left\vert \varphi_{j}\varphi_{j+s}\right\rangle \left\langle \varphi
_{k}\varphi_{k+s}\right\vert +\sum_{1\leq j<k\leq2s}n_{j}n_{k}S_{N-2}\left(
\jmath k\right)  \left\vert \varphi_{j}\varphi_{k}\right\rangle \left\langle
\varphi_{j}\varphi_{k}\right\vert \right\}  , \label{D2}%
\end{align}
where
\begin{equation}
S_{N-1}\left(  \jmath k\right)  =\sum_{\substack{1\leq j_{1}<\cdots
<j_{N-1}\leq s\\j,k\notin\left\{  j_{1},\cdots,j_{N-1}\right\}  }}n_{j_{1}%
}\cdots n_{j_{N-1}}%
\end{equation}
and
\begin{equation}
S_{N-2}\left(  \jmath k\right)  =\sum_{\substack{1\leq j_{1}<\cdots
<j_{N-2}\leq s\\j,k\notin\left\{  j_{1},\cdots,j_{N-2}\right\}  }}n_{j_{1}%
}\cdots n_{j_{N-2}}.
\end{equation}
Note that we define
\begin{align}
S_{N-1}\left(  \jmath k+s\right)   &  =S_{N-1}\left(  \jmath+sk\right)
=S_{N-1}\left(  \jmath+sk+s\right)  =S_{N-1}\left(  kj\right)  =S_{N-1}\left(
\jmath k\right) \nonumber\\
S_{N-2}\left(  \jmath k+s\right)   &  =S_{N-2}\left(  \jmath+sk\right)
=S_{N-2}\left(  \jmath+sk+s\right)  =S_{N-2}\left(  kj\right)  =S_{N-2}\left(
\jmath k\right) \nonumber\\
S_{N-1}\left(  \jmath\jmath\right)   &  =S_{N-2}\left(  \jmath\jmath\right)
=0.
\end{align}
The SORDO eq. $\left(  \ref{D2}\right)  $ can be expanded in terms of
$\mathbf{\theta}$ dependent NGSO's $\left\{  \left.  \left\vert e^{i\alpha
_{j}}\varphi_{j}\right\rangle ,\left\vert e^{i\alpha_{j+s}}\varphi
_{j+s}\right\rangle \right\vert 1\leq j\leq s\right\}  ,$ where $\left\{
\left.  \alpha_{j}+\alpha_{j+s}=\theta_{j}\right\vert 1\leq j\leq s\right\}  $
giving the explicitly orbital phase dependent expression
\begin{align}
&  D^{2}\left(  g^{N}\left(  \mathbf{\theta}\right)  \right)  =\frac{1}{S_{N}%
}\left\{  \sum_{1\leq j\leq s}n_{j}S_{N-1}\left(  \jmath\right)  \left\vert
e^{i\alpha_{j}}\varphi_{j}e^{i\alpha_{j+s}}\varphi_{j+s}\right\rangle
\left\langle e^{i\alpha_{j}}\varphi_{j}e^{i\alpha_{j+s}}\varphi_{j+s}%
\right\vert \right. \nonumber\\
&  \qquad\qquad\qquad\qquad+\sum_{1\leq j,k\leq s}n_{j}^{\frac{1}{2}}%
n_{k}^{\frac{1}{2}}S_{N-1}\left(  \jmath k\right)  \left\vert e^{i\alpha_{j}%
}\varphi_{j}e^{i\alpha_{j+s}}\varphi_{j+s}\right\rangle \left\langle
e^{i\alpha_{k}}\varphi_{k}e^{i\alpha_{k+s}}\varphi_{k+s}\right\vert
\nonumber\\
&  \qquad\qquad\qquad\qquad\left.  +\sum_{1\leq j<k\leq2s}n_{j}n_{k}%
S_{N-2}\left(  \jmath k\right)  \left\vert e^{i\alpha_{j}}\varphi
_{j}e^{i\alpha_{k}}\varphi_{k}\right\rangle \left\langle e^{i\alpha_{j}%
}\varphi_{j}e^{i\alpha_{k}}\varphi_{k}\right\vert \right\}  . \label{en2}%
\end{align}
In terms of CGSO's $\left\{  \psi_{j}\right\}  $ the SORDO can be expressed in
phase free form as
\begin{align}
&  D^{2}\left(  g^{N}\right)  =\frac{1}{S_{N}}\left\{  \sum_{1\leq j\leq
s}n_{j}\left\vert \psi_{j}\psi_{j+s}\right\rangle \left\langle \psi_{j}%
\psi_{j+s}\right\vert \right.  +\nonumber\\
&  \left.  \sum_{1\leq j,k\leq s}n_{j}^{\frac{1}{2}}n_{k}^{\frac{1}{2}}%
S_{N-1}\left(  \jmath k\right)  \left\vert \psi_{j}\psi_{j+s}\right\rangle
\left\langle \psi_{k}\psi_{k+s}\right\vert +\sum_{1\leq j<k\leq2s}n_{j}%
n_{k}S_{N-2}\left(  \jmath k\right)  \left\vert \psi_{j}\psi_{k}\right\rangle
\left\langle \psi_{j}\psi_{k}\right\vert \right\}  . \label{D2C}%
\end{align}


\section{The correspondence between AGP and Geminal Occupation
Numbers\label{transform}}

Any set of $s$ positive numbers $\left\{  \left.  N_{j}\right\vert 1\leq j\leq
s\right\}  $ that have values between $0$ and $1$ and add up to $N,$ can be
interpreted as the occupation numbers of a FORDO that is at least doubly
degenerate and can be associated with an AGP\ state. These numbers are in
$1-1$ correspondence with another set of positive numbers $\left\{  \left.
n_{j}\right\vert 1\leq j\leq s\right\}  ,$ that add up to $1$, which can be
interpreted as the occupation numbers of a FORDO\ of a geminal. The geminal
occupation numbers $\left\{  \left.  n_{j}\right\vert 1\leq j\leq s\right\}  $
together with an arbitrary set of $2s$ orthonormal one electron Canonical
General Spin Orbitals (CGSO's) then \emph{uniquely }determine a $2N$ electron
AGP state. Any state that has a FORDO that has at least doubly degenerate
occupation numbers is covered by a set of AGP\ states i.e. there exists
AGP\ states that have the same FORDO. In most applications e.g. expectation
values of observables, it is sufficient to know the $2N$ electron state
occupation numbers in terms of the geminal occupation numbers, which serve as
a parameterization of the state occupation numbers, and we review this
relationship first. The inverse function i.e. geminal occupation numbers in
terms of state occupation numbers are significant in designing possible
AGP\ states and we present a numerical construction of this function.

\subsection{Transformation of Geminal to AGP Occupation Numbers}

The vector of occupation numbers $\mathbf{N}=\left(  N_{1},\cdots
,N_{s}\right)  $ of an AGP\ state $\left\vert g^{N}\right\rangle $ can be
expressed as a bijective function, $\mathcal{N}\left(  \mathbf{n}\right)  $,
of the geminal occupation numbers $\mathbf{n}=\left(  n_{1},\cdots
,n_{s}\right)  ,$ where
\begin{equation}
\mathcal{N}:\mathfrak{D}_{N}\subset\mathbb{R}^{s}\rightarrow\mathfrak{D}%
_{n}\subset\mathbb{R}^{s}.
\end{equation}
The domain $\mathfrak{D}_{N}$ and range $\mathfrak{D}_{n}$ are given by%
\begin{align}
\mathfrak{D}_{N}  &  =\left(  N_{1},\cdots,N_{s}\right)  ;\quad0\leq N_{j}%
\leq1,\quad1\leq j\leq s\nonumber\\
\sum_{1\leq j\leq s}N_{j}  &  =N
\end{align}
and
\begin{align}
\mathfrak{D}_{n}  &  =\left(  n_{1},\cdots,n_{s}\right)  ;\quad n_{j}%
\geq0,\quad1\leq j\leq s\nonumber\\
\sum_{1\leq j\leq s}n_{j}  &  =1,
\end{align}
where
\begin{equation}
N_{j}=\mathcal{N}_{j}\left(  \mathbf{n}\right)  =n_{j}\frac{S_{N-1}\left(
\jmath\right)  }{S_{N}};\quad1\leq j\leq s
\end{equation}
and
\begin{equation}
S_{N-1}\left(  \jmath\right)  =\sum_{\substack{1\leq j_{1}<\cdots<j_{N-1}\leq
s\\j\notin\left\{  j_{1},\cdots,j_{N-1}\right\}  }}n_{j_{1}}\cdots n_{j_{N-1}%
}.
\end{equation}


\subsection{Transformation of AGP to Geminal Occupation
Numbers\label{AGPoccGemocc}}

The inverse function $\mathcal{N}^{-1}\left(  \mathbf{N}\right)  $ is not
explicitly known, but one can obtain a linear approximation that suffices in
practical applications to numerical determine values of $\mathbf{n}$ from
known values of $\mathbf{N}$, i.e. the determination of $\mathbf{n}_{0}%
+\delta\mathbf{n}$ from $\mathbf{N}_{0}+\delta\mathbf{N}$, when $\mathbf{N}%
_{0}$ is a known initial vector of occupation numbers of the $2N$-electron AGP
state, e.g. an IPS,\ where
\begin{equation}
\mathbf{N}_{0}=\left(  \overset{N}{\overbrace{1,\cdots,1}},\overset
{s-N}{\overbrace{0,\cdots,0}}\right)
\end{equation}
and
\begin{equation}
\mathbf{n}_{0}=\left(  \overset{N}{\overbrace{\frac{1}{N},\cdots,\frac{1}{N}}%
},\overset{s-N}{\overbrace{0,\cdots,0}}\right)  .
\end{equation}
The Taylor series expansion about $\mathbf{N}_{0}$ of $\mathcal{N}^{-1}$ is
given by%
\begin{equation}
\mathcal{N}^{-1}\left(  \mathbf{N}_{0}+\delta\mathbf{N}\right)  =\mathcal{N}%
^{-1}\left(  \mathbf{N}_{0}\right)  +d^{1}\mathcal{N}^{-1}\left(
\mathbf{N}_{0}\right)  \left(  \delta\mathbf{n}\right)  +\cdots,
\end{equation}
where
\begin{equation}
d^{1}\mathcal{N}^{-1}\left(  \mathbf{N}_{0}\right)  =\left[  d^{1}%
\mathcal{N}\left(  \mathbf{n}_{0}\right)  \right]  ^{-1}%
\end{equation}
and the Jacobian, $d^{1}\mathcal{N}\left(  \mathbf{n}_{0}\right)  $, of the
transformation $\mathcal{N}$ at the point $\mathbf{n}_{0}$ is given by the
matrix elements
\begin{align}
\left[  d^{1}\mathcal{N}\left(  \mathbf{n}_{0}\right)  \right]  _{jk}  &
=\frac{\partial\mathcal{N}_{j}\left(  \mathbf{n}_{0}\right)  }{\partial n_{k}%
}=\frac{\partial}{\partial n_{k}}\left\{  n_{j}\frac{S_{N-1}\left(
\jmath\right)  }{S_{N}}\right\} \nonumber\\
&  =\frac{1}{S_{N}}\left(  \delta_{jk}S_{N-1}\left(  \jmath\right)
+n_{j}\frac{\partial S_{N-1}\left(  \jmath\right)  }{\partial n_{k}}%
-n_{j}\frac{S_{N-1}\left(  \jmath\right)  }{S_{N}}\frac{\partial S_{N}%
}{\partial n_{k}}\right) \nonumber\\
&  =\frac{1}{S_{N}}\left(  \delta_{jk}S_{N-1}\left(  \jmath\right)
+n_{j}S_{N-2}\left(  \jmath k\right)  -n_{j}\frac{1}{S_{N}}S_{N-1}\left(
\jmath\right)  S_{N-1}\left(  k\right)  \right)  .
\end{align}
The Jacobian of the inverse transformation $\mathcal{N}^{-1}\left(
\mathbf{N}_{0}\right)  $ at the point $\mathbf{N}_{0}$ is given by the inverse
of the matrix $\left\{  \left.  \left[  d^{1}\mathcal{N}\left(  \mathbf{n}%
_{0}\right)  \right]  _{jk}\right\vert 1\leq j,k\leq s\right\}  $ i.e.
\begin{equation}
d^{1}\mathcal{N}^{-1}\left(  \mathbf{N}_{0}\right)  =\left[  d^{1}%
\mathcal{N}\left(  \mathbf{n}_{0}\right)  \right]  ^{-1}%
\end{equation}
and the incremental change, $\delta\mathbf{n}$, of $\mathbf{n}_{0}$ is
\begin{equation}
d^{1}\mathcal{N}^{-1}\left(  \mathbf{N}_{0}\right)  \left(  \delta
\mathbf{N}\right)  _{j}=\sum_{1\leq k\leq s}\left[  d^{1}\mathcal{N}\left(
\mathbf{n}_{0}\right)  \right]  _{jk}^{-1}\left(  \delta\mathbf{N}\right)
_{k}.
\end{equation}
We can then construct a piecewise linear fit to $\mathcal{N}^{-1}$ as%
\begin{equation}
\mathbf{n}\simeq\mathbf{\tilde{n}}_{m}=\mathcal{N}^{-1}\left(  \mathbf{N}%
\right)  =\mathbf{n}_{0}+\sum_{1\leq k\leq m}\left[  d^{1}\mathcal{N}\left(
\mathbf{n}_{0}+\sum_{1\leq j\leq k-1}\left[  d^{1}\mathcal{N}\left(
\mathbf{n}_{0}\right)  \right]  ^{-1}\left(  \delta\mathbf{N}_{j}\right)
\right)  \right]  ^{-1}\left(  \delta\mathbf{N}_{k}\right)  ,
\end{equation}
where
\begin{equation}
\mathbf{N=N}_{0}+\sum_{1\leq j\leq k-1}\delta\mathbf{N}_{j}%
\end{equation}
and
\begin{equation}
\sum_{1\leq k\leq s}\delta\mathbf{N}_{jk}=0.
\end{equation}
See Appendix \label{linfit} for an example. The accuracy of the fit can be
checked by evaluating
\begin{equation}
\mathbf{\Delta}_{m}=\mathcal{N}\left(  \mathbf{\tilde{n}}_{m}\right)
-\mathbf{N}_{m}%
\end{equation}
and noting if
\begin{equation}
\left\vert \mathbf{\Delta}_{mj}\right\vert \leq\varkappa_{j};\quad1\leq j\leq
s \label{con}%
\end{equation}
If the error $\left(  \left\vert \mathbf{\Delta}_{m1}\right\vert
,\cdots,\left\vert \mathbf{\Delta}_{ms}\right\vert \right)  $, i.e. is too big
i.e. eq. $\left(  \ref{con}\right)  $ is not satisfied one can then decrease
the size of $\delta\mathbf{N}_{m},$ etc.

\section{Energy Functionals \label{Energy}}

In this section we describe the energy of a $2N$-electron AGP state in three
different ways in terms of:

\begin{description}
\item
\begin{enumerate}
\item CGSO's and geminal occupation numbers.

\item A density matrix functional $\mathcal{F}\left(  D^{1}\right)  $ defined
on the equivalence classes $\left[  D^{1}\left(  g^{N}\left(  \mathbf{0}%
\right)  \right)  \right]  .$ The functional $\mathcal{F}\left(  D^{1}\right)
$ is significant in two aspects, it can serve as a rigorous explicitly known
model functional in Density Matrix Functional Theory (DMFT) \cite{DMFT} and it
clearly demonstrates the difference between CGSO's and NGSO's.

\item Group parameters discussed in \cite{weiner2004} that describe the
manifold of AGP\ states. This parameterization allows one to constructively
separate the variation of the phases $\left\{  \theta_{j}\right\}  $, NGSO's
$\left\{  \left\vert \varphi_{j}\right\rangle \right\}  $ and occupation
numbers $\left\{  n_{j}\right\}  $, which is necessary in order to construct
the functionals $\mathcal{F}\left(  D^{1}\right)  $. It is also important in
the analysis of the properties of the generalized Fock operator associated
with optimized AGP\ states and the derivation of a generalized Brillouin condition.
\end{enumerate}
\end{description}

The supporting details are in Appendix \ref{En_Funcs}.

\subsection{Energy in terms of Geminal Occupation Numbers and
CGSO's\label{EnON_CGSO}}

The energy, $\mathcal{E}\left(  \mathbf{n,\psi}\right)  ,$ of a $2N$-electron
AGP\ state for a molecule at a fixed nuclear configuration $\left(
\mathbf{R}_{1},\cdots,\mathbf{R}_{N_{nuc}}\right)  $ in terms of the CGSO's
amplitudes $\left\{  \left.  \left\langle \left.  \mathbf{r,\xi}\right\vert
\psi_{j}\right\rangle \right\vert 1\leq j\leq r\right\}  $ and geminal
occupation numbers $\mathbf{n}$ is
\begin{align}
\mathcal{E}\left(  \mathbf{n,\psi}\right)   &  =\frac{\mathcal{H}\left(
\mathbf{n,\psi}\right)  }{S_{N}\left(  \mathbf{n}\right)  }\nonumber\\
&  =\frac{1}{S_{N}\left(  \mathbf{n}\right)  }\left\{  \sum_{1\leq j\leq
s}n_{j}S_{N-1}\left(  \jmath\right)  \left\{  \left\langle \psi_{j}\left\vert
h_{1}\psi_{j}\right.  \right\rangle +\left\langle \psi_{j+s}\left\vert
h_{1}\psi_{j+s}\right.  \right\rangle +\left\langle \psi_{j}\psi_{j+s}%
||\psi_{j}\psi_{j+s}\right\rangle \right\}  \right. \nonumber\\
&  +\sum_{1\leq j,k\leq s}n_{j}^{\frac{1}{2}}n_{k}^{\frac{1}{2}}S_{N-1}\left(
\jmath k\right)  \left\langle \psi_{j}\psi_{j+s}||\psi_{k}\psi_{k+s}%
\right\rangle \left.  +\sum_{1\leq j<k\leq s}n_{i}n_{j}S_{N-2}\left(  \jmath
k\right)  \left\langle \psi_{j}\psi_{k}|||\psi_{j}\psi_{k}\right\rangle
\right\}  . \label{energy}%
\end{align}
$\lceil$Note: Mathematically the energy functional should really be denoted as
$\mathcal{E}\left(  \mathbf{n,\psi}^{\ast}\mathbf{,\psi}\right)  $ as
$\mathbf{\psi}^{\ast}$ and $\mathbf{\psi}$ are treated as independent
variables even though $\mathbf{\psi}$ determines $\mathbf{\psi}^{\ast}%
$.$\rfloor$ The one-electron matrix elements $\left\{  \left\langle \psi
_{j}\left\vert h_{1}\psi_{k}\right.  \right\rangle \right\}  $ wrt the CGSO's
are
\begin{equation}
\left\langle \psi_{j}\left\vert h_{1}\psi_{k}\right.  \right\rangle =-\int
\psi_{j}^{\ast}\left(  \mathbf{r,\xi}\right)  \left\{  \sum_{_{1\leq l\leq
N_{nuc}}}\frac{Z_{l}}{\left\Vert \mathbf{R}_{l}-\mathbf{r}\right\Vert }%
+\nabla_{\mathbf{r}}^{2}\right\}  \psi_{k}\left(  \mathbf{r,\xi}\right)
d\mathbf{r}d\xi;\quad1\leq j,k\leq2s, \label{h1}%
\end{equation}
the antisymmetric two-electron matrix elements $\left\{  \left\langle
jk||lm\right\rangle \right\}  $
\begin{align}
\left\langle jk||lm\right\rangle  &  =\int\left\{  \psi_{j}^{\ast}\left(
\mathbf{r}_{1}\mathbf{,\xi}_{1}\right)  \psi_{k}^{\ast}\left(  \mathbf{r}%
_{2}\mathbf{,\xi}_{2}\right)  \frac{1}{\left\Vert \mathbf{r}_{1}%
-\mathbf{r}_{2}\right\Vert }\psi_{l}\left(  \mathbf{r}_{1}\mathbf{,\xi}%
_{1}\right)  \psi_{m}\left(  \mathbf{r}_{2}\mathbf{,\xi}_{2}\right)  \right.
\nonumber\\
&  -\left.  \psi_{j}^{\ast}\left(  \mathbf{r}_{1}\mathbf{,\xi}_{1}\right)
\psi_{k}^{\ast}\left(  \mathbf{r}_{2}\mathbf{,\xi}_{2}\right)  \frac
{1}{\left\Vert \mathbf{r}_{1}-\mathbf{r}_{2}\right\Vert }\psi_{l}\left(
\mathbf{r}_{2}\mathbf{,\xi}_{1}\right)  \psi_{m}\left(  \mathbf{r}%
_{1}\mathbf{,\xi}_{2}\right)  \right\}  d\mathbf{r}_{1}d\mathbf{r}_{2}d\xi
_{1}d\xi_{2}\nonumber\\
&  \left.  {}\right.  \qquad\qquad\qquad\qquad\qquad\qquad\qquad\qquad
\qquad1\leq j,k\leq2s;1\leq l,m\leq2s \label{h2}%
\end{align}
and $\left\{  \left\langle jk|||jk\right\rangle \right\}  $
\begin{align}
\left\langle \psi_{j}\psi_{k}|||\psi_{j}\psi_{k}\right\rangle  &
=\left\langle jk||jk\right\rangle +\left\langle j\,k+s||j\,k+s\right\rangle
+\left\langle j+s\,k||j+s\,k\right\rangle +\left\langle
j+s\,k+s||j+s\,k+s\right\rangle \nonumber\\
&  \left.  {}\right.  \qquad\qquad\qquad\qquad\qquad\qquad\qquad\qquad
\qquad\qquad\qquad\qquad\qquad\qquad1\leq j,k\leq s.
\end{align}
We shall often express $\left\langle jk||lm\right\rangle $ in terms of non
antisymmetric matrix elements $\left\{  \left(  \left.  jk\right\vert
h_{2}lm\right)  \right\}  $ as
\begin{equation}
\left\langle jk||lm\right\rangle =\left(  \left.  jk\right\vert h_{2}%
lm\right)  -\left(  \left.  jk\right\vert h_{2}ml\right)  ,
\end{equation}
where
\begin{equation}
\left(  \left.  jk\right\vert h_{2}lm\right)  =\int\psi_{j}^{\ast}\left(
\mathbf{x}_{1}\right)  \psi_{k}^{\ast}\left(  \mathbf{x}_{2}\right)  \frac
{1}{\left\Vert \mathbf{r}_{1}-\mathbf{r}_{2}\right\Vert }\psi_{l}\left(
\mathbf{x}_{1}\right)  \psi_{m}\left(  \mathbf{x}_{2}\right)  d\mathbf{x}%
_{1}d\mathbf{x}_{2}%
\end{equation}
and use the notation
\begin{align}
V_{2jklm}  &  =\left\langle jk||lm\right\rangle \nonumber\\
h_{2jklm}  &  =\left(  \left.  jk\right\vert h_{2}lm\right)  .
\end{align}


\subsection{Energy in terms of NGSO's, Phases and FORDO's\label{EnON_NGSO}}

In terms of $\theta$ dependent NGSO's $\left\{  \left\vert e^{i\alpha_{j}%
}\varphi_{j}\right\rangle ;1\leq j\leq2s\right\}  $ defined in eq. $\left(
\ref{en2}\right)  $ we obtain an energy expression factored through the
equivalence classes $\left[  D^{1}\left(  g^{N}\left(  \mathbf{0}\right)
\right)  \right]  $
\begin{align}
\mathcal{E}\left(  \mathbf{n},\mathbf{\psi},\mathbf{\theta}\right)   &
=\frac{\mathcal{H}\left(  \mathbf{n},\mathbf{\psi},\mathbf{\theta}\right)
}{S_{N}\left(  \mathbf{n}\right)  }=\frac{1}{S_{N}\left(  \mathbf{n}\right)
}\left\{  \sum_{1\leq j\leq s}n_{j}S_{N-1}\left(  \jmath\right)  \left\{
\overset{}{\left\langle \varphi_{j}\left\vert h_{1}\varphi_{j}\right.
\right\rangle +\left\langle \varphi_{j+s}\left\vert h_{1}\varphi_{j+s}\right.
\right\rangle }\right.  \right. \nonumber\\
&  \left.  +\overset{}{\,\left\langle \varphi_{j}\varphi_{j+s}|h_{2}%
\varphi_{j}\varphi_{j+s}\right\rangle }\right\}  +\sum_{1\leq j,k\leq s}%
n_{j}^{\frac{1}{2}}n_{k}^{\frac{1}{2}}e^{i\left(  \theta_{j}-\theta
_{k}\right)  }S_{N-1}\left(  \jmath k\right)  \left\langle \varphi_{j}%
\varphi_{j+s}|h_{2}\varphi_{k}\varphi_{k+s}\right\rangle \nonumber\\
&  \qquad\qquad\qquad\qquad\qquad\left.  +\sum_{1\leq j<k\leq s}n_{i}%
n_{j}S_{N-2}\left(  \jmath k\right)  \left\langle \varphi_{j}\varphi
_{k}|||\varphi_{j}\varphi_{k}\right\rangle \right\}  , \label{thetaen}%
\end{align}
where $\left\vert \mathbf{\varphi}\right\rangle $ is a representative set of
NGSO's for the class $\left[  D^{1}\left(  g^{N}\left(  \mathbf{0}\right)
\right)  \right]  $. As $\left(  \mathbf{n},\mathbf{\psi,\theta}\right)  $
uniquely defines a $2N$-electron AGP-FORDO\ the expression eq. $\left(
\ref{thetaen}\right)  $ can be used to define a density matrix energy
functional
\begin{equation}
\mathcal{F}\left(  D^{1}\right)  =\min_{\mathbf{\theta\in}\Theta}%
\mathcal{E}\left(  \mathbf{n},\mathbf{\varphi},\mathbf{\theta}\right)
\end{equation}
over the manifold of AGP\ states, where
\begin{equation}
\Theta=\left\{  \left.  \theta_{j}\right\vert \theta_{1}=0,\,0<\theta_{j}%
\leq2\pi,2\leq j\leq s\right\}  .
\end{equation}
Noting that each class is determined by $\left(  \mathbf{n},\mathbf{\varphi
}\right)  $ this functional can then be minimized to give
\begin{equation}
\mathcal{F}\left(  D_{\#}^{1}\right)  =\min_{D^{1}\in\mathcal{S}_{1}%
}\mathcal{F}\left(  D^{1}\right)  =\min_{\mathbf{n\in}\mathcal{D}%
_{n},\left\vert \mathbf{\varphi}\right\rangle \in\mathcal{C}}\mathcal{F}%
\left(  D^{1}\left(  \mathbf{n},\mathbf{\varphi}\right)  \right)
=\min_{\left\vert g\right\rangle \in\mathcal{H}^{2}}\mathcal{E}\left(
\left\vert g^{N}\right\rangle \right)  , \label{MinDMF}%
\end{equation}
where the domains $\mathcal{D}_{n}$, $\mathcal{C}$ and $\mathcal{S}_{1}$ are
given as%
\begin{align}
\mathcal{D}_{n}  &  =\left\{  \left.  n_{j}\right\vert ;\sum_{1\leq j\leq
s}n_{j}=1,n_{j}\geq0,\quad1\leq j\leq s\right\}  ,\\
\mathcal{C}  &  =\left\{  \left.  \left\vert \mathbf{\varphi}\right\rangle
\right\vert ;\left\vert \mathbf{\varphi}\right\rangle \neq e^{i\mathbf{\theta
}}\left\vert \mathbf{\varphi}^{\prime}\right\rangle ,\text{ when }\left\vert
\mathbf{\varphi}^{\prime}\right\rangle \in\mathcal{C}\right\}
\end{align}
and
\begin{equation}
\mathcal{S}_{1}=\left\{  \left.  D^{1}\right\vert ;Tr\left\{  D^{1}\right\}
=2N,D^{1}\geq0,D^{1}\text{ at least doubly degenerate}\right\}  .
\end{equation}
The solution to eq. $\left(  \ref{MinDMF}\right)  $ characterizes the
$2N$-electron AGP state that minimizes the total energy.

\subsection{Energy Functional in Terms of Group Parameters\label{Group_Param}}

The manifold of $2N$-electron AGP\ states can be parameterized in terms of a
non degenerate ($c_{j}\neq c_{k}$ for $j\neq k$ ) reference geminal state
$\left\vert g_{0}\right\rangle $ and the subgroups, $U_{orb},$ $U_{phase}$ and
$G_{num}$ of the general linear group $GL\left(  \mathcal{H}^{1}\right)  $ of
one particle Hilbert space $\mathcal{H}^{1}$. The orbital variation special
unitary subgroup $U_{orb}$ is defined as
\begin{equation}
U_{orb}=\left\{  \left.  u\left(  \mathbf{\alpha}\right)  =\exp\left\{
i\sum_{1\leq j\neq k\pm s\leq2s}\alpha_{jk}a_{k}^{\dagger}a_{j}\right\}
\right\vert \alpha_{jk}=\alpha_{kj}^{\ast}\right\}  , \label{orb}%
\end{equation}
where the second quantization operators $\left\{  a_{j}^{\dagger}\right\}  $
are defined by their action on the vacuum state $\left\vert \phi\right\rangle
$ by
\begin{equation}
a_{j}^{\dagger}\left\vert \phi\right\rangle =\left\vert \varphi_{j}%
\right\rangle ;\quad1\leq j\leq2s.
\end{equation}
The phase variation unitary subgroup $U_{phase}$ by%
\begin{equation}
U_{phase}=\left\{  \left.  p\left(  \mathbf{\theta}\right)  =\exp\left\{
\sum_{1\leq j\leq s}\theta_{j}\Omega_{j}\right\}  \right\vert 0\leq\theta
_{j}<2\pi\right\}  , \label{phase}%
\end{equation}
where
\begin{equation}
\Omega_{j}=i\left(  a_{j}^{\dagger}a_{j}+a_{j+s}^{\dagger}a_{j+s}\right)
;\quad1\leq j\leq s.
\end{equation}
The occupation variation non unitary subgroup $G_{num}$
\begin{equation}
G_{num}=\left\{  \left.  k\left(  \mathbf{\eta}\right)  =\exp\left\{
\sum_{1\leq j\leq s}\eta_{j}\Upsilon_{j}\right\}  \right\vert \eta_{j}%
\in\mathbb{R}\right\}  , \label{occup}%
\end{equation}
where
\begin{equation}
\Upsilon_{j}=a_{j}^{\dagger}a_{j}+a_{j+s}^{\dagger}a_{j+s};\quad1\leq j\leq s.
\end{equation}
The manifold of $2N$-electron AGP\ states $\mathcal{M}$ can then be expressed
as
\begin{equation}
\mathcal{M=}\left\{  \left.  k\left(  \mathbf{\eta}\right)  p\left(
\mathbf{\theta}\right)  u\left(  \mathbf{\alpha}\right)  \left\vert g_{0}%
^{N}\right\rangle \right\vert \quad0\leq\theta_{j}<2\pi;\,\eta_{j}%
\in\mathbb{R};\,\alpha_{jk}=\alpha_{kj}^{\ast},1\leq j\neq k\pm s\leq
2s\in\mathbb{C}\right\}
\end{equation}
Explicitly the action on the reference geminal is given by
\begin{equation}
\left\vert g\left(  \mathbf{\eta},\mathbf{\alpha},\mathbf{\theta}\right)
\right\rangle =k\left(  \mathbf{\eta}\right)  p\left(  \mathbf{\theta}\right)
u\left(  \mathbf{\alpha}\right)  \left\vert g_{0}\right\rangle =\sum_{1\leq
l\leq s}e^{\eta_{l}}e^{i\theta_{l}}c_{0l}\left\vert \varphi_{l}\left(
\mathbf{\alpha}\right)  \varphi_{l+s}\left(  \mathbf{\alpha}\right)
\right\rangle ,
\end{equation}
where
\begin{equation}
\left\vert \varphi_{l}\left(  \mathbf{\alpha}\right)  \right\rangle
=\exp\left\{  i\sum_{1\leq j\neq k\pm s\leq2s}\alpha_{jk}a_{j}^{\dagger}%
a_{k}\right\}  \left\vert \varphi_{0l}\right\rangle .
\end{equation}
It should be noted that the special unitary group
\begin{equation}
U_{2s}=\left\{  \left.  \exp i\left(  \sum_{1\leq j\leq s}\kappa_{jj}\left(
a_{j}^{\dagger}a_{j}-a_{j+s}^{\dagger}a_{j+s}\right)  +\kappa_{jj+s}%
a_{j}^{\dagger}a_{j+s}+\kappa_{jj+s}^{\ast}a_{j+s}^{\dagger}a_{j}\right)
\right\vert \kappa_{jj}\in\mathbb{R},\kappa_{jj+s}\in\mathbb{C}\right\}
\label{iso}%
\end{equation}
leaves the geminal $\left\vert g_{0}\right\rangle $ and AGP $\left\vert
g_{0}^{N}\right\rangle $ invariant as
\begin{equation}
\left(  a_{j}^{\dagger}a_{j}-a_{j+s}^{\dagger}a_{j+s}\right)  \left\vert
g_{0}\right\rangle =a_{j}^{\dagger}a_{j+s}\left\vert g_{0}\right\rangle
=a_{j+s}^{\dagger}a_{j}\left\vert g_{0}\right\rangle =0;\,1\leq j\leq s.
\label{iso1}%
\end{equation}


In terms of all these group parameters the energy functional can be expressed
as
\begin{equation}
\mathcal{E}\left(  \mathbf{\eta},\mathbf{\alpha},\mathbf{\theta}\right)
=\frac{\mathcal{H}\left(  \mathbf{\eta},\mathbf{\alpha},\mathbf{\theta
}\right)  }{S_{N}\left(  \mathbf{\eta}\right)  }=\frac{\left\langle \left.
g\left(  \mathbf{\eta},\mathbf{\alpha},\mathbf{\theta}\right)  ^{N}\right\vert
Hg\left(  \mathbf{\eta},\mathbf{\alpha},\mathbf{\theta}\right)  ^{N}%
\right\rangle }{S_{N}\left(  \mathbf{\eta}\right)  }, \label{gp_funct}%
\end{equation}
where
\begin{equation}
\mathcal{H}\left(  \mathbf{\eta},\mathbf{\lambda},\mathbf{\theta}\right)
=\left\langle \left.  k\left(  \mathbf{\eta}\right)  p\left(  \mathbf{\theta
}\right)  u\left(  \mathbf{\alpha}\right)  g_{0}^{N}\right\vert Hk\left(
\mathbf{\eta}\right)  p\left(  \mathbf{\theta}\right)  u\left(  \mathbf{\alpha
}\right)  g_{0}^{N}\right\rangle
\end{equation}
and
\begin{equation}
S_{N}\left(  \mathbf{\eta}\right)  =\left\langle \left.  k\left(
\mathbf{\eta}\right)  g_{0}^{N}\right\vert k\left(  \mathbf{\eta}\right)
g_{0}^{N}\right\rangle =\sum_{1\leq j_{1}<\cdots<j_{N}\leq s}e^{\eta_{j_{1}%
}+\cdots+\eta_{j_{N}}}n_{j_{1}}\cdots n_{j_{N}}.
\end{equation}
The parameters $\left(  \mathbf{\eta,\alpha}\right)  $ uniquely characterize
the FORDO $D^{1}\left(  g^{N}\left(  \mathbf{\eta},\mathbf{\alpha}\right)
\right)  $ i.e.
\begin{equation}
D^{1}\left(  \mathbf{\eta},\mathbf{\alpha}\right)  =D^{1}\left(
\mathbf{n}\left(  \mathbf{\eta}\right)  ,\mathbf{\varphi}\left(
\mathbf{\alpha}\right)  \right)
\end{equation}
so that a DMF can be unambiguously expressed as
\begin{equation}
\mathcal{F}\left(  \left(  \mathbf{\eta},\mathbf{\alpha}\right)  \right)
=\mathcal{F}\left(  D^{1}\left(  \mathbf{\eta},\mathbf{\alpha}\right)
\right)  =\min_{\mathbf{\theta\in}\Theta}\mathcal{E}\left(  \mathbf{\eta
},\mathbf{\alpha},\mathbf{\theta}\right)
\end{equation}
and $\mathcal{F}\left(  \left(  \mathbf{\eta},\mathbf{\alpha}\right)  \right)
$ can be minimized in an unconstrained fashion to give the optimized AGP
energy
\begin{equation}
\mathcal{F}\left(  D_{\min}^{1}\right)  =\mathcal{F}\left(  \left(
\mathbf{\eta}_{\min},\mathbf{\alpha}_{\min}\right)  \right)  =\min
_{\mathbf{\eta},\mathbf{\alpha}}\mathcal{F}\left(  \left(  \mathbf{\eta
},\mathbf{\alpha}\right)  \right)  .
\end{equation}


\section{Necessary Conditions for the Minimization of the Energy
\label{Evariation}}

We first consider the energy functional in terms of $\left(  \mathbf{n,\psi
}\right)  ,$ which produces a constrained optimization problem, then in terms
of the unconstrained variables $\left(  \mathbf{\eta},\mathbf{\alpha
},\mathbf{\theta}\right)  $. The first approach introduces a generalized Fock
operator, while the second approach shows that the Fock operator $F\left(
\mathbf{n,\psi}\right)  $ must be self adjoint when based on an optimized state.

\subsection{Constrained Variational Problem\label{CVP}}

The solution to the constrained optimization problem
\begin{equation}
\min_{\left(  \mathbf{n,\psi}\right)  }\mathcal{E}\left(  \mathbf{n,\psi
}\right)  \label{min}%
\end{equation}
subject to the orthonormality constraints%
\begin{equation}
\left\langle \left.  \psi_{j}\right\vert \psi_{k}\right\rangle =\delta
_{jk};\quad1\leq j,k\leq2s \label{orthonorm}%
\end{equation}
and the constraints of positivity and normalization on the geminal occupation
numbers
\begin{subequations}
\begin{align}
n_{j}  &  \geq0;\quad1\leq j\leq s,\label{pos1}\\
\sum_{1\leq j\leq s}n_{j}  &  =1 \label{gemnorm}%
\end{align}
characterize an optimal $2N$-electron AGP\ state.

The occupation number constraints eqs. $\left(  \ref{pos1}\right)  $, $\left(
\ref{gemnorm}\right)  $ and the orthonormality constraints can be enforced by
the method of Lagrange multipliers and the introduction of a Lagrangian
\end{subequations}
\begin{equation}
\mathcal{L}_{T}\left(  \mathbf{n,\psi,\varepsilon,\xi,}\eta\right)
=\mathcal{E}\left(  \mathbf{n,\psi}\right)  -\sum_{1\leq j,k\leq2s}%
\varepsilon_{jk}\left(  \left\langle \left.  \psi_{j}\right\vert \psi
_{k}\right\rangle -\delta_{jk}\right)  -\eta\left(  \sum_{1\leq j\leq s}%
n_{j}-1\right)  +\sum_{1\leq j\leq s}\xi_{j}n_{j}. \label{Lar}%
\end{equation}
However the occupation number constraints are most effectively invoked by
expressing the geminal occupation numbers as unconstrained functions of group
parameters introduced in section (\ref{Group_Param}). The orbital constraint
part
\begin{equation}
\mathcal{L}\left(  \mathbf{n,\psi}\right)  =\mathcal{E}\left(  \mathbf{n,\psi
}\right)  -\sum_{1\leq j,k\leq2s}\varepsilon_{kj}\left(  \left\langle \left.
\psi_{j}\right\vert \psi_{k}\right\rangle -\delta_{jk}\right)  . \label{larE}%
\end{equation}
of the Lagrangian eq. $\left(  \ref{Lar}\right)  $ reveals some properties of
a generalized Fock operator $F\left(  \mathbf{n,\psi}\right)  $ (see section
\ref{Fock}). The stationarity conditions
\begin{equation}
\frac{\delta\mathcal{L}\left(  \mathbf{n,\psi}\right)  }{\delta\psi_{j}^{\ast
}}=0;\quad1\leq j\leq2s,
\end{equation}
which involve the derivatives
\begin{equation}
\frac{\delta\mathcal{E}\left(  \mathbf{n,\psi}\right)  }{\delta\psi_{j}^{\ast
}}=\frac{1}{S_{N}\left(  \mathbf{n}\right)  }\left\{  \frac{\delta
\mathcal{H}\left(  \mathbf{n,\psi}\right)  }{\delta\psi_{j}^{\ast}%
}-\mathcal{E}\left(  \mathbf{n,\psi}\right)  \frac{\delta S_{N}\left(
\mathbf{n}\right)  }{\delta\psi_{j}^{\ast}}\right\}  =\frac{1}{S_{N}\left(
\mathbf{n}\right)  }\frac{\delta\mathcal{H}\left(  \mathbf{n,\psi}\right)
}{\delta\psi_{j}^{\ast}}%
\end{equation}
produce
\begin{equation}
\frac{\delta\mathcal{E}\left(  \mathbf{n,\psi}\right)  }{\delta\psi_{j}^{\ast
}}=\sum_{1\leq j,k\leq2s}\left\vert \psi_{k}\right\rangle \varepsilon_{jk}.
\label{statcon}%
\end{equation}
We show in section \ref{Fock} that
\begin{equation}
\frac{\delta\mathcal{E}\left(  \mathbf{n,\psi}\right)  }{\delta\psi_{j}^{\ast
}}=F\left(  \mathbf{n,\psi}\right)  \left\vert \psi_{j}\right\rangle
\end{equation}
so eq. $\left(  \ref{statcon}\right)  $ becomes
\begin{equation}
F\left(  \mathbf{n,\psi}\right)  \left\vert \psi_{j}\right\rangle =\sum_{1\leq
j,k\leq2s}\left\vert \psi_{k}\right\rangle \varepsilon_{jk}. \label{Fockeq}%
\end{equation}
As (see section \ref{Fock})%
\begin{equation}
\frac{\delta\mathcal{E}\left(  \mathbf{n,\psi}\right)  }{\delta\psi_{j}%
}=\left(  \frac{\delta\mathcal{E}\left(  \mathbf{n,\psi}\right)  }{\delta
\psi_{j}^{\ast}}\right)  ^{\ast}%
\end{equation}
it is sufficient to consider derivatives wrt $\left\{  \psi_{j}^{\ast
}\right\}  $.

In \emph{contrast} to the IPS\ case, but in \emph{common} with Multi
Configurational Self Consistent Field (MCSCF)\ procedures, eq. $\left(
\ref{Fockeq}\right)  $ is not an eigenvalue equation. However by considering
the stationarity conditions from the unitary group parameterization one can
observe that a stationary point of the orbital variation $F\left(
\mathbf{n,\psi}\right)  $ becomes self adjoint and hence
\begin{equation}
\varepsilon_{jk}=\varepsilon_{kj}^{\ast}%
\end{equation}
and the matrix $\mathbf{\varepsilon}$ of Lagrange parameters is hermitian.

\subsection{Unconstrained Variational Problem\label{UCVP}}

In this case we consider the problem
\begin{equation}
\min_{\left(  \mathbf{\eta,\alpha,\theta}\right)  }\mathcal{E}\left(
\mathbf{\eta,\alpha,\theta}\right)  ,
\end{equation}
where the energy functional is given in eq. $\left(  \ref{gp_funct}\right)  $
in terms of the groups eqs. $\left(  \ref{orb}\right)  ,\left(  \ref{phase}%
\right)  $ and $\left(  \ref{occup}\right)  $ parametrized by $\left(
\mathbf{\eta,\alpha,\theta}\right)  $. Setting
\begin{equation}
g_{0}\equiv g\left(  \mathbf{0,0,0}\right)
\end{equation}
this leads to (see Appendix \ref{Group_Param}) the orbital variation
derivatives
\begin{equation}
\partial_{\mathbf{\lambda}_{jk}}^{1}\mathcal{E}\left(  g_{0}^{N}\right)
=\frac{i}{S_{N}\left(  \mathbf{0}\right)  }\left\langle \left.  g_{0}%
^{N}\right\vert \left[  H,a_{j}^{\dagger}a_{k}\right]  g_{0}^{N}\right\rangle
;\quad1\leq j\neq k\pm s\leq2s,
\end{equation}
the phase variation derivatives
\begin{equation}
\partial_{\mathbf{\theta}_{j}}^{1}\mathcal{E}\left(  g_{0}^{N}\right)
=\frac{i}{S_{N}\left(  \mathbf{0}\right)  }\left\langle \left.  g_{0}%
^{N}\right\vert \left[  H,\left(  a_{j}^{\dagger}a_{j}+a_{j+s}^{\dagger
}a_{j+s}\right)  \right]  g_{0}^{N}\right\rangle ;\quad1\leq j\leq s,
\end{equation}
and the occupation number variations derivatives
\begin{align}
\partial_{\mathbf{\eta}_{j}}^{1}\mathcal{E}\left(  g_{0}^{N}\right)   &
=\frac{1}{S_{N}\left(  \mathbf{0}\right)  }\left\{  \left\langle \left.
g_{0}^{N}\right\vert \left[  H,\left(  a_{j}^{\dagger}a_{j}+a_{j+s}^{\dagger
}a_{j+s}\right)  \right]  _{+}^{N}g_{0}\right\rangle \right. \nonumber\\
&  -\left.  2\mathcal{E}\left(  g_{0}^{N}\right)  \left\langle \left.
g_{0}^{N}\right\vert \left(  a_{j}^{\dagger}a_{j}+a_{j+s}^{\dagger}%
a_{j+s}\right)  g_{0}^{N}\right\rangle \right\}  ;\quad1\leq j\leq s.
\end{align}
All the derivatives can be evaluated at $\left(  \mathbf{0},\mathbf{0}%
,\mathbf{0}\right)  $ by updating the reference geminal $\left\vert
g_{0}\right\rangle \equiv\left\vert g_{old}\right\rangle $ by a
transformation
\begin{equation}
\left\vert g_{new}\right\rangle =k\left(  \mathbf{\eta}\right)  p\left(
\mathbf{\theta}\right)  u\left(  \mathbf{\alpha}\right)  \left\vert
g_{old}\right\rangle
\end{equation}
to produce a new reference $\left\vert g_{0}\right\rangle \equiv\left\vert
g_{new}\right\rangle .$ The stationarity conditions for orbital variation
gives
\begin{equation}
\left\langle \left.  g_{0}^{N}\right\vert \left[  H,a_{j}^{\dagger}%
a_{k}\right]  g_{0}^{N}\right\rangle =0;\quad1\leq j\neq k\pm s\leq2s,
\label{ostat}%
\end{equation}
while phase stationarity gives%
\begin{equation}
\left\langle \left.  g_{0}^{N}\right\vert \left[  H,\left(  a_{j}^{\dagger
}a_{j}+a_{j+s}^{\dagger}a_{j+s}\right)  \right]  g_{0}^{N}\right\rangle
=0;\quad1\leq j\leq s,
\end{equation}
which is equivalent to
\begin{equation}
\left\langle \left.  g_{0}^{N}\right\vert \left[  H,a_{j}^{\dagger}%
a_{j}\right]  g_{0}^{N}\right\rangle =0;\quad1\leq j\leq2s. \label{phasestat}%
\end{equation}
Eqs. $\left(  \ref{ostat}\right)  $ and $\left(  \ref{phasestat}\right)  $
together with the invariance group relationships eq. $\left(  \ref{iso1}%
\right)  $ give
\begin{equation}
\left\langle \left.  g_{0}^{N}\right\vert \left[  H,a_{j}^{\dagger}%
a_{k}\right]  g_{0}^{N}\right\rangle =0;\quad1\leq j,k\leq2s. \label{orbstat}%
\end{equation}


The preceding commutator relationships are the same as are produced by
stationary states, $\left\vert \Psi\right\rangle ,$ in general MCSCF
procedures, i.e.
\begin{equation}
\left\langle \left.  \Psi\right\vert \left[  H,a_{j}^{\dagger}a_{k}\right]
\Psi\right\rangle =0;\quad1\leq j,k\leq2s.
\end{equation}
Optimized AGP states are also stationary to occupation number variations
giving
\begin{equation}
\left\langle \left.  g_{0}^{N}\right\vert \left[  H,\left(  a_{j}^{\dagger
}a_{j}+a_{j+s}^{\dagger}a_{j+s}\right)  \right]  _{+}g_{0}^{N}\right\rangle
-2\mathcal{E}\left(  g_{0}^{N}\right)  \left\langle \left.  g_{0}%
^{N}\right\vert \left(  a_{j}^{\dagger}a_{j}+a_{j+s}^{\dagger}a_{j+s}\right)
g_{0}^{N}\right\rangle =0;\quad1\leq j\leq s,
\end{equation}
which leads to the new relationship
\begin{equation}
\left\langle \left.  g_{0}^{N}\right\vert Ha_{j}^{\dagger}a_{j}g_{0}%
^{N}\right\rangle =\mathcal{E}\left(  g_{0}^{N}\right)  \left\langle \left.
g_{0}^{N}\right\vert a_{j}^{\dagger}a_{j}g_{0}^{N}\right\rangle ;\quad1\leq
j\leq2s,
\end{equation}
which can be recast as
\begin{equation}
\frac{1}{S_{N}\left(  \mathbf{0}\right)  }\left\langle \left.  g_{0}%
^{N}\right\vert Ha_{j}^{\dagger}a_{j}g_{0}^{N}\right\rangle =\mathcal{E}%
\left(  g_{0}^{N}\right)  N_{j};\quad1\leq j\leq2s. \label{occnumstat}%
\end{equation}
In the context of general multi configurational states the generators
$\left\{  a_{j}^{\dagger}a_{j}\right\}  $ produce CI\ coefficients variations,
hence even those cases
\begin{equation}
\frac{1}{\left\langle \left.  \Psi\right\vert \Psi\right\rangle }\left\langle
\left.  \Psi\right\vert Ha_{j}^{\dagger}a_{j}\Psi\right\rangle =\mathcal{E}%
\left(  \Psi\right)  N_{j};\quad1\leq j\leq2s.
\end{equation}


\subsection{Generalized Brillouin Condition for Optimized AGP
States.\label{Brillouin}}

If an AGP\ has non degenerate canonical coefficients i.e. $c_{j}\neq c_{k}$
for $1\leq j\neq k\leq s$ one can \cite{weiner2004}\cite{Goscinski80}%
\cite{Kurtz82}\cite{Ohrn82} always construct operators (see Appendix
\ref{q_op}) $\left\{  \left.  q_{jk}^{\dagger}\right\vert \quad1\leq j\neq
k\pm s\leq2s\right\}  $ that have the property
\begin{equation}
q_{jk}^{\dagger}\left\vert g_{0}^{N}\right\rangle \neq0,\quad q_{jk}\left\vert
g_{0}^{N}\right\rangle =0;\quad1\leq j\neq k\pm s\leq2s,
\end{equation}
such that the orbital transformation group $U_{orb}$ can also be expressed as
\begin{equation}
U_{orb}=\left\{  u\left(  \mathbf{\beta}\right)  =\exp\left\{  i\sum
_{1\leq\left\vert j\right\vert <\left\vert k\right\vert \leq s}\left\{
\beta_{jk}q_{jk}^{\dagger}+\beta_{jk}^{\ast}q_{jk}\right\}  \right\}
\right\}  ,
\end{equation}
where
\begin{equation}
\left\vert j\right\vert =\left\{
\begin{array}
[c]{c}%
j;\qquad1\leq j\leq s\\
j-s;\qquad s+1\leq j\leq2s
\end{array}
\right.  .
\end{equation}
This parameterization of $U_{orb}$ leads to the stationarity conditions
\begin{equation}
\left\langle \left.  g_{0}^{N}\right\vert \left[  H,q_{jk}^{\dagger}\right]
g_{0}^{N}\right\rangle =\left\langle \left.  g_{0}^{N}\right\vert \left[
H,q_{jk}\right]  g_{0}^{N}\right\rangle =0;\quad1\leq\left\vert j\right\vert
<\left\vert k\right\vert \leq s,
\end{equation}
that give the "Brillouin Conditions"
\begin{equation}
\left\langle \left.  g_{0}^{N}\right\vert Hq_{jk}^{\dagger}g_{0}%
^{N}\right\rangle =\left\langle \left.  g_{0}^{N}\right\vert q_{jk}Hg_{0}%
^{N}\right\rangle =0;\quad1\leq\left\vert j\right\vert <\left\vert
k\right\vert \leq s. \label{BC}%
\end{equation}
The $\left\{  q_{jk}^{\dagger},q_{jk}\right\}  $ basis can be transformed into
the $\left\{  a_{j}^{\dagger}a_{k}\right\}  $ basis of the Lie algebra of
$U_{orb}$ by a block diagonal non singular transformation, $\mathbb{T},$ such
that the blocks $\mathbb{T}\left(  jk\right)  $ have the effect
\begin{align}
&  \left(  q_{jk}^{\dagger},q_{j+sk}^{\dagger},q_{jk+s}^{\dagger}%
,q_{j+sk+s}^{\dagger},q_{jk},q_{j+sk},q_{jk+s},q_{j+sk+s}\right)
\mathbb{T}\left(  jk\right) \nonumber\\
&  \qquad\qquad\qquad\qquad\qquad\left.  =\right.  \left(  a_{j}^{\dagger
}a_{k},a_{j}^{\dagger}a_{k+s},a_{j+s}^{\dagger}a_{k},a_{j+s}^{\dagger}%
a_{k+s},a_{k}^{\dagger}a_{j},a_{k}^{\dagger}a_{j+s},a_{k+s}^{\dagger}%
a_{j},a_{k+s}^{\dagger}a_{j+s}\right)
\end{align}
giving
\begin{equation}
a_{l}^{\dagger}a_{m}=\sum_{p,q=0,s}\left\{  q_{j+pk+q}^{\dagger}\sigma\left(
jk\right)  _{pqlm}+q_{j+pk+q}\tau\left(  jk\right)  _{pqlm}\right\}  .
\label{Trans}%
\end{equation}
The Brillouin Conditions eq. $\left(  \ref{BC}\right)  $ combined with eq.
$\left(  \ref{Trans}\right)  $ results in the generalized Brillouin
Conditions
\begin{equation}
\left\langle \left.  g_{0}^{N}\right\vert Ha_{l}^{\dagger}a_{m}g_{0}%
^{N}\right\rangle =0;\quad1\leq l\neq m\pm s\leq2s, \label{GBC}%
\end{equation}
as%
\begin{equation}
\left\langle \left.  g_{0}^{N}\right\vert H\sum_{p,q=0,s}q_{j+pk+q}^{\dagger
}\sigma\left(  jk\right)  _{pqlm}-\sum_{p,q=0,s}q_{j+pk+q}\tau\left(
jk\right)  _{pqlm}Hg_{0}^{N}\right\rangle =0;\quad1\leq l\neq m\pm s\leq2s.
\end{equation}
The remaining off diagonal particle-hole operators $\left\{  \left.
a_{l}^{\dagger}a_{m}\right\vert 1\leq l=m\pm s\leq2s\right\}  $ automatically
satisfy a generalized Brillouin condition
\begin{equation}
\left\langle \left.  g_{0}^{N}\right\vert Ha_{l}^{\dagger}a_{m}g_{0}%
^{N}\right\rangle =0;\quad1\leq l=m\pm s\leq2s
\end{equation}
by the invariance group relationships eq. $\left(  \ref{iso1}\right)  .$

If there are degeneracies of the canonical coefficients i.e. $c_{j}=c_{k}$ for
some $1\leq j\neq k\leq s,$ then the $\left\{  q_{jk}^{\dagger}\right\}  $ and
$\left\{  q_{jk}\right\}  $ associated with those coefficients can be replaced
\cite{weiner2004}\cite{Goscinski82} with a sets of self adjoint operators
$\left\{  u_{jk}\right\}  $ and $\left\{  v_{ij}\right\}  $ respectively,
where
\begin{equation}
u_{jk}\left\vert g_{0}^{N}\right\rangle \neq0,\quad v_{jk}\left\vert g_{0}%
^{N}\right\rangle =0.
\end{equation}
and one can deduce that a generalized Brillouin condition of the form eq.
$\left(  \ref{GBC}\right)  $ \emph{does not} hold for $\left\{  a_{j}%
^{\dagger}a_{k}\right\}  $ associated with those indices.

\section{Generalized Fock Operator \label{Fock} for the Variation wrt
$\mathbf{\psi}$}

\subsection{Functional Derivatives\label{FD}}

The energy expression eq. $\left(  \ref{energy}\right)  $ can be
differentiated wrt $\mathbf{\psi}^{\ast},$ $\mathbf{\psi}$ to produce the
functional derivatives (Appendix \ref{variational})
\begin{equation}
\left\{  \frac{\delta\mathcal{E}\left(  \mathbf{n,\psi}\right)  }{\delta
\psi_{j}^{\ast}},\frac{\delta\mathcal{E}\left(  \mathbf{n,\psi}\right)
}{\delta\psi_{j}};\quad1\leq j\leq s\right\}  ,
\end{equation}
which can be expressed as%
\begin{align}
\frac{\delta\mathcal{E}\left(  \mathbf{n,\psi}\right)  }{\delta\psi_{j}^{\ast
}}  &  =\frac{1}{S_{N}\left(  \mathbf{n}\right)  }\frac{\delta\mathcal{H}%
\left(  \mathbf{n,\psi}\right)  }{\delta\psi_{j}^{\ast}}\\
\frac{\delta\mathcal{E}\left(  \mathbf{n,\psi}\right)  }{\delta\psi_{j}}  &
=\frac{1}{S_{N}\left(  \mathbf{n}\right)  }\frac{\delta\mathcal{H}\left(
\mathbf{n,\psi}\right)  }{\delta\psi_{j}}%
\end{align}
in terms of the derivatives of the unnormalized energy $\mathcal{H}\left(
\mathbf{n,\psi}\right)  ,$ as $S_{N}\left(  \mathbf{n}\right)  $ does not
depend on $\mathbf{\psi}$ when $\left\{  \left.  \left\vert \psi
_{j}\right\rangle \right\vert \quad1\leq j\leq2s\right\}  $ are orthonormal.
It is easy to see that
\begin{equation}
\frac{\delta\mathcal{E}\left(  \mathbf{n,\psi}\right)  }{\delta\psi_{j}^{\ast
}}=\left(  \frac{\delta\mathcal{E}\left(  \mathbf{n,\psi}\right)  }{\delta
\psi_{j}}\right)  ^{\ast}%
\end{equation}
so we only need consider $\left\{  \frac{\delta\mathcal{H}\left(
\mathbf{n,\psi}\right)  }{\delta\psi_{j}^{\ast}}\right\}  .$ These derivatives
can be expressed $\left(  \text{see Appendix \ref{variational}}\right)  $ in
terms of a generalized Fock operator $F\left(  \mathbf{n,\psi}\right)  $ as
\begin{equation}
\frac{\delta\mathcal{E}\left(  \mathbf{n,\psi}\right)  }{\delta\psi_{j}^{\ast
}}=F\left(  \mathbf{n,\psi}\right)  \left\vert \psi_{j}\right\rangle
,\quad1\leq j\leq2s, \label{H_deriv}%
\end{equation}
that depends on the occupation numbers $\mathbf{n,}$ and the GSO's
$\mathbf{\psi}$ and is given by
\begin{equation}
F\left(  \mathbf{n,\psi}\right)  =\sum_{1\leq m,n\leq2s}\left\{  \left(
\mathbb{H}_{1}\left(  \mathbf{\psi}\right)  \mathbb{D}^{1}\left(
\mathbf{n}\right)  \right)  _{mn}+\frac{1}{2}\left(  L_{2}^{1}\left(
\mathbb{V}_{2}\left(  \mathbf{\psi}\right)  \mathbb{D}^{2}\left(
\mathbf{n}\right)  \right)  \right)  _{mn}\right\}  \left\vert \psi
_{m}\right\rangle \left\langle \psi_{n}\right\vert . \label{genfock}%
\end{equation}
The contraction operator $L_{2}^{1}$ is defined as
\begin{equation}
L_{2}^{1}\left(  \mathbb{A}\right)  _{jk}=\sum_{1\leq l\leq2s}A_{jlkl},
\end{equation}
$\mathbb{H}_{1}\left(  \mathbf{\psi}\right)  $ is the matrix of the one
particle Hamiltonian $h_{1}$, $\mathbb{V}_{2}\left(  \mathbf{\psi}\right)  $
is the matrix of the two particle Hamiltonian $h_{2}$, $\mathbb{D}^{1}\left(
\mathbf{n}\right)  $ the FORDM and $\mathbb{D}^{2}\left(  \mathbf{n}\right)  $
the SORDM of the $2N$-elecron AGP\ state (see in eq. $\left(  \ref{D2C}%
\right)  ),$ wrt the basis $\left\{  \left\vert \psi_{j}\right\rangle ;1\leq
j\leq2s\right\}  .$ One should recall that in L\"{o}wdin normalization and
using an unordered basis of two electron space
\begin{equation}
\mathbb{D}^{1}\left(  \mathbf{n}\right)  =\left(  2N-1\right)  ^{-1}L_{2}%
^{1}\left(  \mathbb{D}^{2}\left(  \mathbf{n}\right)  \right)  .
\end{equation}
The expression eq. $\left(  \ref{genfock}\right)  $ is in fact valid for all
$N$-representable. SORDM $\mathbb{D}^{2}$ and its associated FORDM
$\mathbb{D}^{1}$ and defines a generalized Fock operator for \emph{any }MCSCF
state and always produces the orbital functional derivatives eq. $\left(
\ref{H_deriv}\right)  .$ In general this generalized Fock operator is
\emph{not }self adjoint\emph{ }unless the state that it is based on fully
optimizes the energy functional.

\subsection{Properties of a Stationary Point Fock Operator\label{Stat_Fock}}

We show in Appendix \ref{VarGroupPar} that the derivatives wrt the group
parameters controlling the GSO variations are
\begin{equation}
\frac{\partial\mathcal{E}\left(  \mathbb{D}^{2},\mathbf{0}\right)  }%
{\partial\alpha_{pq}}=i\left(  \left[  \mathbb{D}^{1},\mathbb{H}_{1}\left(
\mathbf{\varphi}\left(  \mathbf{0}\right)  \right)  \right]  +\frac{1}{2}%
L_{2}^{1}\left(  \left[  \mathbb{D}^{2},\mathbb{V}_{2}\left(  \mathbf{\varphi
}\left(  \mathbf{0}\right)  \right)  \right]  \right)  \right)  _{qp}%
\end{equation}
and in Appendix \ref{Deriv_Fock} that the generalized Fock operator matrix
elements are
\begin{equation}
F\left(  \mathbb{D}^{2},\mathbf{\varphi}\right)  _{qp}=\operatorname{Tr}%
\left\{  \left(  \mathbb{H}_{1}\left(  \mathbf{\varphi}\right)  \mathbb{D}%
^{1}\right)  _{qp}+\frac{1}{2}\left(  L_{2}^{1}\left(  \mathbb{V}_{2}\left(
\mathbf{\varphi}\right)  \mathbb{D}^{2}\right)  \right)  _{qp}\right\}
\end{equation}
for \emph{any} state described by a SORDM $\mathbb{D}^{2}$. These two
equations give that
\begin{equation}
\frac{\partial\mathcal{E}\left(  \mathbb{D}^{2},\mathbf{0}\right)  }%
{\partial\alpha_{pq}}=i\left(  F\left(  \mathbb{D}^{2},\mathbf{\varphi}\left(
\mathbf{0}\right)  \right)  ^{\dagger}-F\left(  \mathbb{D}^{2},\mathbf{\varphi
}\left(  \mathbf{0}\right)  \right)  \right)  _{qp},
\end{equation}
which show if the state is stationary wrt orbital transformations then the
generalized Fock operator is self adjoint as
\begin{equation}
\frac{\partial\mathcal{E}\left(  \mathbb{D}^{2},\mathbf{0}\right)  }%
{\partial\alpha_{pq}}=0=i\left(  F\left(  \mathbb{D}^{2},\mathbf{\varphi
}\left(  \mathbf{0}\right)  \right)  ^{\dagger}-F\left(  \mathbb{D}%
^{2},\mathbf{\varphi}\left(  \mathbf{0}\right)  \right)  \right)  _{qp}%
;\quad1\leq p,q\leq2s.
\end{equation}


\subsection{Properties of an Optimized AGP\ Fock Operator\label{AGP_Fock}}

The generalized Fock operator associated with an optimized AGP\ state has
further significant properties, in addition to being self adjoint, originating from

\begin{enumerate}
\item the stationarity wrt occupation number variations shown in eq. $\left(
\ref{occnumstat}\right)  $ as
\begin{equation}
\frac{1}{S_{N}\left(  g^{N}\right)  }\left\langle \left.  g^{N}\right\vert
Ha_{p}^{\dagger}a_{p}g^{N}\right\rangle =\mathcal{E}\left(  g^{N}\right)
N_{p};\quad1\leq p\leq2s,
\end{equation}
and

\item the expression for the generalized Fock operator matrix elements
\begin{equation}
F\left(  \mathbb{D}^{2},\mathbf{\varphi}\right)  _{qp}=\frac{1}%
{\operatorname{Tr}\left\{  D^{2N}\right\}  }\operatorname{Tr}\left\{
a_{p}^{\dagger}\left[  a_{q},H\right]  D^{2N}\right\}
\end{equation}
obtained in Appendix \ref{SQ_Fock}, which are valid for any state described by
the positive operator $D^{2N}$. For AGP\ states this becomes
\begin{equation}
F\left(  \mathbf{n},\mathbf{\psi}\right)  _{qp}=\frac{1}{S_{N}\left(
g^{N}\right)  }\left\langle \left.  g^{N}\right\vert a_{p}^{\dagger}%
a_{q}H-a_{p}^{\dagger}Ha_{q}g^{N}\right\rangle .
\end{equation}

\end{enumerate}

Combining items $\left(  1\right)  $ and $\left(  2\right)  $ together shows
that the diagonal elements of the optimized AGP\ Fock operator are given by
\begin{equation}
F\left(  \mathbf{n},\mathbf{\psi}\right)  _{pp}=\frac{1}{S_{N}\left(
g^{N}\right)  }\left\langle \left.  g^{N}\right\vert a_{p}^{\dagger}%
a_{p}H-a_{p}^{\dagger}Ha_{p}g^{N}\right\rangle =\mathcal{E}\left(
g^{N}\right)  N_{p}-\frac{1}{S_{N}\left(  g^{N}\right)  }\left\langle \left.
a_{p}g^{N}\right\vert Ha_{p}g^{N}\right\rangle ,
\end{equation}
which can be developed to give
\begin{equation}
F\left(  \mathbf{n},\mathbf{\psi}\right)  _{pp}=N_{p}\left\{  \mathcal{E}%
\left(  g^{N}\right)  -\frac{\left\langle \left.  a_{p}g^{N}\right\vert
Ha_{p}g^{N}\right\rangle }{\left\langle \left.  a_{p}g^{N}\right\vert
a_{p}g^{N}\right\rangle }\right\}  =N_{p}\left\{  \mathcal{E}\left(
g^{N}\right)  -\mathcal{E}_{2N-1}\left(  \psi_{p}\right)  \right\}  ,
\end{equation}
where $\mathcal{E}_{2N-1}\left(  \psi_{p}\right)  $ is the energy of the
normalized $2N-1$ electron state produced from the state $\frac{1}%
{S_{N}\left(  g^{N}\right)  }\left\vert g^{N}\right\rangle $ by removing an
electron described by the CGSO\ $\left\vert \psi_{p}\right\rangle $. One can
observe from the form of $F\left(  \mathbf{n},\mathbf{\psi}\right)  _{pp}$
that
\begin{equation}
F\left(  \mathbf{n},\mathbf{\psi}\right)  _{pp}=F\left(  \mathbf{n}%
,\mathbf{\psi}\right)  _{p+sp+s}\Rightarrow I\left(  \psi_{p}\right)
=I\left(  \psi_{p+s}\right)  ;\quad1\leq p\leq s.
\end{equation}
In Appendix \ref{Fock&Energy} we show that the total energy can be expressed
in terms of the Fock operator $F\left(  \mathbf{n},\mathbf{\psi}\right)
\ $as
\begin{equation}
\mathcal{E}\left(  \mathbf{n},\mathbf{\psi}\right)  =\frac{1}{2}\left\{
\overset{}{\operatorname{Tr}\left\{  F\left(  \mathbf{n},\mathbf{\psi}\right)
\right\}  +\mathcal{H}_{1}\left(  \mathbf{n},\mathbf{\psi}\right)  }\right\}
,
\end{equation}
where the expectation value of the one electron part of the Hamiltonian
$\mathcal{H}_{1}\left(  \mathbf{n},\mathbf{\psi}\right)  $ is given by
\begin{equation}
\mathcal{H}_{1}\left(  \mathbf{n},\mathbf{\psi}\right)  =\sum_{1\leq p\leq
2s}N_{p}\left(  \left.  \psi_{p}\right\vert h_{1}\psi_{p}\right)  ,
\end{equation}
as the FORDM $\mathbb{D}^{1}\left(  \mathbf{n},\mathbf{\psi}\right)  $ is
diagonal in the CGSO's $\left\{  \left\vert \psi_{p}\right\rangle \right\}  .$
The energy $\mathcal{E}\left(  \mathbf{n},\mathbf{\psi}\right)  $ can be
expressed as
\begin{align}
\mathcal{E}\left(  \mathbf{n},\mathbf{\psi}\right)   &  =\sum_{1\leq p\leq
2s}\left\{  N_{p}I\left(  \psi_{p}\right)  +N_{p}\left(  \left.  \psi
_{p}\right\vert h_{1}\psi_{p}\right)  \right\}  \\
&  =\sum_{1\leq p\leq2s}\left\{  N_{p}I\left(  \psi_{p}\right)  +N_{p}\left(
\left.  \psi_{p}\right\vert h_{1}\psi_{p}\right)  \right\}  \equiv\frac{1}{2}%
%TCIMACRO{\tsum \limits_{1\leq l\leq2s}}%
%BeginExpansion
{\textstyle\sum\limits_{1\leq l\leq2s}}
%EndExpansion
N_{l}\varsigma\left(  \psi_{l}\right)  ,
\end{align}
where we have introduced the new quantity,
\begin{equation}
\varsigma\left(  \psi_{l}\right)  =I\left(  \psi_{l}\right)  +\left(  \left.
\psi_{l}\right\vert h_{1}\psi_{l}\right)  ,
\end{equation}
which is the sum of the ionization potential, the kinetic energy, the electron
- nuclear potential energy of attraction and potential energy of interaction
due to any other external forces of an electron described by a
GSO\ $\left\vert \psi_{l}\right\rangle $. 

Although the AGP\ Fock operator is self adjoint its matrix, which is
hermitian, in the CGSO basis is in general non diagonal. If the canonical
coefficients of the optimized AGP\ are non degenerate then the generalized
Brillouin conditions eq. $\left(  \ref{GBC}\right)  $ hold and
\begin{equation}
F\left(  \mathbf{n},\mathbf{\psi}\right)  _{qp}=\frac{1}{S_{N}\left(
g^{N}\right)  }\left\{  \left\langle \left.  g^{N}\right\vert a_{p}^{\dagger
}a_{p}Hg^{N}\right\rangle \delta_{qp}-\left\langle \left.  a_{p}%
g^{N}\right\vert Ha_{q}g^{N}\right\rangle \right\}  .
\end{equation}
The matrix $\left\{  F\left(  \mathbf{n},\mathbf{\psi}\right)  _{qp}\right\}
$ can be diagonalized by a unitary transformation $u$ of the CGSO basis to
produce
\begin{equation}
\tilde{F}\left(  \mathbf{n},\mathbf{\psi}\right)  _{qp}=\frac{1}{S_{N}\left(
g^{N}\right)  }\left\{  \left\langle \left.  g^{N}\right\vert b_{p}^{\dagger
}b_{p}Hg^{N}\right\rangle -\left\langle \left.  b_{p}g^{N}\right\vert
Hb_{p}g^{N}\right\rangle \right\}  \delta_{qp},
\end{equation}
where
\begin{equation}
b_{p}^{\dagger}=ua_{p}^{\dagger}u^{\dagger}=\sum_{1\leq j\leq2s}a_{j}%
^{\dagger}u_{jp}%
\end{equation}
and
\begin{equation}
\left\vert \varsigma_{p}\right\rangle =\sum_{1\leq j\leq2s}\left\vert \psi
_{j}\right\rangle u_{jp}.
\end{equation}
The eigenvalues
\begin{align}
\varepsilon_{p}  &  =\tilde{F}\left(  \mathbf{n},\mathbf{\psi}\right)
_{pp}=\frac{1}{S_{N}\left(  g^{N}\right)  }\left\{  \left\langle \left.
g^{N}\right\vert b_{p}^{\dagger}b_{p}Hg^{N}\right\rangle -\left\langle \left.
b_{p}g^{N}\right\vert Hb_{p}g^{N}\right\rangle \right\} \nonumber\\
&  =\frac{1}{S_{N}\left(  g^{N}\right)  }\mathcal{N}_{p}\left\{
\mathcal{E}\left(  g^{N}\right)  -\frac{\left\langle \left.  b_{p}%
g^{N}\right\vert Hb_{p}g^{N}\right\rangle }{\left\langle \left.  b_{p}%
g^{N}\right\vert b_{p}g^{N}\right\rangle }\right\}  ,
\end{align}
where
\begin{equation}
\mathcal{N}_{p}=\left\langle \left.  \varsigma_{p}\right\vert D^{1}\left(
\mathbf{n},\mathbf{\psi}\right)  \varsigma_{p}\right\rangle .
\end{equation}
Hence
\begin{equation}
\varepsilon_{p}=\frac{\mathcal{N}_{p}I\left(  \varsigma_{p}\right)  }%
{S_{N}\left(  g^{N}\right)  },
\end{equation}
where $\left\{  I\left(  \varsigma_{p}\right)  \right\}  $ are the ionization
potentials associated with the GSO's $\left\{  \left\vert \varsigma
_{p}\right\rangle \right\}  .$ The optimized geminal\ \emph{is not} in general
in canonical form wrt the basis $\left\{  \left\vert \varsigma_{p}%
\right\rangle \right\}  $ , nor is $D^{1}\left(  \mathbf{n},\mathbf{\psi
}\right)  $ diagonal so $\mathcal{N}_{p}$ are \emph{not} occupation numbers.

\section{Reduction to the IPS \label{IPS}}

The standard basic model in chemistry is the IPS description, as we have noted
in the preceding, this is a special case of the AGP state model which relaxes
to the IPS\ for integer occupation numbers i.e.
\begin{align}
N_{j}  &  =N_{j+s}=1,\quad1\leq j\leq N\nonumber\\
N_{j}  &  =N_{j+s}=0,\quad N+1\leq j\leq s
\end{align}
and%
\begin{align}
n_{j}  &  =n_{j+s}=1,\quad1\leq j\leq N\nonumber\\
n_{j}  &  =n_{j+s}=0,\quad N+1\leq j\leq s,
\end{align}
where the FORDM\ is given by
\begin{equation}
D_{jk}^{1}=N_{j}\delta_{jk};\quad1\leq j,k\leq2s
\end{equation}
The SORDM for an IPS is given by
\begin{equation}
\mathbb{D}^{2}=\mathbb{D}^{1}\wedge\mathbb{D}^{1},
\end{equation}
where
\begin{equation}
\left(  \mathbb{D}^{1}\wedge\mathbb{D}^{1}\right)  _{ijkl}=D_{ik}^{1}%
D_{jl}^{1}-D_{il}^{1}D_{jk}^{1}.
\end{equation}
The generalized Fock operator that we have introduced in this paper
\begin{equation}
F\left(  \mathbb{D}^{2},\mathbf{\varphi}\right)  =\sum_{1\leq m,n\leq
2s}\left\{  \left(  \mathbb{H}_{1}\left(  \mathbf{\varphi}\right)
\mathbb{D}^{1}\right)  _{mn}+\left(  L_{2}^{1}\left(  \mathbb{H}_{2}\left(
\mathbf{\varphi}\right)  \mathbb{D}^{2}\right)  \right)  _{mn}\right\}
\left\vert \varphi_{m}\right\rangle \left\langle \varphi_{n}\right\vert ,
\label{F1}%
\end{equation}
can then be simplified by noting
\begin{align}
&  \left(  L_{2}^{1}\left(  \mathbb{H}_{2}\left(  \mathbf{\varphi}\right)
\mathbb{D}^{2}\right)  \right)  _{mn}=\sum_{1\leq k\leq2s}\left(
\mathbb{H}_{2}\left(  \mathbf{\varphi}\right)  \mathbb{D}^{2}\right)
_{mknk}=\sum_{1\leq p,q,k\leq2s}h_{2mkpq}\left(  \mathbf{\varphi}\right)
D_{pqnk}^{2}\nonumber\\
&  =\sum_{1\leq p,q,k\leq2s}h_{2mkpq}\left(  \mathbf{\varphi}\right)  \left(
D_{pn}^{1}D_{qk}^{1}-D_{pk}^{1}D_{qn}^{1}\right)  =\sum_{1\leq p,q,k\leq
2s}\left(  h_{2mkpq}\left(  \mathbf{\varphi}\right)  D_{qk}^{1}-h_{2mkqp}%
\left(  \mathbf{\varphi}\right)  D_{qk}^{1}\right)  D_{pn}^{1}%
\end{align}
and defining the coulomb matrix
\begin{equation}
J_{mp}\left(  \mathbf{\varphi}\right)  =\sum_{1\leq q,k\leq2s}h_{2mkpq}\left(
\mathbf{\varphi}\right)  D_{qk}^{1}%
\end{equation}
and exchange matrix
\begin{equation}
K_{mp}\left(  \mathbf{\varphi}\right)  =\sum_{1\leq q,k\leq2s}h_{2mkqp}\left(
\mathbf{\varphi}\right)  D_{qk}^{1}%
\end{equation}
leading to
\begin{equation}
\left(  L_{2}^{1}\left(  \mathbb{H}_{2}\left(  \mathbf{\varphi}\right)
\mathbb{D}^{2}\right)  \right)  _{mn}=\sum_{1\leq p\leq2s}\left\{
J_{mp}\left(  \mathbf{\varphi}\right)  -K_{mp}\left(  \mathbf{\varphi}\right)
\right\}  D_{pn}^{1}.
\end{equation}
The generalized Fock operator eq. $\left(  \ref{F1}\right)  $ can then be
expressed as
\begin{align}
F\left(  \mathbf{\varphi}\right)   &  =\sum_{1\leq m,n\leq2s}\left\{
\sum_{1\leq q\leq2s}\left\{  \left(  h_{1mq}\left(  \mathbf{\varphi}\right)
\right)  +J_{mq}\left(  \mathbf{\varphi}\right)  -K_{mq}\left(
\mathbf{\varphi}\right)  \right\}  D_{qn}^{1}\right\}  \left\vert \varphi
_{m}\right\rangle \left\langle \varphi_{n}\right\vert \nonumber\\
&  =\sum_{1\leq m,n\leq2s}\left(  \mathbf{f}\left(  \mathbf{\varphi}\right)
\mathbb{D}^{1}\right)  _{mn}\left\vert \varphi_{m}\right\rangle \left\langle
\varphi_{n}\right\vert ,
\end{align}
where we have introduced the standard Fock matrix $\mathbf{f}\left(
\mathbb{D}^{1},\mathbf{\varphi}\right)  $ with matrix elements
\begin{equation}
f\left(  \mathbf{\varphi}\right)  _{mq}=\left(  h_{1mq}\left(  \mathbf{\varphi
}\right)  \right)  +J_{mq}\left(  \mathbf{\varphi}\right)  -K_{mq}\left(
\mathbf{\varphi}\right)
\end{equation}
corresponding to the Fock operator
\begin{equation}
\mathfrak{f}\left(  \mathbf{\varphi}\right)  =\sum_{1\leq m,n\leq2s}f\left(
\mathbf{\varphi}\right)  _{mn}\left\vert \varphi_{m}\right\rangle \left\langle
\varphi_{n}\right\vert =\mathfrak{f}\left(  \mathbf{\varphi}\right)
^{\dagger},
\end{equation}
which is always self adjoint. In the IPS case the generalized Fock operator
thus corresponds to the product of the conventional Fock operator and the
FORDO i.e.
\begin{equation}
F\left(  \mathbf{\varphi}\right)  =\mathfrak{f}\left(  \mathbf{\varphi
}\right)  D^{1}.
\end{equation}


\section{Discussion\label{Discussion}}

We have extended the IPS one particle model of molecular quantum mechanics
described by IPS's to a more accurate correlated one particle model described
by AGP\ states for an even number of electrons and GAGP states for an odd
number. All of these states are determined by a set of one particle occupation
numbers and a set of CGSO's. In the IPS case the occupation numbers are one or
zero, while in the AGP case they must be at least doubly degenerate and for
the GAGP case there must be at most one non degenerate integer occupation
number and the rest at least doubly degenerate.

The occupation numbers and CGSO's totally determine the many electron state, a
necessary criterion for a one particle model. In the AGP\ case the GSO's
produce weighted Probability Distributions (PD's)\ that produce pictorial
models of a molecules that generalize the standard ones used in chemistry in
that they contain the effects of correlation. It is true that multi
configurational states also produce weighted Probability Distributions (PD's)
but they are not determined by them, in contrast to the AGP\ case.

Another important aspect of the IPS model is the role of energies of the one
particle orbitals, which are given by the eigenvalues of the Fock operator.
When these are combined with the shape of PD's one obtains a description of
chemical properties i.e. reactivity and reaction mechanisms. Physically these
energies, are the poles of a first order approximation to the electron
propagator and correspond to the ionization potentials of the electrons from
the occupied orbitals. Mathematically they reflect the effect of the
orthonormality constraints in the IP-SCF\ procedure i.e. they give the rate of
change of the total energy with respect to changes in the normalization of the
corresponding orbitals. The aufbau principle is also based on these
eigenvalues i.e. the optimized IPS\ is in general composed of the orbitals
corresponding to most negative ionization potentials.

These properties are generalized by the optimized AGP-Fock operator with two
important differences. The generalized Fock operator is not diagonal in the
basis of CGSO's and the diagonal elements give weighted ionization potentials.
This is a reflection of our contention that the proper entity to generalize to
correlated states in order to obtain orbital derivatives is the product
$\mathfrak{f}\left(  \mathbf{\varphi}\right)  D^{1}$. In the IPS\ case the
diagonal elements of $\mathfrak{f}\left(  \mathbf{\varphi}\right)  D^{1}$ and
$\mathfrak{f}\left(  \mathbf{\varphi}\right)  $ in the basis $\left\{
\left\vert \varphi_{j}\right\rangle \right\}  $, referring to the occupied
GSO's, are identical but this is not the case for AGP and MCSCF\ states.
Moreover the Fock operator $\mathfrak{f}\left(  \mathbf{\varphi}\right)  $ is
always self adjoint and both $\mathfrak{f}\left(  \mathbf{\varphi}\right)
D^{1}$ and $\mathfrak{f}\left(  \mathbf{\varphi}\right)  $ are diagonal in
this basis, thus the diagonal elements are the eigenvalues of these operators.
In the AGP and MCSCF cases the generalized Fock operator $F\left(
\mathbb{D}^{2},\mathbf{\varphi}\right)  $ is only self adjoint at stationary
points of the GSO variation and not diagonal.

Another aspect of the IPS theory that is generalized by optimized AGP's is the
Brillouin Condition
\begin{equation}
\left\langle \left.  g^{N}\right\vert Ha_{j}^{\dagger}a_{k}g^{N}\right\rangle
=0 \label{GBCD}%
\end{equation}
for $j=k\pm s$ and for all $j\neq k$ such that $c_{j}\neq c_{k}$. For non
degenerate AGP's i.e. where all the canonical coefficients $\left\{
c_{j}\right\}  $ are distinct the condition eq. $\left(  \ref{GBCD}\right)  $
holds for all $j\neq k$. In terms of CGSO's $\mathbf{\psi}$ the AGP-SORDM can
be expressed solely in terms of the geminal occupation numbers i.e.
$\mathbb{D}^{2}=\mathbb{D}^{2}\left(  \mathbf{n}\right)  $ and if these
occupation numbers are also optimized and non degenerate the Fock operator
$F\left(  \mathbf{n},\mathbf{\psi}\right)  $ can be diagonalized to obtain
eigenvalues that are again weighted ionization potentials associated with
eigenvectors $\left\{  \left\vert \varsigma_{p}\right\rangle \right\}  $, but
these eigenvectors are not NGSO's nor CGSO's of the optimized AGP state.

By using subgroups of the general linear group of one electron Hilbert space
we have shown how to separate the variation over AGP states into occupation
number, phase and NGSO variations. The phase dependence is a reflection of the
fact the AGP FORDO (determined by the occupation numbers NGSO's) does not
determine a unique AGP\ state, but only an equivalence class of states. By
optimizing over the phase for fixed reference FORDO's one obtains an explicit
Density Matrix Function (DMF)\ that can be used as a model functional for DMFT.

\section{Acknowledgements}

\appendix

\section{Example of Linear Fit\label{LinfitEX}}

As an illustration of the numerical procedure described in section
\ref{AGPoccGemocc} consider the change of AGP occupation numbers from
$\mathbf{N}_{0}$ to $\mathbf{N}_{0}+\delta\mathbf{N}_{1}+\delta\mathbf{N}%
_{2}$
\begin{equation}
\mathbf{n}_{2}=\mathbf{n}_{0}+\left[  d\mathcal{N}\left(  \mathbf{n}%
_{0}\right)  \right]  ^{-1}\left(  \delta\mathbf{N}_{1}\right)  +\left[
d\mathcal{N}\left(  \mathbf{n}_{0}+\left[  d\mathcal{N}\left(  \mathbf{n}%
_{0}\right)  \right]  ^{-1}\left(  \delta\mathbf{N}_{1}\right)  \right)
\right]  ^{-1}\left(  \left(  \delta\mathbf{N}_{2}\right)  \right)
\end{equation}
producing the change of geminal occupation numbers $\mathbf{n}_{2}%
-\mathbf{n}_{0},$ where
\begin{equation}
\mathbf{N}_{1}=\left(
\begin{array}
[c]{c}%
1\\
.9\\
.8\\
.1\\
.2\\
0
\end{array}
\right)  =\left(
\begin{array}
[c]{c}%
1\\
1\\
1\\
0\\
0\\
0
\end{array}
\right)  +\left(
\begin{array}
[c]{c}%
0\\
-.1\\
-.2\\
.1\\
.2\\
0
\end{array}
\right)  =\mathbf{N}_{0}+\delta\mathbf{N}_{1}%
\end{equation}
giving the first order changes in the geminal occupation numbers
\begin{equation}
\mathbf{n}_{1}=\mathcal{N}^{-1}\left(
\begin{array}
[c]{c}%
1\\
.9\\
.8\\
0\\
.2\\
.1
\end{array}
\right)  =\left(
\begin{array}
[c]{c}%
\frac{1}{3}\\
\frac{1}{3}\\
\frac{1}{3}\\
0\\
0\\
0
\end{array}
\right)  +\left[  d\mathcal{N}\left(
\begin{array}
[c]{c}%
\frac{1}{3}\\
\frac{1}{3}\\
\frac{1}{3}\\
0\\
0\\
0
\end{array}
\right)  \right]  ^{-1}\left(
\begin{array}
[c]{c}%
0\\
-.1\\
-.2\\
.1\\
.2\\
0
\end{array}
\right)  .
\end{equation}
The next incremental change of the AGP occupation numbers, $\delta
\mathbf{N}_{2}$, added to $\delta\mathbf{N}_{1}$ gives
\[
\mathbf{n}_{2}=\mathcal{N}^{-1}\left(
\begin{array}
[c]{c}%
.9\\
.8\\
.8\\
.1\\
.2\\
.2
\end{array}
\right)  =\left(
\begin{array}
[c]{c}%
\frac{1}{3}\\
\frac{1}{3}\\
\frac{1}{3}\\
0\\
0\\
0
\end{array}
\right)  +\left[  d\mathcal{N}\left(
\begin{array}
[c]{c}%
\frac{1}{3}\\
\frac{1}{3}\\
\frac{1}{3}\\
0\\
0\\
0
\end{array}
\right)  \right]  ^{-1}\left(
\begin{array}
[c]{c}%
0\\
-.1\\
-.2\\
.1\\
.2\\
0
\end{array}
\right)
\]%
\begin{equation}
+\left[  d\mathcal{N}\left(  \left(
\begin{array}
[c]{c}%
\frac{1}{3}\\
\frac{1}{3}\\
\frac{1}{3}\\
0\\
0\\
0
\end{array}
\right)  +\left[  d\mathcal{N}\left(
\begin{array}
[c]{c}%
\frac{1}{3}\\
\frac{1}{3}\\
\frac{1}{3}\\
0\\
0\\
0
\end{array}
\right)  \right]  ^{-1}\left(
\begin{array}
[c]{c}%
0\\
-.1\\
-.2\\
.1\\
.2\\
0
\end{array}
\right)  \right)  \right]  ^{-1}\left(
\begin{array}
[c]{c}%
-.1\\
-.1\\
0\\
.1\\
0\\
.1
\end{array}
\right)  .
\end{equation}
At each step one can calculate $\mathbf{N}$ from $\mathbf{n}$ to check whether
the increment $\delta\mathbf{N}$ was small enough.

\section{Energy Functionals\label{En_Funcs}}

The energy of a $2N$-electron state, represented by the multi electron
positive operator, $D^{2N}$, is given by
\begin{equation}
\mathcal{E}\left(  D^{2N}\right)  =\frac{1}{\operatorname{Tr}\left\{
D^{2N}\right\}  }\operatorname{Tr}\left\{  \left\{  \sum_{1\leq i,k\leq
2s}h_{1ik}a_{i}^{\dagger}a_{k}+\frac{1}{4}\sum_{1\leq i,j,k,l\leq2s}%
V_{2ijkl}a_{i}^{\dagger}a_{j}^{\dagger}a_{l}a_{k}\right\}  D^{2N}\right\}  ,
\label{en1}%
\end{equation}
where
\begin{align}
h_{1ik}  &  =\left\langle \varphi_{i}\left\vert h_{1}\varphi_{k}\right.
\right\rangle \nonumber\\
V_{2ijkl}  &  =\left\langle ij||kl\right\rangle =\left\langle \varphi
_{i}\varphi_{j}\left\vert h_{2}\varphi_{k}\varphi_{l}\right.  \right\rangle
\end{align}
and the second quantized operators $\left\{  \left.  a_{j}^{\dagger
}\right\vert 1\leq j\leq2s\right\}  $ refer to the orthonormal reference basis
$\left\{  \left.  \left\vert \varphi_{j}\right\rangle \right\vert 1\leq
j\leq2s\right\}  $. For pure states
\begin{equation}
D^{2N}=\frac{\left\vert \Psi\right\rangle \left\langle \Psi\right\vert
}{\left\langle \Psi\left\vert \Psi\right.  \right\rangle }\equiv
\frac{\left\vert \Psi\right\rangle \left\langle \Psi\right\vert }%
{\mathcal{S}\left(  \Psi\right)  },
\end{equation}
where $\left\vert \Psi\right\rangle $ is in general an \emph{unnormalized
}vector\emph{ }state belonging to the $2N$-electron Hilbert space
$\mathcal{H}^{2N}.$ The expression eq. $\left(  \ref{en1}\right)  $ can be
reformulated as
\begin{equation}
\mathcal{E}\left(  \mathbb{D}^{2}\right)  =\operatorname{Tr}\left\{
\mathbb{H}_{1}\mathbb{D}^{1}\right\}  +\frac{1}{4}\operatorname{Tr}\left\{
\mathbb{V}_{2}\mathbb{D}^{2}\right\}
\end{equation}
in terms of the first and second order reduced density matrices $\mathbb{D}%
^{1}$ and $\mathbb{D}^{2}$ defined by
\begin{equation}
D_{klij}^{2}=\frac{1}{\operatorname{Tr}\left\{  D^{2N}\right\}  }%
\operatorname{Tr}\left\{  a_{i}^{\dagger}a_{j}^{\dagger}a_{l}a_{k}%
D^{2N}\right\}  ;\quad1\leq i,j,k,l\leq2s
\end{equation}
and
\begin{equation}
\mathbb{D}^{1}=\left(  2N-1\right)  ^{-1}L_{2}^{1}\left(  \mathbb{D}%
^{2}\right)  .
\end{equation}
The matrices $\mathbb{H}_{1}$ and $\mathbb{V}_{2}$ have the matrix elements
$\left\{  h_{1ik}\right\}  $ and $\left\{  \left\langle ij||kl\right\rangle
\right\}  $ respectively. The energy expression eq. $\left(  \ref{en1}\right)
$ can be used to define a functional, $\mathcal{E}\left(  \mathbf{\alpha
}\right)  $, of orthonormal GSO's by considering a fixed reference state
$D_{0}^{2N}$ and one electron unitary transformations $U\left(  \mathbf{\alpha
}\right)  $ by setting%
\begin{equation}
\mathcal{E}\left(  \mathbf{\alpha}\right)  =\frac{1}{\operatorname{Tr}\left\{
D_{0}^{2N}\right\}  }\operatorname{Tr}\left\{  \left\{  \sum_{1\leq i,k\leq
2s}h_{1ik}a_{i}^{\dagger}a_{k}+\frac{1}{4}\sum_{1\leq i,j,k,l\leq2s}%
V_{2ijkl}a_{i}^{\dagger}a_{j}^{\dagger}a_{l}a_{k}\right\}  U\left(
\mathbf{\alpha}\right)  ^{\dagger}D_{0}^{2N}U\left(  \mathbf{\alpha}\right)
\right\}  ,
\end{equation}
where
\begin{equation}
U\left(  \mathbf{\alpha}\right)  =\exp\left\{  i\sum_{1\leq p,q\leq2s}%
\alpha_{pq}a_{p}^{\dagger}a_{q}\right\}  ;\qquad\alpha_{pq}=\alpha_{qp}^{\ast
}\quad1\leq p,q\leq2s.
\end{equation}
This leads to
\begin{align}
\mathcal{E}\left(  \mathbf{\alpha}\right)   &  =\frac{1}{\operatorname{Tr}%
\left\{  D_{0}^{2N}\right\}  }\operatorname{Tr}\left\{  \left\{  \sum_{1\leq
i,k\leq2s}h_{1ik}U\left(  \mathbf{\alpha}\right)  a_{i}^{\dagger}U\left(
\mathbf{\alpha}\right)  ^{\dagger}U\left(  \mathbf{\alpha}\right)
a_{k}U\left(  \mathbf{\alpha}\right)  ^{\dagger}\right.  \right. \nonumber\\
&  +\left.  \left.  \frac{1}{4}\sum_{1\leq i,j,k,l\leq2s}V_{2ijkl}U\left(
\mathbf{\alpha}\right)  a_{i}^{\dagger}U\left(  \mathbf{\alpha}\right)
^{\dagger}U\left(  \mathbf{\alpha}\right)  a_{j}^{\dagger}U\left(
\mathbf{\alpha}\right)  ^{\dagger}U\left(  \mathbf{\alpha}\right)
a_{l}U\left(  \mathbf{\alpha}\right)  ^{\dagger}U\left(  \mathbf{\alpha
}\right)  a_{k}U\left(  \mathbf{\alpha}\right)  ^{\dagger}D_{0}^{2N}\right\}
\right\}  ,
\end{align}
which can be rearranged to give
\begin{align}
\mathcal{E}\left(  \mathbf{\alpha}\right)   &  =\frac{1}{\operatorname{Tr}%
\left\{  D_{0}^{2N}\right\}  }\operatorname{Tr}\left\{  \left\{  \sum_{1\leq
i,k\leq2s}\left\langle \left.  \varphi_{i}\left(  \mathbf{\alpha}\right)
\right\vert h_{1}\varphi_{k}\left(  \mathbf{\alpha}\right)  \right\rangle
a_{i}^{\dagger}a_{k}\right.  \right. \nonumber\\
&  \qquad\qquad\qquad+\left.  \left.  \frac{1}{4}\sum_{1\leq i,j,k,l\leq
2s}\left\langle \left.  \varphi_{i}\left(  \mathbf{\alpha}\right)  \varphi
_{j}\left(  \mathbf{\alpha}\right)  \right\vert h_{2}\varphi_{k}\left(
\mathbf{\alpha}\right)  \varphi_{l}\left(  \mathbf{\alpha}\right)
\right\rangle a_{i}^{\dagger}a_{j}^{\dagger}a_{l}a_{k}\right\}  D_{0}%
^{2N}\right\}
\end{align}
and
\begin{equation}
\mathcal{E}\left(  \mathbf{\alpha}\right)  =\operatorname{Tr}\left\{
\mathbb{H}_{1}\left(  \mathbf{\varphi}\left(  \mathbf{\alpha}\right)  \right)
\mathbb{D}_{0}^{1}\right\}  +\frac{1}{4}\operatorname{Tr}\left\{
\mathbb{V}_{2}\left(  \mathbf{\varphi}\left(  \mathbf{\alpha}\right)  \right)
\mathbb{D}_{0}^{2}\right\}  , \label{enmat}%
\end{equation}
where
\begin{equation}
\left\vert \varphi_{i}\left(  \mathbf{\alpha}\right)  \right\rangle =U\left(
\mathbf{\alpha}\right)  \left\vert \varphi_{i}\right\rangle .
\end{equation}
The expression eq. $\left(  \ref{enmat}\right)  $ utilizes the GSO dependent
matrices $\mathbb{H}_{1}\left(  \mathbf{\varphi}\left(  \mathbf{\alpha
}\right)  \right)  $ and $\mathbb{V}_{2}\left(  \mathbf{\varphi}\left(
\mathbf{\alpha}\right)  \right)  $ whose matrix elements are
\begin{subequations}
\begin{align}
\mathbb{H}_{1}\left(  \mathbf{\varphi}\left(  \mathbf{\alpha}\right)  \right)
&  \leftrightarrow\left\{  \left\langle \left.  \varphi_{i}\left(
\mathbf{\alpha}\right)  \right\vert h_{1}\varphi_{k}\left(  \mathbf{\alpha
}\right)  \right\rangle \right\} \\
\mathbb{V}_{2}\left(  \mathbf{\varphi}\left(  \mathbf{\alpha}\right)  \right)
&  \leftrightarrow\left\{  \left\langle \left.  \varphi_{i}\left(
\mathbf{\alpha}\right)  \varphi_{j}\left(  \mathbf{\alpha}\right)  \right\vert
h_{2}\varphi_{k}\left(  \mathbf{\alpha}\right)  \varphi_{l}\left(
\mathbf{\alpha}\right)  \right\rangle \right\}  .
\end{align}
Each new non zero value of $\mathbf{\alpha}$ determines a new state
$D^{2N}\left(  \mathbf{\alpha}\right)  =U\left(  \mathbf{\alpha}\right)
^{\dagger}D_{0}^{2N}U\left(  \mathbf{\alpha}\right)  $ and new first and
second order reduced density matrices given by
\end{subequations}
\begin{align}
D^{2}\left(  \mathbf{\alpha}\right)  _{klij}  &  =\frac{1}{\operatorname{Tr}%
\left\{  D_{0}^{2N}\right\}  }\operatorname{Tr}\left\{  U\left(
\mathbf{\alpha}\right)  a_{i}^{\dagger}U\left(  \mathbf{\alpha}\right)
^{\dagger}U\left(  \mathbf{\alpha}\right)  a_{j}^{\dagger}U\left(
\mathbf{\alpha}\right)  ^{\dagger}U\left(  \mathbf{\alpha}\right)
a_{l}U\left(  \mathbf{\alpha}\right)  ^{\dagger}U\left(  \mathbf{\alpha
}\right)  a_{k}U\left(  \mathbf{\alpha}\right)  ^{\dagger}D_{0}^{2N}\right\}
\nonumber\\
&  =\frac{1}{\operatorname{Tr}\left\{  D_{0}^{2N}\right\}  }\operatorname{Tr}%
\left\{  a_{i}^{\dagger}a_{j}^{\dagger}a_{l}a_{k}U\left(  \mathbf{\alpha
}\right)  ^{\dagger}D_{0}^{2N}U\left(  \mathbf{\alpha}\right)  \right\}
;\quad1\leq i,j,k,l\leq2s
\end{align}
and
\begin{align}
D^{1}\left(  \mathbf{\alpha}\right)  _{ji}  &  =\frac{1}{\operatorname{Tr}%
\left\{  D_{0}^{2N}\right\}  }\operatorname{Tr}\left\{  U\left(
\mathbf{\alpha}\right)  a_{i}^{\dagger}U\left(  \mathbf{\alpha}\right)
^{\dagger}U\left(  \mathbf{\alpha}\right)  a_{j}U\left(  \mathbf{\alpha
}\right)  ^{\dagger}D_{0}^{2N}\right\} \nonumber\\
&  =\frac{1}{\operatorname{Tr}\left\{  D_{0}^{2N}\right\}  }\operatorname{Tr}%
\left\{  a_{i}^{\dagger}a_{j}U\left(  \mathbf{\alpha}\right)  ^{\dagger}%
D_{0}^{2N}U\left(  \mathbf{\alpha}\right)  \right\}  .
\end{align}


The energy functional eq. $\left(  \ref{enmat}\right)  $ could be extended to
a full non redundent energy functional by varying the reference
SORDM\ $\mathbb{D}_{0}^{2}$ in such a way as not to duplicate GSO changes.
However as the $N$-representability problem is still unsolved so the domain of
variation of $\mathbb{D}_{0}^{2}$ is in general unknown, except in the AGP
state case, which we study in this paper, where $N$-representable SORDM's can
be explicitly parameterized in terms of the occupation numbers $\left\{
\left.  n_{j}\right\vert 1\leq j\leq s\right\}  $. The energy functionals eqs.
$\left(  \ref{RDO_en}\right)  $ and $\left(  \ref{RDO_en1}\right)  $ are of
course also valid for SORDM's explicitly constructed from multi
configurational states, i.e. in terms of CI\ coefficients.

\section{Orbital Derivatives and Derivation of the Fock Operator}

\subsection{Direct Variation of GSO's}

In order to obtain derivative of terms in the energy expression it helpful to
utilize matrix elements wrt non antisymmetric states
\begin{equation}
\left\vert \varphi_{j}\varphi_{k}\right)  =\left\vert \varphi_{j}\right)
\otimes\left\vert \varphi_{k}\right)  ,
\end{equation}
whose amplitudes are
\begin{equation}
\left(  \left.  \mathbf{r}_{1}\mathbf{,\xi}_{1}\mathbf{r}_{2}\mathbf{,\xi}%
_{2}\right\vert \varphi_{j}\varphi_{k}\right)  =\varphi_{j}\left(
\mathbf{r}_{1}\mathbf{,\xi}_{1}\right)  \varphi_{k}\left(  \mathbf{r}%
_{2}\mathbf{,\xi}_{2}\right)  ,
\end{equation}
where
\begin{equation}
\varphi_{j}\left(  \mathbf{r}_{1}\mathbf{,\xi}_{1}\right)  =\left(  \left.
\mathbf{r}_{1}\mathbf{,\xi}_{1}\right\vert \varphi_{j}\right)  =\left\langle
\left.  \mathbf{r}_{1}\mathbf{,\xi}_{1}\right\vert \varphi_{j}\right\rangle .
\end{equation}
We will consider the energy in terms of the SORDM\ $\mathbb{D}^{2}$ eq.
$\left(  \ref{sqe3}\right)  $ and GSO's $\mathbf{\varphi}$ given by
\begin{equation}
\mathcal{E}\left(  \mathbb{D}^{2},\mathbf{\varphi}\right)  =\operatorname{Tr}%
\left\{  \mathbb{H}_{1}\left(  \mathbf{\varphi}\right)  \mathbb{D}%
^{1}\right\}  +\frac{1}{4}\operatorname{Tr}\left\{  \mathbb{V}_{2}\left(
\mathbf{\varphi}\right)  \mathbb{D}^{2}\right\}  , \label{e1}%
\end{equation}
and we will concentrate on the direct variation of the GSO's. The first order
change in the energy due to changes $\left\{  \delta\mathbf{\varphi}^{\ast
},\delta\mathbf{\varphi}\right\}  $ in the GSO's is given by the Taylors
series as
\begin{align}
\delta^{1}\mathcal{E}\left(  \mathbb{D}^{2},\mathbf{\varphi}\right)   &
=\left(  \partial_{\mathbf{\varphi}^{\ast}}^{1}\mathcal{E}\left(
\mathbb{D}^{2},\mathbf{\varphi}\right)  ,\partial_{\mathbf{\varphi}}%
^{1}\mathcal{E}\left(  \mathbb{D}^{2},\mathbf{\varphi}\right)  \right)
\left(
\begin{array}
[c]{c}%
\delta\mathbf{\varphi}^{\ast}\\
\delta\mathbf{\varphi}%
\end{array}
\right) \nonumber\\
&  =\left(  \partial_{\mathbf{\varphi}^{\ast}}^{1}\mathcal{E}\left(
\mathbb{D}^{2},\mathbf{\varphi}\right)  ,\partial_{\mathbf{\varphi}^{\ast}%
}^{1}\mathcal{E}\left(  \mathbb{D}^{2},\mathbf{\varphi}\right)  ^{\ast
}\right)  \left(
\begin{array}
[c]{c}%
\delta\mathbf{\varphi}^{\ast}\\
\delta\mathbf{\varphi}%
\end{array}
\right) \nonumber\\
&  =2\operatorname{Re}\left\{  \partial_{\mathbf{\varphi}^{\ast}}%
^{1}\mathcal{E}\left(  \mathbb{D}^{2},\mathbf{\varphi}\right)  ^{\dagger
}\delta\mathbf{\varphi}^{\ast}\right\}  ,
\end{align}
where
\begin{equation}
\partial_{\varphi_{\mu}^{\ast}}^{1}\mathcal{E}\left(  \mathbb{D}%
^{2},\mathbf{\varphi}\right)  =\frac{\delta\mathcal{E}\left(  \mathbb{D}%
^{2},\mathbf{\varphi}\right)  }{\delta\varphi_{\mu}^{\ast}}.
\end{equation}
It is clear that
\begin{equation}
\partial_{\varphi_{\mu}^{\ast}}^{1}\mathcal{E}\left(  \mathbb{D}%
^{2},\mathbf{\varphi}\right)  =\left[  \partial_{\varphi_{\mu}}^{1}%
\mathcal{E}\left(  \mathbb{D}^{2},\mathbf{\varphi}\right)  \right]  ^{\ast}%
\end{equation}
as $\delta^{1}\mathcal{E}\left(  \mathbb{D}^{2},\mathbf{\varphi}\right)  $
must be real valued for all $\left(  \delta\mathbf{\varphi}^{\ast}%
,\delta\mathbf{\varphi}\right)  $. In order to find the explicit expressions
for the derivatives of the energy one observes that:

\begin{enumerate}
\item The one electron functional derivatives are given by
\begin{equation}
\frac{\delta}{\delta\varphi_{\mu}^{\ast}}\left\langle \left.  \varphi
_{a}\right\vert h_{1}\varphi_{b}\right\rangle =h_{1}\left\vert \varphi
_{b}\right\rangle \delta_{a\mu}=\left\{  h_{1}\left\vert \varphi
_{b}\right\rangle \left\langle \varphi_{a}\right\vert \right\}  \left\vert
\varphi_{\mu}\right\rangle . \label{1partfuncder_o}%
\end{equation}
After noting that
\begin{equation}
h_{1}\left\vert \varphi_{b}\right\rangle \left\langle \varphi_{a}\right\vert
=\sum_{1\leq j,k\leq2s}h_{1jk}\left\vert \varphi_{j}\right\rangle \left\langle
\left.  \varphi_{k}\right\vert \varphi_{b}\right\rangle \left\langle
\varphi_{a}\right\vert =\sum_{1\leq j\leq2s}h_{1jb}\left\vert \varphi
_{j}\right\rangle \left\langle \varphi_{a}\right\vert
\end{equation}
eq. $\left(  \ref{1partfuncder_o}\right)  $ can be expanded as
\begin{equation}
\frac{\delta}{\delta\varphi_{\mu}^{\ast}}\left\langle \left.  \varphi
_{a}\right\vert h_{1}\varphi_{b}\right\rangle =\left\{  \sum_{1\leq m,n\leq
2s}h_{1mb}\delta_{an}\left\vert \varphi_{m}\right\rangle \left\langle
\varphi_{n}\right\vert \right\}  \left\vert \varphi_{\mu}\right\rangle .
\label{1partfuncder_ex}%
\end{equation}


\item The two electron functional derivatives are
\begin{align}
&  \frac{\delta}{\delta\varphi_{\mu}^{\ast}}\left(  \left.  \varphi_{a}%
\varphi_{b}\right\vert h_{2}\varphi_{c}\varphi_{d}\right)  =\left(  \left.
\varphi_{b}\right\vert h_{2}\varphi_{d}\right)  \left\vert \varphi
_{c}\right\rangle \delta_{a\mu}+\left(  \left.  \varphi_{a}\right\vert
h_{2}\varphi_{c}\right)  \left\vert \varphi_{d}\right\rangle \delta_{b\mu
}\nonumber\\
&  \qquad\qquad\qquad\left.  =\right.  \left\{  \left(  \left.  \varphi
_{b}\right\vert h_{2}\varphi_{d}\right)  \left\vert \varphi_{c}\right\rangle
\left\langle \varphi_{a}\right\vert +\left(  \left.  \varphi_{a}\right\vert
h_{2}\varphi_{c}\right)  \left\vert \varphi_{d}\right\rangle \left\langle
\varphi_{b}\right\vert \right\}  \left\vert \varphi_{\mu}\right\rangle ,
\label{2partfuncder_o}%
\end{align}
where the operators $\left\{  \left(  \left.  \varphi_{b}\right\vert
h_{2}\varphi_{d}\right)  ;1\leq b,d\leq2s\right\}  $ are defined by
\begin{equation}
\left(  \left.  \varphi_{b}\right\vert h_{2}\varphi_{d}\right)  =\sum_{1\leq
m,n\leq2s}\left(  \left.  \varphi_{m}\varphi_{b}\right\vert h_{2}\varphi
_{n}\varphi_{d}\right)  \left\vert \varphi_{m}\right\rangle \left\langle
\varphi_{n}\right\vert .
\end{equation}

\end{enumerate}

Eq. $\left(  \ref{1partfuncder_ex}\right)  $ can be further expanded to give
\begin{align}
\frac{\delta}{\delta\varphi_{\mu}^{\ast}}\left(  \left.  \varphi_{a}%
\varphi_{b}\right\vert h_{2}\varphi_{c}\varphi_{d}\right)   &  =\left\{
\sum_{1\leq k\leq2s}\left\{  h_{2kbcd}\left\vert \varphi_{k}\right\rangle
\left\langle \varphi_{a}\right\vert +h_{2kadc}\left\vert \varphi
_{k}\right\rangle \left\langle \varphi_{b}\right\vert \right\}  \right\}
\left\vert \varphi_{\mu}\right\rangle \nonumber\\
&  =\left\{  \sum_{1\leq m,n\leq2s}\left\{  h_{2mbcd}\delta_{an}%
+h_{2madc}\delta_{bn}\right\}  \left\vert \varphi_{m}\right\rangle
\left\langle \varphi_{n}\right\vert \right\}  \left\vert \varphi_{\mu
}\right\rangle . \label{2partfuncder_ex}%
\end{align}


The derivatives $\left\{  \frac{\delta\mathcal{E}\left(  \mathbb{D}%
^{2},\mathbf{\varphi}\right)  }{\delta\varphi_{\mu}^{\ast}}\right\}  $ can be
evaluated by using eqs. $\left(  \ref{1partfuncder_ex}\right)  $ and $\left(
\ref{2partfuncder_ex}\right)  $ to give
\begin{align}
&  \frac{\delta}{\delta\varphi_{\mu}^{\ast}}\left\{  \sum_{1\leq i,k\leq
2s}\left\langle \left.  \varphi_{i}\right\vert h_{1}\varphi_{k}\right\rangle
D_{ki}^{1}+\frac{1}{4}\sum_{1\leq i,j,k,l\leq2s}\left\{  \left(  \left.
\varphi_{i}\varphi_{j}\right\vert h_{2}\varphi_{k}\varphi_{l}\right)  -\left(
\left.  \varphi_{i}\varphi_{j}\right\vert h_{2}\varphi_{l}\varphi_{k}\right)
\right\}  D_{klij}^{2}\right\} \nonumber\\
&  =\sum_{1\leq m,n\leq2s}\left\{  \left\{  \sum_{1\leq k\leq2s}h_{1mk}%
D_{kn}^{1}\left\vert \varphi_{m}\right\rangle \left\langle \varphi
_{n}\right\vert \right.  \right. \nonumber\\
&  \qquad\qquad\qquad\qquad\qquad+\left.  \left.  \frac{1}{2}\sum_{1\leq
j,k,l\leq2s}\left\{  h_{2mjkl}D_{klnj}^{2}-h_{2mjlk}D_{klnj}^{2}\right\}
\right\}  \left\vert \varphi_{m}\right\rangle \left\langle \varphi
_{n}\right\vert \right\}  \left\vert \varphi_{\mu}\right\rangle ,
\end{align}
which can be written as
\begin{equation}
\frac{\delta\mathcal{E}\left(  \mathbb{D}^{2},\mathbf{\varphi}\right)
}{\delta\varphi_{\mu}^{\ast}}=\left\{  \sum_{1\leq m,n\leq2s}\left\{  \left(
\mathbb{H}_{1}\left(  \mathbf{\varphi}\right)  \mathbb{D}^{1}\right)
_{mn}+\left(  L_{2}^{1}\left(  \mathbb{H}_{2}\left(  \mathbf{\varphi}\right)
\mathbb{D}^{2}\right)  \right)  _{mn}\right\}  \left\vert \varphi
_{m}\right\rangle \left\langle \varphi_{n}\right\vert \right\}  \left\vert
\varphi_{\mu}\right\rangle ,
\end{equation}
where $\mathbb{H}_{2}\left(  \mathbf{\varphi}\right)  $ is the \emph{non
}antisymmetric matrix
\begin{equation}
\mathbb{H}_{2}\left(  \mathbf{\varphi}\right)  \leftrightarrow\left\{
h_{2ijkl}\right\}
\end{equation}
and
\begin{equation}
V_{2ijkl}\left(  \mathbf{\varphi}\right)  =\left(  h_{2ijkl}\left(
\mathbf{\varphi}\right)  -h_{2ijlk}\left(  \mathbf{\varphi}\right)  \right)  .
\end{equation}
This allows us to define a generalized Fock operator
\begin{subequations}
\label{Fock_Op}%
\begin{align}
F\left(  \mathbb{D}^{2},\mathbf{\varphi}\right)   &  =\sum_{1\leq m,n\leq
2s}\left\{  \left(  \mathbb{H}_{1}\left(  \mathbf{\varphi}\right)
\mathbb{D}^{1}\right)  _{mn}+\frac{1}{2}\left(  L_{2}^{1}\left(
\mathbb{V}_{2}\left(  \mathbf{\varphi}\right)  \mathbb{D}^{2}\right)  \right)
_{mn}\right\}  \left\vert \varphi_{m}\right\rangle \left\langle \varphi
_{n}\right\vert \label{Fock_Op_a}\\
&  =\sum_{1\leq m,n\leq2s}\left\{  \left(  \mathbb{H}_{1}\left(
\mathbf{\varphi}\right)  \mathbb{D}^{1}\right)  _{mn}+\left(  L_{2}^{1}\left(
\mathbb{H}_{2}\left(  \mathbf{\varphi}\right)  \mathbb{D}^{2}\right)  \right)
_{mn}\right\}  \left\vert \varphi_{m}\right\rangle \left\langle \varphi
_{n}\right\vert , \label{Fock_Op_b}%
\end{align}
that produces the functional derivatives
\end{subequations}
\begin{equation}
\frac{\delta\mathcal{E}\left(  \mathbb{D}^{2},\mathbf{\varphi}\right)
}{\delta\varphi_{\mu}^{\ast}}=F\left(  \mathbb{D}^{2},\mathbf{\varphi}\right)
\left\vert \varphi_{\mu}\right\rangle ;\quad1\leq\mu\leq2s.
\label{FuncDerFock}%
\end{equation}
As
\begin{equation}
L_{2}^{1}\left(  \mathbb{A}^{\dagger}\right)  =L_{2}^{1}\left(  \mathbb{A}%
\right)  ^{\dagger}.
\end{equation}
the adjoint of the Fock operator eq. $\left(  \ref{Fock_Op}\right)  $ is given
by
\begin{equation}
F\left(  \mathbb{D}^{2},\mathbf{\varphi}\right)  ^{\dagger}=\sum_{1\leq
m,n\leq2s}\left\{  \left(  \mathbb{D}^{1}\mathbb{H}_{1}\left(  \mathbf{\varphi
}\right)  \right)  _{mn}+\left(  L_{2}^{1}\left(  \mathbb{D}^{2}\mathbb{H}%
_{2}\left(  \mathbf{\varphi}\right)  \right)  \right)  _{mn}\right\}
\left\vert \varphi_{m}\right\rangle \left\langle \varphi_{n}\right\vert .
\label{Fock_Op_adj}%
\end{equation}
In general the generalized Fock operator is not self adjoint i.e.
\begin{equation}
F\left(  \mathbb{D}^{2},\mathbf{\varphi}\right)  ^{\dagger}\neq F\left(
\mathbb{D}^{2},\mathbf{\varphi}\right)  .
\end{equation}


\subsection{Matrix Elements of $F\left(  \mathbf{\varphi}\right)  $ in Terms
of Second Quantized Operators\label{SQ_Fock}}

In order to obtain an expression for the matrix elements $\left\{  \left.
F\left(  \mathbb{D}^{2},\mathbf{\varphi}\right)  _{jk}\right\vert 1\leq
j,k\leq2s\right\}  $ in terms of the second quantized operators $\left\{
\left.  a_{j}^{\dagger}\right\vert 1\leq j,k\leq2s\right\}  ,$ that refer to
the GSO's $\left\{  \left.  \left\vert \varphi_{j}\right\rangle \right\vert
1\leq j,k\leq2s\right\}  $, we consider the expectation value of the
commutators $\left\{  \left.  a_{p}^{\dagger}\left[  a_{q},H\right]
\right\vert 1\leq j,k\leq2s\right\}  ,$ where
\begin{equation}
H=h_{1}+\frac{1}{4}V_{2}=\sum_{1\leq i,k\leq2s}h_{1ik}a_{i}^{\dagger}%
a_{k}+\frac{1}{4}\sum_{1\leq i,j,k,l\leq2s}V_{2ijkl}a_{i}^{\dagger}%
a_{j}^{\dagger}a_{l}a_{k},
\end{equation}%
\begin{equation}
h_{1}=\sum_{1\leq i,k\leq2s}h_{1ik}a_{i}^{\dagger}a_{k}%
\end{equation}
and
\begin{equation}
V_{2}=\sum_{1\leq i,j,k,l\leq2s}V_{2ijkl}a_{i}^{\dagger}a_{j}^{\dagger}%
a_{l}a_{k}.
\end{equation}
The one electron part of $H$ gives
\begin{equation}
a_{p}^{\dagger}\sum_{1\leq i,k\leq2s}h_{1ik}\left[  a_{q},a_{i}^{\dagger}%
a_{k}\right]  =\sum_{1\leq k\leq2s}h_{1qk}a_{p}^{\dagger}a_{k},
\end{equation}
while the two electron interaction part gives
\begin{equation}
a_{p}^{\dagger}\sum_{1\leq i,j,k,l\leq2s}V_{2ijkl}\left[  a_{q},a_{i}%
^{\dagger}a_{j}^{\dagger}a_{l}a_{k}\right]  =2\sum_{1\leq j,k,l\leq
2s}V_{2qjkl}a_{p}^{\dagger}a_{j}^{\dagger}a_{l}a_{k}.
\end{equation}
The expectation values of the commutators can be expanded as
\begin{align}
&  \frac{1}{\operatorname{Tr}\left\{  D^{2N}\right\}  }\operatorname{Tr}%
\left\{  a_{p}^{\dagger}\left[  a_{q},H\right]  D^{2N}\right\}  =\nonumber\\
&  \qquad\qquad\qquad\frac{1}{\operatorname{Tr}\left\{  D^{2N}\right\}
}\operatorname{Tr}\left\{  \left\{  \sum_{1\leq k\leq2s}h_{1qk}a_{p}^{\dagger
}a_{k}+\frac{1}{2}\sum_{1\leq j,k,l\leq2s}V_{2qjkl}a_{p}^{\dagger}%
a_{j}^{\dagger}a_{l}a_{k}\right\}  D^{2N}\right\}  ,
\end{align}
showing that
\begin{equation}
F\left(  \mathbb{D}^{2},\mathbf{\varphi}\right)  _{qp}=\frac{1}%
{\operatorname{Tr}\left\{  D^{2N}\right\}  }\operatorname{Tr}\left\{
a_{p}^{\dagger}\left[  a_{q},H\right]  D^{2N}\right\}  =\operatorname{Tr}%
\left\{  \left(  \mathbb{H}_{1}\left(  \mathbf{\varphi}\right)  \mathbb{D}%
^{1}\right)  _{qp}+\frac{1}{2}\left(  L_{2}^{1}\left(  \mathbb{V}_{2}\left(
\mathbf{\varphi}\right)  \mathbb{D}^{2}\right)  \right)  _{qp}\right\}  .
\label{Fockelem}%
\end{equation}


\section{Variation of Group Parameters\label{VarGroupPar}}

The energy of a general $2N$-electron state can be expressed (see eqs.
$\left(  \ref{Rham_1elec}\right)  $ and $\left(  \ref{RDO_en}\right)  $) as
\begin{equation}
\mathcal{E}\left(  \mathbb{D}^{2},\mathbf{\alpha}\right)  =\operatorname{Tr}%
\left\{  \left\{  \sum_{1\leq i,j\leq2s}h_{1ij}a_{i}^{\dagger}a_{j}+\frac
{1}{4}\sum_{1\leq i,j,k,l\leq2s}V_{2ijkl}a_{i}^{\dagger}a_{j}^{\dagger}%
a_{l}a_{k}\right\}  U\left(  \mathbf{\alpha}\right)  ^{\dagger}D^{2N}U\left(
\mathbf{\alpha}\right)  \right\}  ,
\end{equation}
where
\begin{equation}
U\left(  \mathbf{\alpha}\right)  =\exp\left\{  i\sum_{1\leq p,q\leq2s}%
\alpha_{pq}a_{p}^{\dagger}a_{q}\right\}  ;\qquad\alpha_{pq}=\alpha_{qp}^{\ast
}\quad1\leq p,q\leq2s.
\end{equation}
in terms of the unitary group $U\left(  \mathcal{H}^{1}\right)  $ of
transformations $\left\{  U\left(  \mathbf{\alpha}\right)  \right\}  $ of the
Hilbert space of one electron states $\mathcal{H}^{1}$ and SORDM's
$\mathbb{D}^{2}$. The Taylor series expansion gives
\begin{align}
&  \mathcal{E}\left(  \mathbb{D}^{2},\mathbf{\alpha}\right)  \simeq
\operatorname{Tr}\left\{  \sum_{1\leq i,j\leq2s}h_{1ij}a_{i}^{\dagger}%
a_{j}+\frac{1}{4}\sum_{1\leq i,j,k,l\leq2s}V_{2ijkl}\operatorname{Tr}%
a_{i}^{\dagger}a_{j}^{\dagger}a_{l}a_{k}\right. \nonumber\\
&  \qquad\qquad\qquad\qquad\qquad\left.  \left\{  1-i\sum_{1\leq p,q\leq
2s}\alpha_{pq}a_{p}^{\dagger}a_{q}\right\}  D^{2N}\left\{  1+i\sum_{1\leq
p,q\leq2s}\alpha_{pq}a_{p}^{\dagger}a_{q}\right\}  \right\} \nonumber\\
&  \simeq\mathcal{E}\left(  \mathbb{D}^{2},\mathbf{0}\right)  -i\sum_{1\leq
p,q\leq2s}\operatorname{Tr}\left\{  \left\{  \sum_{1\leq i,j\leq2s}%
h_{1ij}\left[  a_{i}^{\dagger}a_{j},a_{p}^{\dagger}a_{q}\right]  \right.
\right. \nonumber\\
&  \qquad\qquad\qquad\qquad\qquad\qquad\left.  \left.  +\frac{1}{4}\sum_{1\leq
i,j,k,l\leq2s}V_{2ijkl}\left[  a_{i}^{\dagger}a_{j}^{\dagger}a_{l}a_{k,}%
a_{p}^{\dagger}a_{q}\right]  \right\}  D^{2N}\right\}  \alpha_{pq}.
\end{align}
Thus showing that the first order change of the energy functional wrt
\textbf{$\alpha$} evaluated around \textbf{$\alpha$ }$=\mathbf{0}$ is
\begin{align}
\delta^{1}\mathcal{E}\left(  \mathbb{D}^{2},\mathbf{0}\right)   &
=-i\sum_{1\leq p,q\leq2s}\operatorname{Tr}\left\{  \left\{  \sum_{1\leq
i,j\leq2s}h_{1ij}\left[  a_{i}^{\dagger}a_{j},a_{p}^{\dagger}a_{q}\right]
\right.  \right. \nonumber\\
&  \qquad\qquad+\left.  \left.  \frac{1}{4}\sum_{1\leq i,j,k,l\leq2s}%
V_{2ijkl}\left[  a_{i}^{\dagger}a_{j}^{\dagger}a_{l}a_{k,}a_{p}^{\dagger}%
a_{q}\right]  \right\}  \left.  D^{2N}\right\}  \right\}  \alpha_{pq},
\end{align}
showing that the derivatives wrt $\left\{  \left.  \alpha_{pq}\right\vert
1\leq p,q\leq2s\right\}  $ are
\begin{align}
\frac{\partial\mathcal{E}\left(  \mathbb{D}^{2},\mathbf{0}\right)  }%
{\partial\alpha_{pq}}  &  =-i\operatorname{Tr}\left\{  \left\{  \sum_{1\leq
i,j\leq2s}h_{1ij}\left[  a_{i}^{\dagger}a_{j},a_{p}^{\dagger}a_{q}\right]
\right.  \right. \nonumber\\
&  \qquad\qquad+\left.  \left.  \frac{1}{4}\sum_{1\leq i,j,k,l\leq2s}%
V_{2ijkl}\left[  a_{i}^{\dagger}a_{j}^{\dagger}a_{l}a_{k},a_{p}^{\dagger}%
a_{q}\right]  \right\}  D^{2N}\right\}  .
\end{align}
The commutators give
\begin{subequations}
\begin{align}
\left[  a_{i}^{\dagger}a_{j},a_{p}^{\dagger}a_{q}\right]   &  =a_{i}^{\dagger
}a_{q}\delta_{pj}-a_{p}^{\dagger}a_{j}\delta_{pj}\\
\left[  a_{i}^{\dagger}a_{j}^{\dagger}a_{l}a_{k},a_{p}^{\dagger}a_{q}\right]
&  =a_{i}^{\dagger}a_{j}^{\dagger}a_{q}a_{k}\delta_{pl}-a_{p}^{\dagger}%
a_{j}^{\dagger}a_{l}a_{k}\delta_{iq}-a_{i}^{\dagger}a_{p}^{\dagger}a_{l}%
a_{k}\delta_{jq}+a_{i}^{\dagger}a_{j}^{\dagger}a_{l}a_{q}\delta_{pk}%
\end{align}
leading to
\end{subequations}
\begin{align}
&  \frac{\partial\mathcal{E}\left(  \mathbb{D}^{2},\mathbf{0}\right)
}{\partial\alpha_{pq}}=-i\operatorname{Tr}\left\{  \left\{  \sum_{1\leq
i,j\leq2s}h_{1ij}\left(  a_{i}^{\dagger}a_{q}\delta_{pj}-a_{p}^{\dagger}%
a_{j}\delta_{iq}\right)  \right.  \right.  +\nonumber\\
&  \left.  \left.  \frac{1}{4}\sum_{1\leq i,j,k,l\leq2s}V_{2ijkl}\left(
a_{i}^{\dagger}a_{j}^{\dagger}a_{q}a_{k}\delta_{pl}-a_{p}^{\dagger}%
a_{j}^{\dagger}a_{l}a_{k}\delta_{iq}-a_{i}^{\dagger}a_{p}^{\dagger}a_{l}%
a_{k}\delta_{jq}+a_{i}^{\dagger}a_{j}^{\dagger}a_{l}a_{q}\delta_{pk}\right)
\right\}  D^{2N}\right\}  .
\end{align}
This produces%
\begin{align}
\frac{\partial\mathcal{E}\left(  \mathbb{D}^{2},\mathbf{0}\right)  }%
{\partial\alpha_{pq}}=-  &  i\operatorname{Tr}\left\{  \sum_{1\leq j\leq
2s}\left\{  \left\{  h_{1jp}a_{j}^{\dagger}a_{q}-h_{1qj}a_{p}^{\dagger}%
a_{j}\right\}  \right.  \right. \nonumber\\
&  +\left.  \left.  \frac{1}{2}\sum_{1\leq j,k,l\leq2s}\left\{  V_{2ijkl}%
a_{l}^{\dagger}a_{j}^{\dagger}a_{q}a_{k}-V_{2ijkl}a_{p}^{\dagger}%
a_{j}^{\dagger}a_{l}a_{k}\right\}  \right\}  D^{2N}\right\}  ,
\end{align}
giving%
\begin{equation}
\frac{\partial\mathcal{E}\left(  \mathbb{D}^{2},\mathbf{0}\right)  }%
{\partial\alpha_{pq}}=-i\left\{  \sum_{1\leq j\leq2s}\left\{  D_{qj}%
^{1}h_{1jp}-h_{1qj}D_{jp}^{1}\right\}  +\frac{1}{2}\sum_{1\leq i,j,k\leq
2s}\left\{  D_{kqlj}^{2}V_{2ijkl}-V_{2ijkl}D_{klpj}^{2}\right\}  \right\}
\end{equation}
and finally
\begin{equation}
\frac{\partial\mathcal{E}\left(  \mathbb{D}^{2},\mathbf{0}\right)  }%
{\partial\alpha_{pq}}=-i\left(  \left[  \mathbb{D}^{1},\mathbb{H}_{1}\left(
\mathbf{\varphi}\left(  \mathbf{\alpha}\right)  \right)  \right]  +\frac{1}%
{2}L_{2}^{1}\left(  \left[  \mathbb{D}^{2},\mathbb{V}_{2}\left(
\mathbf{\varphi}\left(  \mathbf{\alpha}\right)  \right)  \right]  \right)
\right)  _{qp}. \label{derpq}%
\end{equation}
The derivatives $\left\{  \frac{\partial\mathcal{E}\left(  \mathbb{D}%
^{2},\mathbf{0}\right)  }{\partial\alpha_{pq}}\right\}  $ can also be
expressed in terms of the non antisymmetric matrix $\mathbb{H}_{2}\left(
\mathbf{\varphi}\left(  \mathbf{\alpha}\right)  \right)  $ as
\begin{equation}
\frac{\partial\mathcal{E}\left(  \mathbb{D}^{2},\mathbf{0}\right)  }%
{\partial\alpha_{pq}}=-i\left(  \left[  \mathbb{D}^{1},\mathbb{H}_{1}\left(
\mathbf{\varphi}\left(  \mathbf{\alpha}\right)  \right)  \right]  +L_{2}%
^{1}\left(  \left[  \mathbb{D}^{2},\mathbb{H}_{2}\left(  \mathbf{\varphi
}\left(  \mathbf{\alpha}\right)  \right)  \right]  \right)  \right)  _{qp}.
\label{derpq1}%
\end{equation}
If we compare eq. $\left(  \ref{derpq}\right)  $ with eqs. $\left(
\ref{Fock_Op}\right)  $ and $\left(  \ref{Fock_Op_adj}\right)  $ one observes
the important result that
\begin{equation}
\frac{\partial\mathcal{E}\left(  \mathbb{D}^{2},\mathbf{0}\right)  }%
{\partial\alpha_{pq}}=i\left(  F\left(  \mathbf{\varphi}\left(  \mathbf{\alpha
}\right)  \right)  -F\left(  \mathbf{\varphi}\left(  \mathbf{\alpha}\right)
\right)  ^{\dagger}\right)  _{qp},
\end{equation}
where
\begin{equation}
\left\vert \varphi_{j}\left(  \mathbf{\alpha}\right)  \right\rangle =U\left(
\mathbf{\alpha}\right)  \left\vert \varphi_{j}\right\rangle ;\quad1\leq
j\leq2s.
\end{equation}


\section{Stationarity Conditions\label{StatCon}}

At a stationary point wrt unitary group parameters $\mathbf{\alpha}$ the
derivatives eq. $\left(  \ref{derpq}\right)  $ must equal zero hence
\begin{equation}
\left(  F\left(  \mathbf{\varphi}\left(  \mathbf{\alpha}\right)  \right)
-F\left(  \mathbf{\varphi}\left(  \mathbf{\alpha}\right)  \right)  ^{\dagger
}\right)  _{qp}=0\quad1\leq p,q\leq2s
\end{equation}
and the generalized Fock operator is self adjoint i.e.
\begin{equation}
F\left(  \mathbf{\varphi}\left(  \mathbf{\alpha}\right)  \right)  =F\left(
\mathbf{\varphi}\left(  \mathbf{\alpha}\right)  \right)  ^{\dagger}.
\label{adjcon}%
\end{equation}
In order to constrain the GSO's to be orthonormal in a direct non unitary
group variational procedure one considers the Lagrangian
\begin{equation}
\mathcal{L}\left(  \mathbb{D}^{2},\mathbf{\varphi}\right)  =\mathcal{E}\left(
\mathbb{D}^{2},\mathbf{\varphi}\right)  -\sum_{1\leq j,k\leq2s}\varepsilon
_{kj}\left\{  \left\langle \left.  \varphi_{j}\right\vert \varphi
_{k}\right\rangle -\delta_{jk}\right\}  ,
\end{equation}
where $\left\{  \left.  \varepsilon_{jk}\right\vert 1\leq j,k\leq2s\right\}  $
are Lagrange multipliers, which produces the constrained stationarity
condition
\begin{equation}
\frac{\delta\mathcal{L}\left(  \mathbb{D}^{2},\mathbf{\varphi}\right)
}{\delta\varphi_{j}^{\ast}}=\frac{\delta\mathcal{E}\left(  \mathbb{D}%
^{2},\mathbf{\varphi}\right)  }{\delta\varphi_{j}^{\ast}}-\sum_{1\leq k\leq
2s}\left\vert \varphi_{k}\right\rangle \varepsilon_{kj}=0;\quad1\leq j\leq2s.
\label{Lagran}%
\end{equation}
Using eq. $\left(  \ref{FuncDerFock}\right)  $ eq. $\left(  \ref{Lagran}%
\right)  $ can be expressed as
\begin{equation}
F\left(  \mathbb{D}^{2},\mathbf{\varphi}\right)  \left\vert \varphi
_{j}\right\rangle =\sum_{1\leq k\leq2s}\left\vert \varphi_{k}\right\rangle
\varepsilon_{kj};\quad1\leq j\leq2s,
\end{equation}
which in combination with eq. $\left(  \ref{adjcon}\right)  $ shows that the
matrix of Lagrange multipliers is hermitian when $\mathbf{\varphi}$ is a
constrained stationary point.

\section{Derivatives of the AGP\ Energy Functional\label{variational}}

The $2N$-AGP\ energy functional in terms of the NGSO's $\left\{  \left.
\varphi_{j}\right\vert 1\leq j\leq2s\right\}  $, the geminal occupation
numbers $\left\{  \left.  n_{j}\right\vert 1\leq j\leq s\right\}  $ and the
phases $\left\{  \left.  \theta_{j}\right\vert 1\leq j\leq s\right\}  $ can be
expressed as the quotient%
\begin{align}
&  \mathcal{E}\left(  \mathbf{n},\mathbf{\varphi},\mathbf{\theta}\right)
=\frac{\mathcal{H}\left(  \mathbf{n},\mathbf{\varphi},\mathbf{\theta}\right)
}{S_{N}\left(  \mathbf{n}\right)  }\nonumber\\
&  =\frac{1}{S_{N}\left(  \mathbf{n}\right)  }\left\{  \sum_{1\leq j\leq
s}n_{j}\left\{  \left\langle \varphi_{j}\left\vert h_{1}\varphi_{j}\right.
\right\rangle +\left\langle \varphi_{j+s}\left\vert h_{1}\varphi_{j+s}\right.
\right\rangle +\left\langle \varphi_{j}\varphi_{j+s}\left\Vert \varphi
_{j}\varphi_{j+s}\right.  \right\rangle \right\}  \right. \nonumber\\
&  +\sum_{1\leq j,k\leq s}n_{j}^{\frac{1}{2}}n_{k}^{\frac{1}{2}}e^{i\left(
\theta_{j}-\theta_{k}\right)  }S_{N-1}\left(  \jmath k\right)  \varphi
_{j}\varphi_{j+s}|\left\vert \varphi_{k}\varphi_{k+s}\right\rangle \left.
+\frac{1}{2}\sum_{1\leq j,k\leq s}n_{j}n_{k}S_{N-2}\left(  \jmath k\right)
\left\langle \varphi_{j}\varphi_{k}|||\varphi_{j}\varphi_{k}\right\rangle
\right\}  .
\end{align}
The first order variation $\delta^{1}\mathcal{E}\left(  \mathbf{n}%
,\mathbf{\varphi},\mathbf{\theta}\right)  $ of the energy is
\begin{equation}
\delta^{1}\mathcal{E}\left(  \mathbf{n},\mathbf{\varphi},\mathbf{\theta
}\right)  =\left(
\begin{array}
[c]{cccc}%
\partial_{\mathbf{n}}^{1}\mathcal{E}\left(  \mathbf{n},\mathbf{\varphi
},\mathbf{\theta}\right)  & \partial_{\mathbf{\varphi}^{\ast}}^{1}%
\mathcal{E}\left(  \mathbf{n},\mathbf{\varphi},\mathbf{\theta}\right)  &
\partial_{\mathbf{\varphi}}^{1}\mathcal{E}\left(  \mathbf{n},\mathbf{\varphi
},\mathbf{\theta}\right)  & \partial_{\mathbf{\theta}}^{1}\mathcal{E}\left(
\mathbf{n},\mathbf{\varphi},\mathbf{\theta}\right)
\end{array}
\right)  \left(
\begin{array}
[c]{c}%
\delta\mathbf{n}\\
\delta\mathbf{\varphi}^{\ast}\\
\delta\mathbf{\varphi}\\
\delta\mathbf{\theta}%
\end{array}
\right)  ,
\end{equation}
where
\begin{subequations}
\label{Fderiv}%
\begin{align}
\partial_{\mathbf{n}}^{1}\mathcal{E}\left(  \mathbf{n},\mathbf{\varphi
},\mathbf{\theta}\right)   &  =\frac{1}{S_{N}\left(  \mathbf{n}\right)
}\left\{  \partial_{\mathbf{n}}^{1}\mathcal{H}\left(  \mathbf{n}%
,\mathbf{\varphi},\mathbf{\theta}\right)  -\mathcal{E}\left(  \mathbf{n}%
,\mathbf{\varphi},\mathbf{\theta}\right)  \partial_{\mathbf{n}}^{1}%
S_{N}\left(  \mathbf{n}\right)  \right\} \\
\partial_{\mathbf{\psi}^{\ast}}^{1}\mathcal{E}\left(  \mathbf{n}%
,\mathbf{\varphi},\mathbf{\theta}\right)   &  =\frac{1}{S_{N}\left(
\mathbf{n}\right)  }\left\{  \partial_{\mathbf{\psi}^{\ast}}^{1}%
\mathcal{H}\left(  \mathbf{n},\mathbf{\varphi},\mathbf{\theta}\right)
-\mathcal{E}\left(  \mathbf{n},\mathbf{\varphi},\mathbf{\theta}\right)
\partial_{\mathbf{\psi}^{\ast}}^{1}S_{N}\left(  \mathbf{n}\right)  \right\}
=\left[  \partial_{\mathbf{\psi}}^{1}\mathcal{E}\left(  \mathbf{n}%
,\mathbf{\varphi},\mathbf{\theta}\right)  \right]  ^{\ast}\\
\partial_{\mathbf{\theta}}^{1}\mathcal{E}\left(  \mathbf{n},\mathbf{\varphi
},\mathbf{\theta}\right)   &  =\frac{1}{S_{N}\left(  \mathbf{n}\right)
}\left\{  \partial_{\mathbf{\theta}}^{1}\mathcal{H}\left(  \mathbf{n}%
,\mathbf{\varphi},\mathbf{\theta}\right)  -\mathcal{E}\left(  \mathbf{n}%
,\mathbf{\varphi},\mathbf{\theta}\right)  \partial_{\mathbf{\theta}}^{1}%
S_{N}\left(  \mathbf{n}\right)  \right\}
\end{align}
and
\end{subequations}
\begin{subequations}
\begin{align}
\partial_{\mathbf{\varphi}^{\ast}}^{1}\mathcal{E}\left(  \mathbf{n}%
,\mathbf{\varphi},\mathbf{\theta}\right)  \cdot\delta\mathbf{\varphi}^{\ast}
&  =\sum_{1\leq j\leq2s}\partial_{\varphi_{j}^{\ast}}^{1}\mathcal{E}\left(
\mathbf{n},\mathbf{\varphi},\mathbf{\theta}\right)  \left(  \delta\varphi
_{j}^{\ast}\right)  =\sum_{1\leq j\leq2s}\int\frac{\delta\mathcal{E}\left(
\mathbf{n},\mathbf{\varphi},\mathbf{\theta}\right)  }{\delta\varphi_{j}^{\ast
}\left(  \mathbf{r,}\zeta\right)  }\delta\varphi_{j}^{\ast}\left(
\mathbf{r,}\zeta\right)  d\mathbf{r}d\zeta\label{funcder}\\
\partial_{\mathbf{n}}^{1}\mathcal{E}\left(  \mathbf{n},\mathbf{\varphi
},\mathbf{\theta}\right)  \cdot\delta\mathbf{n}  &  =\sum_{1\leq j\leq s}%
\frac{\partial\mathcal{E}\left(  \mathbf{n},\mathbf{\varphi},\mathbf{\theta
}\right)  }{\partial n_{j}}\delta n_{j}\\
\partial_{\mathbf{\theta}}^{1}\mathcal{E}\left(  \mathbf{n},\mathbf{\varphi
},\mathbf{\theta}\right)  \cdot\delta\mathbf{\theta}  &  =\sum_{1\leq j\leq
s}\frac{\partial\mathcal{E}\left(  \mathbf{n},\mathbf{\varphi},\mathbf{\theta
}\right)  }{\partial\theta_{j}}\delta\theta_{j}.
\end{align}
If the variations $\left\{  \delta\mathbf{\varphi}^{\ast},\delta
\mathbf{\varphi}\right\}  $ maintain the orthonormality of the NGSO's i.e.
\end{subequations}
\begin{equation}
\left\langle \left.  \varphi_{j}\right\vert \varphi_{k}\right\rangle
=\delta_{jk}%
\end{equation}
then
\begin{subequations}
\label{fderiv2}%
\begin{align}
\partial_{\mathbf{\psi}^{\ast}}^{1}\mathcal{E}\left(  \mathbf{n}%
,\mathbf{\varphi},\mathbf{\theta}\right)   &  =\frac{1}{S_{N}\left(
\mathbf{n}\right)  }\partial_{\mathbf{\psi}^{\ast}}^{1}\mathcal{H}\left(
\mathbf{n},\mathbf{\varphi},\mathbf{\theta}\right) \\
\partial_{\mathbf{\theta}}^{1}\mathcal{E}\left(  \mathbf{n},\mathbf{\varphi
},\mathbf{\theta}\right)   &  =\frac{1}{S_{N}\left(  \mathbf{n}\right)
}\partial_{\mathbf{\theta}}^{1}\mathcal{H}\left(  \mathbf{n},\mathbf{\varphi
},\mathbf{\theta}\right)
\end{align}
as $\partial_{\mathbf{\varphi}^{\ast}}^{1}S_{N}\left(  \mathbf{n}\right)
=\partial_{\mathbf{\varphi}}^{1}S_{N}\left(  \mathbf{n}\right)  =\partial
_{\mathbf{\theta}}^{1}S_{N}\left(  \mathbf{n}\right)  =\mathbf{0}$.

Variation in terms of the variables $\left\{  \mathbf{n},\mathbf{\varphi
},\mathbf{\theta}\right\}  $ leads to a constrained variational problem, which
can be avoided if one uses the group parameters $\left\{  \mathbf{\eta
},\mathbf{\lambda},\mathbf{\theta}\right\}  $ and the energy functional
\end{subequations}
\begin{equation}
\mathcal{E}\left(  \mathbf{\eta},\mathbf{\lambda},\mathbf{\theta}\right)
=\frac{\mathcal{H}\left(  \mathbf{\eta},\mathbf{\lambda},\mathbf{\theta
}\right)  }{S_{N}\left(  \mathbf{\eta}\right)  }=\frac{\left\langle \left.
g\left(  \mathbf{\eta},\mathbf{\lambda},\mathbf{\theta}\right)  ^{N}%
\right\vert Hg\left(  \mathbf{\eta},\mathbf{\lambda},\mathbf{\theta}\right)
^{N}\right\rangle }{S_{N}\left(  \mathbf{\eta}\right)  }%
\end{equation}
and considers the derivatives
\begin{subequations}
\label{GFderiv}%
\begin{align}
\partial_{\mathbf{\eta}}^{1}\mathcal{E}\left(  \mathbf{\eta},\mathbf{\lambda
},\mathbf{\theta}\right)   &  =\frac{1}{S_{N}\left(  \mathbf{\eta}\right)
}\left\{  \partial_{\mathbf{\eta}}^{1}\mathcal{H}\left(  \mathbf{\eta
},\mathbf{\lambda},\mathbf{\theta}\right)  -\mathcal{E}\left(  \mathbf{\eta
},\mathbf{\lambda},\mathbf{\theta}\right)  \partial_{\mathbf{\eta}}^{1}%
S_{N}\left(  \mathbf{\eta}\right)  \right\} \\
\partial_{\mathbf{\lambda}}^{1}\mathcal{E}\left(  \mathbf{\eta}%
,\mathbf{\lambda},\mathbf{\theta}\right)   &  =\frac{1}{S_{N}\left(
\mathbf{\eta}\right)  }\partial_{\mathbf{\lambda}}^{1}\mathcal{H}\left(
\mathbf{\eta},\mathbf{\lambda},\mathbf{\theta}\right) \\
\partial_{\mathbf{\theta}}^{1}\mathcal{E}\left(  \mathbf{\eta},\mathbf{\lambda
},\mathbf{\theta}\right)   &  =\frac{1}{S_{N}\left(  \mathbf{\eta}\right)
}\partial_{\mathbf{\theta}}^{1}\mathcal{H}\left(  \mathbf{\eta}%
,\mathbf{\lambda},\mathbf{\theta}\right)  .
\end{align}
The generators of the GSO and the phase variations produce
\end{subequations}
\begin{equation}
\partial_{\mathbf{\lambda}_{jk}}^{1}\mathcal{E}\left(  \mathbf{\eta
},\mathbf{\lambda},\mathbf{\theta}\right)  =\frac{i}{S_{N}\left(
\mathbf{\eta}\right)  }\left\langle \left.  g\left(  \mathbf{\eta
},\mathbf{\lambda},\mathbf{\theta}\right)  ^{N}\right\vert \left[
H,a_{j}^{\dagger}a_{k}\right]  g\left(  \mathbf{\eta},\mathbf{\lambda
},\mathbf{\theta}\right)  ^{N}\right\rangle ;\quad1\leq j\neq k\pm s\leq2s
\label{Dlamda}%
\end{equation}
and
\begin{equation}
\partial_{\mathbf{\theta}_{j}}^{1}\mathcal{E}\left(  \mathbf{\eta
},\mathbf{\lambda},\mathbf{\theta}\right)  =\frac{i}{S_{N}\left(
\mathbf{\eta}\right)  }\left\langle \left.  g\left(  \mathbf{\eta
},\mathbf{\lambda},\mathbf{\theta}\right)  ^{N}\right\vert \left[  H,\left(
a_{j}^{\dagger}a_{j}+a_{j+s}^{\dagger}a_{j+s}\right)  \right]  g\left(
\mathbf{\eta},\mathbf{\lambda},\mathbf{\theta}\right)  ^{N}\right\rangle
;\quad1\leq j\leq s, \label{Dtheta}%
\end{equation}
while the generators of the occupation number variations give the derivatives
\begin{align}
\partial_{\mathbf{\eta}_{j}}^{1}\mathcal{E}\left(  \mathbf{\eta}%
,\mathbf{\lambda},\mathbf{\theta}\right)   &  =\frac{1}{S_{N}\left(
\mathbf{\eta}\right)  }\left\{  \left\langle \left.  g\left(  \mathbf{\eta
},\mathbf{\lambda},\mathbf{\theta}\right)  ^{N}\right\vert \left[  H,\left(
a_{j}^{\dagger}a_{j}+a_{j+s}^{\dagger}a_{j+s}\right)  \right]  _{+}g\left(
\mathbf{\eta},\mathbf{\lambda},\mathbf{\theta}\right)  ^{N}\right\rangle
\right. \nonumber\\
&  -\left.  2\mathcal{E}\left(  \mathbf{\eta},\mathbf{\lambda},\mathbf{\theta
}\right)  \left\langle \left.  g\left(  \mathbf{\eta},\mathbf{\lambda
},\mathbf{\theta}\right)  ^{N}\right\vert \left(  a_{j}^{\dagger}a_{j}%
+a_{j+s}^{\dagger}a_{j+s}\right)  g\left(  \mathbf{\eta},\mathbf{\lambda
},\mathbf{\theta}\right)  ^{N}\right\rangle \right\}  ;\quad1\leq j\leq s.
\label{Deta}%
\end{align}


\section{Stationarity Conditions for an Optimized AGP\label{StatAGPCond}}

At a stationary point of the unconstrained variation wrt $\left(
\mathbf{\eta},\mathbf{\lambda},\mathbf{\theta}\right)  $
\begin{equation}
\partial_{\mathbf{\theta}_{j}}^{1}\mathcal{E}\left(  \mathbf{\eta
},\mathbf{\lambda},\mathbf{\theta}\right)  =\frac{i}{S_{N}\left(
\mathbf{\eta}\right)  }\left\langle \left.  g\left(  \mathbf{\eta
},\mathbf{\lambda},\mathbf{\theta}\right)  ^{N}\right\vert \left[  H,\left(
a_{j}^{\dagger}a_{j}+a_{j+s}^{\dagger}a_{j+s}\right)  \right]  g\left(
\mathbf{\eta},\mathbf{\lambda},\mathbf{\theta}\right)  ^{N}\right\rangle
=0;\quad1\leq j\leq s, \label{s1}%
\end{equation}
and%
\begin{equation}
\partial_{\mathbf{\lambda}_{jk}}^{1}\mathcal{E}\left(  \mathbf{\eta
},\mathbf{\lambda},\mathbf{\theta}\right)  =\frac{i}{S_{N}\left(
\mathbf{\eta}\right)  }\left\langle \left.  g\left(  \mathbf{\eta
},\mathbf{\lambda},\mathbf{\theta}\right)  ^{N}\right\vert \left[
H,a_{j}^{\dagger}a_{k}\right]  g\left(  \mathbf{\eta},\mathbf{\lambda
},\mathbf{\theta}\right)  ^{N}\right\rangle =0;\quad1\leq j\neq k\pm s\leq2s.
\label{s2}%
\end{equation}
Eqs. $\left(  \ref{s1}\right)  $ and $\left(  \ref{s2}\right)  $ combined with
the commutator relationships deduced from the invariance group (see eq.
$\left(  \ref{iso}\right)  $) i.e.
\begin{align}
\left\langle \left.  g\left(  \mathbf{\eta},\mathbf{\lambda},\mathbf{\theta
}\right)  ^{N}\right\vert \left[  H,\left(  a_{j}^{\dagger}a_{j}%
-a_{j+s}^{\dagger}a_{j+s}\right)  \right]  g\left(  \mathbf{\eta
},\mathbf{\lambda},\mathbf{\theta}\right)  ^{N}\right\rangle  &  =\nonumber\\
\left\langle \left.  g\left(  \mathbf{\eta},\mathbf{\lambda},\mathbf{\theta
}\right)  ^{N}\right\vert \left[  H,a_{j}^{\dagger}a_{j+s}\right]  g\left(
\mathbf{\eta},\mathbf{\lambda},\mathbf{\theta}\right)  ^{N}\right\rangle  &
=\left\langle \left.  g\left(  \mathbf{\eta},\mathbf{\lambda},\mathbf{\theta
}\right)  ^{N}\right\vert \left[  H,a_{j+s}^{\dagger}a_{j}\right]  g\left(
\mathbf{\eta},\mathbf{\lambda},\mathbf{\theta}\right)  ^{N}\right\rangle =0
\end{align}
give
\begin{equation}
\left\langle \left.  g\left(  \mathbf{\eta},\mathbf{\lambda},\mathbf{\theta
}\right)  ^{N}\right\vert \left[  H,a_{j}^{\dagger}a_{k}\right]  g\left(
\mathbf{\eta},\mathbf{\lambda},\mathbf{\theta}\right)  ^{N}\right\rangle
=0;\quad1\leq j,k\leq2s. \label{s3}%
\end{equation}
The stationarity of the occupation number variations gives
\begin{align}
\partial_{\mathbf{\eta}_{j}}^{1}\mathcal{E}\left(  \mathbf{\eta}%
,\mathbf{\lambda},\mathbf{\theta}\right)   &  =\frac{1}{S_{N}\left(
\mathbf{\eta}\right)  }\left\{  \left\langle \left.  g\left(  \mathbf{\eta
},\mathbf{\lambda},\mathbf{\theta}\right)  ^{N}\right\vert \left[  H,\left(
a_{j}^{\dagger}a_{j}+a_{j+s}^{\dagger}a_{j+s}\right)  \right]  _{+}g\left(
\mathbf{\eta},\mathbf{\lambda},\mathbf{\theta}\right)  ^{N}\right\rangle
\right. \nonumber\\
&  -\left.  2\mathcal{E}\left(  \mathbf{\eta},\mathbf{\lambda},\mathbf{\theta
}\right)  \left\langle \left.  g\left(  \mathbf{\eta},\mathbf{\lambda
},\mathbf{\theta}\right)  ^{N}\right\vert \left(  a_{j}^{\dagger}a_{j}%
+a_{j+s}^{\dagger}a_{j+s}\right)  g\left(  \mathbf{\eta},\mathbf{\lambda
},\mathbf{\theta}\right)  ^{N}\right\rangle \right\}  =0;\quad1\leq j\leq s,
\end{align}
which combined with eq. $\left(  \ref{s3}\right)  $ and noting that
\begin{equation}
\left(  a_{j}^{\dagger}a_{j}+a_{j+s}^{\dagger}a_{j+s}\right)  \left\vert
g\left(  \mathbf{\eta},\mathbf{\lambda},\mathbf{\theta}\right)  ^{N}%
\right\rangle =2a_{j}^{\dagger}a_{j}\left\vert g\left(  \mathbf{\eta
},\mathbf{\lambda},\mathbf{\theta}\right)  ^{N}\right\rangle
\end{equation}
leads to
\begin{equation}
\left\langle \left.  g\left(  \mathbf{\eta},\mathbf{\lambda},\mathbf{\theta
}\right)  ^{N}\right\vert Ha_{j}^{\dagger}a_{j}g\left(  \mathbf{\eta
},\mathbf{\lambda},\mathbf{\theta}\right)  ^{N}\right\rangle =\mathcal{E}%
\left(  \mathbf{\eta},\mathbf{\lambda},\mathbf{\theta}\right)  \left\langle
\left.  g\left(  \mathbf{\eta},\mathbf{\lambda},\mathbf{\theta}\right)
^{N}\right\vert a_{j}^{\dagger}a_{j}g\left(  \mathbf{\eta},\mathbf{\lambda
},\mathbf{\theta}\right)  ^{N}\right\rangle ;\quad1\leq j\leq2s,
\end{equation}
which can be written as
\begin{equation}
\left\langle \left.  g\left(  \mathbf{\eta},\mathbf{\lambda},\mathbf{\theta
}\right)  ^{N}\right\vert Ha_{j}^{\dagger}a_{j}g\left(  \mathbf{\eta
},\mathbf{\lambda},\mathbf{\theta}\right)  ^{N}\right\rangle =\mathcal{E}%
\left(  \mathbf{\eta},\mathbf{\lambda},\mathbf{\theta}\right)  N_{j}.
\end{equation}


\section{Stationary AGP Fock Operator\label{AGPFock}}

The matrix elements $\left\{  \left.  F\left(  g\left(  \mathbf{\eta
},\mathbf{\psi}\right)  ^{N}\right)  _{qp}\right\vert 1\leq p,q\leq2s\right\}
$ of the generalized Fock operator are
\begin{align}
F\left(  g\left(  \mathbf{\eta},\mathbf{\psi}\right)  ^{N}\right)  _{qp} &
=\frac{1}{S_{N}\left(  \mathbf{\eta}\right)  }\left\langle \left.  g\left(
\mathbf{\eta},\mathbf{\psi}\right)  ^{N}\right\vert a_{p}^{\dagger}\left[
a_{q},H\right]  g\left(  \mathbf{\eta},\mathbf{\psi}\right)  ^{N}\right\rangle
\nonumber\\
&  =\frac{1}{S_{N}\left(  \mathbf{\eta}\right)  }\left\langle \left.  g\left(
\mathbf{\eta},\mathbf{\psi}\right)  ^{N}\right\vert \left\{  a_{p}^{\dagger
}a_{q}H-a_{p}^{\dagger}Ha_{q}\right\}  g\left(  \mathbf{\eta},\mathbf{\psi
}\right)  ^{N}\right\rangle .
\end{align}
The diagonal elements of $F\left(  g\left(  \mathbf{\eta},\mathbf{\psi
}\right)  ^{N}\right)  $ for an optimized AGP\ have the form
\begin{align}
F\left(  g\left(  \mathbf{\eta},\mathbf{\psi}\right)  ^{N}\right)  _{qq} &
=\frac{N_{q}}{S_{N}\left(  \mathbf{\eta}\right)  }\left\{  \mathcal{E}\left(
g\left(  \mathbf{\eta},\mathbf{\psi}\right)  ^{N}\right)  -\frac{\left\langle
\left.  g\left(  \mathbf{\eta},\mathbf{\psi}\right)  ^{N}\right\vert
a_{q}^{\dagger}Ha_{q}g\left(  \mathbf{\eta},\mathbf{\psi}\right)
^{N}\right\rangle }{\left\langle \left.  g\left(  \mathbf{\eta},\mathbf{\psi
}\right)  ^{N}\right\vert a_{q}^{\dagger}a_{q}g\left(  \mathbf{\eta
},\mathbf{\psi}\right)  ^{N}\right\rangle }\right\}  \nonumber\\
&  =\frac{N_{q}}{S_{N}\left(  \mathbf{\eta}\right)  }\left\{  \mathcal{E}%
\left(  g\left(  \mathbf{\eta},\mathbf{\psi}\right)  ^{N}\right)
-\mathcal{E}_{q}\left(  g\left(  \mathbf{\eta},\mathbf{\psi}\right)
^{N}\right)  \right\}  ,
\end{align}
where $\mathcal{E}_{q}\left(  g\left(  \mathbf{\eta},\mathbf{\psi}\right)
^{N}\right)  $ is the energy of the normalized $2N-1$ electron state formed by
removing an electron described by a CGSO $\left\vert \psi_{q}\right\rangle $
\begin{equation}
\mathcal{E}_{q}\left(  g\left(  \mathbf{\eta},\mathbf{\psi}\right)
^{N}\right)  =\frac{\left\langle \left.  g\left(  \mathbf{\eta},\mathbf{\psi
}\right)  ^{N}\right\vert a_{q}^{\dagger}Ha_{q}g\left(  \mathbf{\eta
},\mathbf{\psi}\right)  ^{N}\right\rangle }{\left\langle \left.  a_{q}g\left(
\mathbf{\eta},\mathbf{\psi}\right)  ^{N}\right\vert a_{q}g\left(
\mathbf{\eta},\mathbf{\psi}\right)  ^{N}\right\rangle }.
\end{equation}
It is easy to observe that
\begin{equation}
\mathcal{E}_{q+s}\left(  g\left(  \mathbf{\eta},\mathbf{\psi}\right)
^{N}\right)  =\mathcal{E}_{q}\left(  g\left(  \mathbf{\eta},\mathbf{\psi
}\right)  ^{N}\right)
\end{equation}
and
\begin{equation}
F\left(  g\left(  \mathbf{\eta},\mathbf{\psi}\right)  ^{N}\right)
_{qq}=F\left(  g\left(  \mathbf{\eta},\mathbf{\psi}\right)  ^{N}\right)
_{q+sq+s};\quad1\leq q\leq s.
\end{equation}
The off diagonal elements are given by
\begin{equation}
F\left(  g\left(  \mathbf{\eta},\mathbf{\psi}\right)  ^{N}\right)  _{qp}%
=\frac{1}{S_{N}\left(  \mathbf{\eta}\right)  }\left\langle \left.  g\left(
\mathbf{\eta},\mathbf{\psi}\right)  ^{N}\right\vert a_{p}^{\dagger}%
Ha_{q}g\left(  \mathbf{\eta},\mathbf{\psi}\right)  ^{N}\right\rangle ,
\end{equation}
which are \emph{not }in general zero in the CGSO basis.  

\section{AGP Excitation Operators\label{q_op}}

One can define operators $\left\{  \left.  q_{kj}^{\dagger}\right\vert 1\leq
j\neq k\pm s\leq2s\right\}  $ by
\begin{equation}
q_{jk}^{\dagger}\left(  \mathbf{\eta,\psi}\right)  =c_{k}\left(  \mathbf{\eta
}\right)  a_{j}^{\dagger}a_{k}-\sigma\left(  j\right)  \sigma\left(  k\right)
c_{j}\left(  \mathbf{\eta}\right)  a_{\left\vert k\right\vert }^{\dagger
}a_{\left\vert j\right\vert }\text{ for }n_{j}\neq n_{k},
\end{equation}
where%
\begin{equation}
\left\vert j\right\vert =\left\{
\begin{array}
[c]{c}%
j;\qquad1\leq j\leq s\\
j-s;\qquad s+1\leq j\leq2s
\end{array}
\right.
\end{equation}
and
\begin{equation}
\sigma\left(  j\right)  =\left\{
\begin{array}
[c]{c}%
-1;\qquad1\leq j\leq s\quad\quad\\
\quad1;\qquad s+1\leq j\leq2s
\end{array}
\right.  .
\end{equation}
the operators $\left\{  a_{j}^{\dagger}\right\}  $ refer to basis formed by
the optimal CGSO's $\left\{  \left\vert \psi_{j}\right\rangle \right\}  $. The
importance of these operators lies in the observation that their adjoints
$\left\{  \left.  q_{kj}\right\vert 1\leq j\neq k\pm s\leq2s\right\}  $ have
the special annihilation property%
\begin{equation}
q_{jk}\left\vert g\left(  \mathbf{\eta},\mathbf{\varphi}\right)
^{N}\right\rangle =0;\quad1\leq j\neq k\pm s\leq2s.
\end{equation}
One can show \cite{Weiner&Goscinski83} that
\begin{equation}
\left(
\begin{array}
[c]{c}%
q_{jk}^{\dagger}\\
q_{j+sk}^{\dagger}\\
q_{jk+s}^{\dagger}\\
q_{j+sk+s}^{\dagger}\\
q_{jk}\\
q_{j+sk}\\
q_{jk+s}\\
q_{j+sk+s}%
\end{array}
\right)  =\left(
\begin{array}
[c]{cccccccc}%
c_{k} & 0 & \cdots & \cdots & \cdots & \cdots & 0 & -c_{j}\\
0 & c_{k} & 0 & \cdots & \cdots & 0 & c_{j} & 0\\
\vdots & 0 & c_{k} & 0 & 0 & c_{j} & 0 & \vdots\\
\vdots & \vdots & 0 & c_{k} & -c_{j} & 0 & \vdots & \vdots\\
\vdots & \vdots & 0 & -c_{j} & c_{k} & 0 & \vdots & \vdots\\
\vdots & 0 & c_{j} & 0 & 0 & c_{k} & 0 & \vdots\\
0 & c_{j} & 0 & \cdots & \cdots & 0 & c_{k} & 0\\
-c_{j} & 0 & \cdots & \cdots & \cdots & \cdots & 0 & c_{k}%
\end{array}
\right)  \left(
\begin{array}
[c]{c}%
a_{j}^{\dagger}a_{k}\\
a_{j}^{\dagger}a_{k+s}\\
a_{j+s}^{\dagger}a_{k}\\
a_{j+s}^{\dagger}a_{k+s}\\
a_{k}^{\dagger}a_{j}\\
a_{k}^{\dagger}a_{j+s}\\
a_{k+s}^{\dagger}a_{j}\\
a_{k+s}^{\dagger}a_{j+s}%
\end{array}
\right)  ,
\end{equation}
leading to the transpose relationship
\begin{align}
&  \left(  q_{jk}^{\dagger},q_{j+sk}^{\dagger},q_{jk+s}^{\dagger}%
,q_{j+sk+s}^{\dagger},q_{jk},q_{j+sk},q_{jk+s},q_{j+sk+s}\right) \nonumber\\
&  \qquad\qquad\qquad\qquad\qquad\left.  =\right.  \left(  a_{j}^{\dagger
}a_{k},a_{j}^{\dagger}a_{k+s},a_{j+s}^{\dagger}a_{k},a_{j+s}^{\dagger}%
a_{k+s},a_{k}^{\dagger}a_{j},a_{k}^{\dagger}a_{j+s},a_{k+s}^{\dagger}%
a_{j},a_{k+s}^{\dagger}a_{j+s}\right)  \mathbb{S}\left(  jk\right)
\end{align}
and that the transformation matrix $\mathbb{S}\left(  jk\right)  $ is non singular.

\section{Fock Operator and the Total Energy\label{Fock&Energy}}

\subsection{Homogeneous Functions\label{HomFunc}}

Polynomial functions can be expanded as a sum of homogeneous mononomial
functions as%

\begin{equation}
f\left(  \mathbf{x}\right)  =f_{0}+f_{1}\left(  \mathbf{x}\right)
+\cdots+f_{n}\left(  \mathbf{x}\right)  =%
%TCIMACRO{\tsum \limits_{0\leq m\leq n}}%
%BeginExpansion
{\textstyle\sum\limits_{0\leq m\leq n}}
%EndExpansion
f_{m}\left(  \mathbf{x}\right)  ,
\end{equation}
where
\begin{equation}
f_{m}\left(  \mathbf{x}\right)  =%
%TCIMACRO{\tsum \limits_{1\leq j_{1},\cdots,j_{m}\leq r}}%
%BeginExpansion
{\textstyle\sum\limits_{1\leq j_{1},\cdots,j_{m}\leq r}}
%EndExpansion
c_{j_{1}\cdots j_{m}}x_{j_{1}}\cdots x_{j_{m}}.
\end{equation}
The variables $\mathbf{x}$ can be complex and in particular $\mathbf{x=}%
\left(  \mathbf{z}^{\ast}\mathbf{,z}\right)  $, where we consider
$\mathbf{z}^{\ast}\mathbf{,z}$ as independent variables. Consider
\begin{equation}
\frac{\partial f_{m}}{\partial x_{l}}\left(  \mathbf{x}\right)  =%
%TCIMACRO{\tsum \limits_{1\leq j_{1},\cdots,j_{m-1}\neq l\leq r}}%
%BeginExpansion
{\textstyle\sum\limits_{1\leq j_{1},\cdots,j_{m-1}\neq l\leq r}}
%EndExpansion
c_{j_{1}\cdots j_{l}\cdots j_{m}}x_{j_{1}}\cdots x_{j_{m-1}}%
\end{equation}
then
\begin{equation}
x_{l}\frac{\partial f_{m}}{\partial x_{l}}\left(  \mathbf{x}\right)  =%
%TCIMACRO{\tsum \limits_{1\leq j_{1},\cdots,j_{m-1}\neq l\leq r}}%
%BeginExpansion
{\textstyle\sum\limits_{1\leq j_{1},\cdots,j_{m-1}\neq l\leq r}}
%EndExpansion
c_{j_{1}\cdots j_{l}\cdots j_{m}}x_{j_{1}}\cdots x_{j_{m-1}}x_{l},
\end{equation}
which gives
\begin{equation}%
%TCIMACRO{\tsum \limits_{1\leq l\leq r}}%
%BeginExpansion
{\textstyle\sum\limits_{1\leq l\leq r}}
%EndExpansion
x_{l}\frac{\partial f_{m}}{\partial x_{l}}\left(  \mathbf{x}\right)
=mf_{m}\left(  \mathbf{x}\right)  .
\end{equation}
Hence%
\begin{align}%
%TCIMACRO{\tsum \limits_{1\leq l\leq r}}%
%BeginExpansion
{\textstyle\sum\limits_{1\leq l\leq r}}
%EndExpansion
x_{l}\frac{\partial f}{\partial x_{l}}\left(  \mathbf{x}\right)   &  =%
%TCIMACRO{\tsum \limits_{1\leq l\leq r}}%
%BeginExpansion
{\textstyle\sum\limits_{1\leq l\leq r}}
%EndExpansion
x_{l}\frac{\partial}{\partial x_{l}}\left\{
%TCIMACRO{\tsum \limits_{0\leq m\leq n}}%
%BeginExpansion
{\textstyle\sum\limits_{0\leq m\leq n}}
%EndExpansion
f_{m}\left(  \mathbf{x}\right)  \right\}  =%
%TCIMACRO{\tsum \limits_{\substack{1\leq l\leq r\\0\leq m\leq n}}}%
%BeginExpansion
{\textstyle\sum\limits_{\substack{1\leq l\leq r\\0\leq m\leq n}}}
%EndExpansion
x_{l}\frac{\partial f_{m}\left(  \mathbf{x}\right)  }{\partial x_{l}%
}\nonumber\\
&  =%
%TCIMACRO{\tsum \limits_{\substack{0\leq m\leq n\\1\leq l\leq r}}}%
%BeginExpansion
{\textstyle\sum\limits_{\substack{0\leq m\leq n\\1\leq l\leq r}}}
%EndExpansion
x_{l}\frac{\partial f_{m}\left(  \mathbf{x}\right)  }{\partial x_{l}}=%
%TCIMACRO{\tsum \limits_{0\leq m\leq n}}%
%BeginExpansion
{\textstyle\sum\limits_{0\leq m\leq n}}
%EndExpansion
mf_{m}\left(  \mathbf{x}\right)  =f_{1}\left(  \mathbf{x}\right)
+2f_{2}\left(  \mathbf{x}\right)  +\cdots+nf_{n}\left(  \mathbf{x}\right)  .
\label{homex}%
\end{align}


\subsection{Fock Operator and Energy\label{Fock&En}}

The energy of \emph{any} $2N$-electron state is given by
\begin{equation}
\mathcal{E}\left(  \mathbb{D}^{2},\mathbf{\varphi}\right)  =\operatorname{Tr}%
\left\{  \mathbb{H}_{1}\left(  \mathbf{\varphi}\right)  \mathbb{D}%
^{1}\right\}  +\frac{1}{4}\operatorname{Tr}\left\{  \mathbb{V}_{2}\left(
\mathbf{\varphi}\right)  \mathbb{D}^{2}\right\}  \equiv\frac{1}%
{\operatorname{Tr}\left\{  D^{2N}\right\}  }\left\{  \mathcal{H}_{1}\left(
\mathbb{D}^{1},\mathbf{\varphi}\right)  +\mathcal{H}_{2}\left(  \mathbb{D}%
^{2},\mathbf{\varphi}\right)  \right\}  ,
\end{equation}
where
\begin{equation}
\mathcal{H}_{1}\left(  \mathbb{D}^{1},\mathbf{\varphi}\right)
=\operatorname{Tr}\left\{  \mathbb{H}_{1}\left(  \mathbf{\varphi}\right)
\mathbb{D}^{1}\right\}
\end{equation}
and
\begin{equation}
\mathcal{H}_{2}\left(  \mathbb{D}^{2},\mathbf{\varphi}\right)  =\frac{1}%
{4}\operatorname{Tr}\left\{  \mathbb{V}_{2}\left(  \mathbf{\varphi}\right)
\mathbb{D}^{2}\right\}  .
\end{equation}
As
\begin{equation}
\frac{\delta\mathcal{E}\left(  \mathbb{D}^{2},\mathbf{\varphi}\right)
}{\delta\varphi_{j}^{\ast}}=F\left(  \mathbb{D}^{2},\mathbf{\varphi}\right)
\left\vert \varphi_{j}\right\rangle
\end{equation}
one obtains
\begin{equation}%
%TCIMACRO{\tsum \limits_{1\leq j\leq2s}}%
%BeginExpansion
{\textstyle\sum\limits_{1\leq j\leq2s}}
%EndExpansion
\int\varphi_{j}^{\ast}\left(  \mathbf{r,\xi}\right)  \frac{\delta
\mathcal{E}\left(  \mathbb{D}^{2},\mathbf{\varphi}\right)  }{\delta\varphi
_{j}^{\ast}\left(  \mathbf{r,\xi}\right)  }d\mathbf{r}d\mathbf{\xi=}%
%TCIMACRO{\tsum \limits_{1\leq j\leq2s}}%
%BeginExpansion
{\textstyle\sum\limits_{1\leq j\leq2s}}
%EndExpansion
\left\langle \left.  \varphi_{j}\right\vert F\left(  \mathbb{D}^{2}%
,\mathbf{\varphi}\right)  \varphi_{j}\right\rangle =\operatorname{Tr}\left\{
F\left(  \mathbb{D}^{2},\mathbf{\varphi}\right)  \right\}  .
\end{equation}
The normalization $\operatorname{Tr}\left\{  D^{2N}\right\}  $ does not depend
on $\left\{  \varphi_{j}\right\}  $ as they are constrained to be orthonormal
thus%
\begin{equation}
\frac{\delta\mathcal{E}\left(  \mathbb{D}^{2},\mathbf{\varphi}\right)
}{\delta\varphi_{j}^{\ast}\left(  \mathbf{r,\xi}\right)  }=\frac
{1}{\operatorname{Tr}\left\{  D^{2N}\right\}  }\frac{\delta\mathcal{H}\left(
\mathbb{D}^{2},\mathbf{\varphi}\right)  }{\delta\varphi_{j}^{\ast}\left(
\mathbf{r,\xi}\right)  }%
\end{equation}
giving%
\begin{align}
\operatorname{Tr}\left\{  F\left(  \mathbb{D}^{2},\mathbf{\varphi}\right)
\right\}   &  =\frac{1}{\operatorname{Tr}\left\{  D^{2N}\right\}  }%
%TCIMACRO{\tsum \limits_{1\leq j\leq2s}}%
%BeginExpansion
{\textstyle\sum\limits_{1\leq j\leq2s}}
%EndExpansion
\int\varphi_{j}^{\ast}\left(  \mathbf{r,\xi}\right)  \frac{\delta
\mathcal{H}\left(  \mathbb{D}^{2},\mathbf{\varphi}\right)  }{\delta\varphi
_{j}^{\ast}\left(  \mathbf{r,\xi}\right)  }d\mathbf{r}d\mathbf{\xi}\\
&  \mathbf{=}\frac{1}{\operatorname{Tr}\left\{  D^{2N}\right\}  }%
%TCIMACRO{\tsum \limits_{1\leq j\leq2s}}%
%BeginExpansion
{\textstyle\sum\limits_{1\leq j\leq2s}}
%EndExpansion
\int\varphi_{j}^{\ast}\left(  \mathbf{r,\xi}\right)  \left\{  \frac
{\delta\mathcal{H}_{1}\left(  \mathbb{D}^{2},\mathbf{\varphi}\right)  }%
{\delta\varphi_{j}^{\ast}\left(  \mathbf{r,\xi}\right)  }+\frac{\delta
\mathcal{H}_{2}\left(  \mathbb{D}^{2},\mathbf{\varphi}\right)  }{\delta
\varphi_{j}^{\ast}\left(  \mathbf{r,\xi}\right)  }\right\}  d\mathbf{r}%
d\mathbf{\xi,}%
\end{align}
which by eq. $\left(  \ref{homex}\right)  $ produces
\begin{equation}
\operatorname{Tr}\left\{  F\left(  \mathbb{D}^{2},\mathbf{\varphi}\right)
\right\}  =\frac{1}{\operatorname{Tr}\left\{  D^{2N}\right\}  }\left\{
\mathcal{H}_{1}\left(  \mathbb{D}^{2},\mathbf{\varphi}\right)  +2\mathcal{H}%
_{2}\left(  \mathbb{D}^{2},\mathbf{\varphi}\right)  \right\}  .
\end{equation}
Noting that
\begin{equation}
\mathcal{H}\left(  \mathbb{D}^{2},\mathbf{\varphi}\right)  =\mathcal{H}%
_{1}\left(  \mathbb{D}^{2},\mathbf{\varphi}\right)  +\mathcal{H}_{2}\left(
\mathbb{D}^{2},\mathbf{\varphi}\right)
\end{equation}
and
\begin{equation}
\mathcal{E}\left(  \mathbb{D}^{2},\mathbf{\varphi}\right)  =\frac
{\mathcal{H}\left(  \mathbb{D}^{2},\mathbf{\varphi}\right)  }%
{\operatorname{Tr}\left\{  D^{2N}\right\}  }%
\end{equation}
one obtains
\begin{equation}
\operatorname{Tr}\left\{  F\left(  \mathbb{D}^{2},\mathbf{\varphi}\right)
\right\}  =\frac{1}{\operatorname{Tr}\left\{  D^{2N}\right\}  }\left\{
\mathcal{H}\left(  \mathbb{D}^{2},\mathbf{\varphi}\right)  +\mathcal{H}%
_{2}\left(  \mathbb{D}^{2},\mathbf{\varphi}\right)  \right\}  =\frac
{1}{\operatorname{Tr}\left\{  D^{2N}\right\}  }\left\{  2\mathcal{H}\left(
\mathbb{D}^{2},\mathbf{\varphi}\right)  -\mathcal{H}_{1}\left(  \mathbb{D}%
^{2},\mathbf{\varphi}\right)  \right\}  .
\end{equation}
Hence the important relationships
\begin{subequations}
\label{FE}%
\begin{align}
\mathcal{E}\left(  \mathbb{D}^{2},\mathbf{\varphi}\right)   &
=\operatorname{Tr}\left\{  F\left(  \mathbb{D}^{2},\mathbf{\varphi}\right)
\right\}  -\mathcal{H}_{2}\left(  \mathbb{D}^{2},\mathbf{\varphi}\right)
\label{FEa}\\
\mathcal{E}\left(  \mathbb{D}^{2},\mathbf{\varphi}\right)   &  =\frac{1}%
{2}\left\{  \overset{}{\operatorname{Tr}\left\{  F\left(  \mathbb{D}%
^{2},\mathbf{\varphi}\right)  \right\}  +\mathcal{H}_{1}\left(  \mathbb{D}%
^{2},\mathbf{\varphi}\right)  }\right\}  . \label{FEb}%
\end{align}


Recalling eq. $\left(  \ref{Fockelem}\right)  $
\end{subequations}
\begin{equation}
F\left(  \mathbb{D}^{2},\mathbf{\varphi}\right)  _{qp}=\frac{1}%
{\operatorname{Tr}\left\{  D^{2N}\right\}  }\operatorname{Tr}\left\{
\overset{}{\left\{  a_{p}^{\dagger}a_{q}H-a_{q}Ha_{p}^{\dagger}\right\}
}D^{2N}\right\}
\end{equation}
one observes that for \emph{any }state $D^{2N}$
\begin{align}
\operatorname{Tr}\left\{  F\left(  \mathbb{D}^{2},\mathbf{\varphi}\right)
\right\}   &  =\sum_{1\leq p\leq2s}\frac{1}{\operatorname{Tr}\left\{
D^{2N}\right\}  }\operatorname{Tr}\left\{  \overset{}{\left\{  a_{p}^{\dagger
}a_{p}H-a_{p}Ha_{p}^{\dagger}\right\}  }D^{2N}\right\} \nonumber\\
&  =\frac{1}{\operatorname{Tr}\left\{  D^{2N}\right\}  }\operatorname{Tr}%
\left\{  \mathfrak{N}HD^{2N}\right\}  -\sum_{1\leq p\leq2s}N_{p}\left(
\mathbf{\varphi}\right)  \frac{\operatorname{Tr}\left\{  a_{p}Ha_{p}^{\dagger
}D^{2N}\right\}  }{\operatorname{Tr}\left\{  a_{p}^{\dagger}a_{p}%
D^{2N}\right\}  }, \label{TrF1}%
\end{align}
where
\begin{equation}
\mathfrak{N=}\sum_{1\leq p\leq2s}a_{p}^{\dagger}a_{p}%
\end{equation}
is the number operator and
\begin{equation}
N_{p}\left(  \mathbb{D}^{2},\mathbf{\varphi}\right)  =\frac{\operatorname{Tr}%
\left\{  a_{p}^{\dagger}a_{p}D^{2N}\right\}  }{\operatorname{Tr}\left\{
D^{2N}\right\}  }%
\end{equation}
is the diagonal element of the FORDM of $D^{2N}$ in basis $\left\{  \left.
\varphi_{j}\right\vert 1\leq j\leq2s\right\}  $ i.e. the occupation of the GSO
$\left\vert \varphi_{p}\right\rangle .$ Eq. $\left(  \ref{TrF1}\right)  $ can
be further developed to give
\begin{equation}
\operatorname{Tr}\left\{  F\left(  \mathbb{D}^{2},\mathbf{\varphi}\right)
\right\}  =2N\mathcal{E}\left(  \mathbb{D}^{2},\mathbf{\varphi}\right)
-\sum_{1\leq p\leq2s}N_{p}\left(  \mathbb{D}^{2},\mathbf{\varphi}\right)
\mathcal{E}_{p}\left(  D^{2N},\mathbf{\varphi}\right)  ,
\end{equation}
where the energy of the $2N$ electron state $D^{2N}$ is
\begin{equation}
\mathcal{E}\left(  \mathbb{D}^{2},\mathbf{\varphi}\right)  =\frac
{\operatorname{Tr}\left\{  HD^{2N}\right\}  }{\operatorname{Tr}\left\{
D^{2N}\right\}  }%
\end{equation}
and the energy of the $2N-1$ electron state $a_{p}^{\dagger}D^{2N}a_{p}$ is
\begin{equation}
\mathcal{E}_{p}\left(  D^{2N},\mathbf{\varphi}\right)  =\frac
{\operatorname{Tr}\left\{  a_{p}Ha_{p}^{\dagger}D^{2N}\right\}  }%
{\operatorname{Tr}\left\{  a_{p}^{\dagger}a_{p}D^{2N}\right\}  }.
\end{equation}
This gives
\begin{equation}
\operatorname{Tr}\left\{  F\left(  \mathbb{D}^{2},\mathbf{\varphi}\right)
\right\}  =\sum_{1\leq p\leq2s}N_{p}\left(  \mathbb{D}^{2},\mathbf{\varphi
}\right)  \left\{  \mathcal{E}\left(  \mathbb{D}^{2},\mathbf{\varphi}\right)
-\mathcal{E}_{p}\left(  \mathbb{D}^{2},\mathbf{\varphi}\right)  \right\}
=\sum_{1\leq p\leq2s}N_{p}\left(  \mathbb{D}^{2},\mathbf{\varphi}\right)
I\left(  \mathbb{D}^{2},\varphi_{p}\right)  , \label{TrF}%
\end{equation}
where the ionization potential associated with the GSO\ $\left\vert
\varphi_{p}\right\rangle $ from the state $D^{2N}$ is
\begin{equation}
I\left(  \mathbb{D}^{2},\varphi_{p}\right)  =\mathcal{E}\left(  \mathbb{D}%
^{2},\mathbf{\varphi}\right)  -\mathcal{E}_{p}\left(  \mathbb{D}%
^{2},\mathbf{\varphi}\right)  .
\end{equation}
Eqs. $\left(  \ref{FE}\right)  $ then give
\begin{subequations}
\begin{align}
\mathcal{E}\left(  \mathbb{D}^{2},\mathbf{\varphi}\right)   &  =\sum_{1\leq
p\leq2s}N_{p}\left(  \mathbb{D}^{2},\mathbf{\varphi}\right)  I\left(
\mathbb{D}^{2},\varphi_{p}\right)  -\mathcal{H}_{2}\left(  \mathbb{D}%
^{2},\mathbf{\varphi}\right) \label{FE1a}\\
\mathcal{E}\left(  \mathbb{D}^{2},\mathbf{\varphi}\right)   &  =\frac{1}%
{2}\left\{  \sum_{1\leq p\leq2s}N_{p}\left(  \mathbb{D}^{2},\mathbf{\varphi
}\right)  I\left(  \mathbb{D}^{2},\varphi_{p}\right)  +\mathcal{H}_{1}\left(
\mathbb{D}^{2},\mathbf{\varphi}\right)  \right\}  . \label{FE1b}%
\end{align}


If $\left\{  \left\vert \varphi_{p}\right\rangle \right\}  $ are the NGSO's of
the state $D^{2N}$ then the one electron energy $\mathcal{H}_{1}\left(
\mathbb{D}^{2},\mathbf{\varphi}\right)  $ can be expressed as
\end{subequations}
\begin{equation}
\mathcal{H}_{1}\left(  \mathbb{D}^{2},\mathbf{\varphi}\right)  =%
%TCIMACRO{\tsum \limits_{1\leq p\leq2s}}%
%BeginExpansion
{\textstyle\sum\limits_{1\leq p\leq2s}}
%EndExpansion
N_{p}\left(  \mathbb{D}^{2},\mathbf{\varphi}\right)  \left\langle \left.
\varphi_{p}\right\vert h_{1}\varphi_{p}\right\rangle
\end{equation}
and we obtain
\begin{equation}
\mathcal{E}\left(  \mathbb{D}^{2},\mathbf{\varphi}\right)  =\frac{1}{2}%
%TCIMACRO{\tsum \limits_{1\leq l\leq2s}}%
%BeginExpansion
{\textstyle\sum\limits_{1\leq l\leq2s}}
%EndExpansion
N_{p}\left(  \mathbb{D}^{2},\mathbf{\varphi}\right)  \left\{  I\left(
\mathbb{D}^{2},\varphi_{p}\right)  +\left\langle \left.  \varphi
_{p}\right\vert h_{1}\varphi_{p}\right\rangle \right\}  \equiv\frac{1}{2}%
%TCIMACRO{\tsum \limits_{1\leq p\leq2s}}%
%BeginExpansion
{\textstyle\sum\limits_{1\leq p\leq2s}}
%EndExpansion
N_{p}\left(  \mathbb{D}^{2},\mathbf{\varphi}\right)  \varsigma\left(
\mathbb{D}^{2},\varphi_{p}\right)  ,
\end{equation}
where
\begin{equation}
\varsigma\left(  \mathbb{D}^{2},\varphi_{p}\right)  =I\left(  \mathbb{D}%
^{2},\varphi_{p}\right)  +\left\langle \left.  \varphi_{p}\right\vert
h_{1}\varphi_{p}\right\rangle .
\end{equation}


One should observe that \emph{only} in the case of an optimized AGP\ state are
the diagonal elements $F\left(  \mathbf{n},\mathbf{\psi}\right)  _{pp}$ of the
Fock operator $F\left(  \mathbf{n},\mathbf{\psi}\right)  $ equal to weighted
ionization potentials $N_{p}\left(  \mathbf{\psi}\right)  I\left(  \psi
_{p}\right)  $ even though the trace of the Fock operator $F\left(
\mathbb{D}^{2},\mathbf{\varphi}\right)  $ for any state can \emph{always }be
expressed as the sum of weighted ionization potentials $N_{p}\left(
\mathbf{\psi}\right)  I\left(  \psi_{p}\right)  $ as displayed in eq. $\left(
\ref{TrF}\right)  $.

\subsection{Fock Identities\label{FockID}}

Using eq. $\left(  \ref{FE1b}\right)  $ we obtain
\begin{subequations}
\begin{equation}
\mathcal{E}\left(  \mathbb{D}^{2},\mathbf{\varphi}\right)  =N\mathcal{E}%
\left(  \mathbb{D}^{2},\mathbf{\varphi}\right)  -\frac{1}{2}\left\{
\sum_{1\leq p\leq2s}N_{p}\left(  \mathbb{D}^{2},\mathbf{\varphi}\right)
\mathcal{E}_{p}\left(  \mathbb{D}^{2},\mathbf{\varphi}\right)  -\mathcal{H}%
_{1}\left(  \mathbb{D}^{2},\mathbf{\varphi}\right)  \right\}  ,
\end{equation}
which leads to
\end{subequations}
\begin{equation}
\mathcal{E}\left(  \mathbb{D}^{2},\mathbf{\varphi}\right)  =\frac{1}{2\left(
N-1\right)  }\left\{  \sum_{1\leq p\leq2s}N_{p}\left(  \mathbb{D}%
^{2},\mathbf{\varphi}\right)  \mathcal{E}_{p}\left(  \mathbb{D}^{2}%
,\mathbf{\varphi}\right)  -\mathcal{H}_{1}\left(  \mathbb{D}^{2}%
,\mathbf{\varphi}\right)  \right\}  . \label{EnFunc}%
\end{equation}
Eq. $\left(  \ref{EnFunc}\right)  $ can also be expressed as
\begin{equation}
\mathcal{E}\left(  \mathbb{D}^{2},\mathbf{\varphi}\right)  =\frac{1}{2\left(
N-1\right)  }\left\{  \sum_{1\leq p\leq2s}N_{p}\left(  \mathbb{D}%
^{2},\mathbf{\varphi}\right)  \left\{  \mathcal{E}_{p}\left(  \mathbb{D}%
^{2},\mathbf{\varphi}\right)  -\left\langle \left.  \varphi_{p}\right\vert
h_{1}\varphi_{p}\right\rangle \right\}  \right\}  ,
\end{equation}
leading to
\begin{equation}
\mathcal{E}\left(  \mathbb{D}^{2},\mathbf{\varphi}\right)  =\frac{1}{2\left(
N-1\right)  }\sum_{1\leq p\leq2s}N_{p}\left(  \mathbb{D}^{2},\mathbf{\varphi
}\right)  \kappa_{p}\left(  \mathbb{D}^{2},\mathbf{\varphi}\right)  ,
\end{equation}
where
\begin{equation}
\kappa_{p}\left(  \mathbb{D}^{2},\mathbf{\varphi}\right)  =\mathcal{E}%
_{p}\left(  \mathbb{D}^{2},\mathbf{\varphi}\right)  -\left\langle \left.
\varphi_{p}\right\vert h_{1}\varphi_{p}\right\rangle .
\end{equation}%
\begin{equation}
\mathcal{E}\left(  \mathbb{D}^{2}\left(  \eta\right)  ,\mathbf{\varphi
}\right)  =\frac{1}{2\left(  N-1\right)  }\left\{  \sum_{1\leq p\leq2s}%
\frac{\operatorname{Tr}\left\{  a_{p}^{\dagger}a_{p}k\left(  \eta\right)
D^{2N}k\left(  \eta\right)  \right\}  }{\operatorname{Tr}\left\{  k\left(
\eta\right)  D^{2N}k\left(  \eta\right)  \right\}  }\left\{  \frac
{\operatorname{Tr}\left\{  a_{p}^{\dagger}Ha_{p}k\left(  \eta\right)
D^{2N}k\left(  \eta\right)  \right\}  }{\operatorname{Tr}\left\{
a_{p}^{\dagger}a_{p}k\left(  \eta\right)  D^{2N}k\left(  \eta\right)
\right\}  }-\left\langle \left.  \varphi_{p}\right\vert h_{1}\varphi
_{p}\right\rangle \right\}  \right\}
\end{equation}%
\begin{equation}
\mathcal{E}\left(  \mathbb{D}^{2},\mathbf{\varphi}\right)  =\frac{1}{2\left(
N-1\right)  }\left\{  \sum_{1\leq p\leq2s}N_{p}\left(  \mathbb{D}%
^{2},\mathbf{\varphi}\right)  \left\{  \frac{\operatorname{Tr}\left\{
a_{p}^{\dagger}Ha_{p}D^{2N}\right\}  }{\operatorname{Tr}\left\{
a_{p}^{\dagger}a_{p}D^{2N}\right\}  }-\left\langle \left.  \varphi
_{p}\right\vert h_{1}\varphi_{p}\right\rangle \right\}  \right\}
\end{equation}


If $D^{2N}$ is an optimized state then in particular we obtain
\begin{equation}
\frac{\delta\mathcal{E}\left(  \mathbb{D}^{2},\mathbf{\varphi}\right)
}{\delta\varphi_{j}^{\ast}}=F\left(  \mathbb{D}^{2},\mathbf{\varphi}\right)
\left\vert \varphi_{j}\right\rangle =\sum_{1\leq k\leq2s}\left\vert
\varphi_{k}\right\rangle \varepsilon_{kj}=h_{1}\left\vert \varphi
_{j}\right\rangle +\sum_{1\leq p\leq2s}N_{p}\frac{\delta\mathcal{E}_{p}\left(
D^{2N},\mathbf{\varphi}\right)  }{\delta\varphi_{j}^{\ast}};\quad1\leq
j\leq2s,
\end{equation}
where
\begin{equation}
\varepsilon_{kj}=\varepsilon_{jk}^{\ast}.
\end{equation}
We can express the derivatives of $2N-1$ electron states $\left\{  a_{p}%
D^{2N}a_{p}^{\dagger}\right\}  $ in terms of a generalized Fock operator%
\begin{equation}
\frac{\delta\mathcal{E}_{p}\left(  D^{2N},\mathbf{\varphi}\right)  }%
{\delta\varphi_{j}^{\ast}}=F\left(  \mathbb{D}^{2}\left(  p\right)
,\mathbf{\varphi}\right)  \left\vert \varphi_{j}\right\rangle ,
\end{equation}
where
\begin{equation}
D^{2}\left(  p\right)  =\frac{1}{\operatorname{Tr}\left\{  a_{p}^{\dagger
}a_{p}D^{2N}\right\}  }L_{N}^{2}\left(  a_{p}D^{2N}a_{p}^{\dagger}\right)  .
\end{equation}
The derivatives wrt occupation number formally give
\begin{equation}
\frac{\partial\mathcal{E}\left(  \mathbb{D}^{2},\mathbf{\varphi}\right)
}{\partial N_{j}}=\frac{1}{2\left(  N-1\right)  }\left\{  \left\{
\mathcal{E}_{p}\left(  D^{2N},\mathbf{\varphi}\right)  +\left\langle \left.
\varphi_{p}\right\vert h_{1}\varphi_{p}\right\rangle \right\}  +\frac
{\partial\mathcal{E}_{p}\left(  D^{2N},\mathbf{\varphi}\right)  }{\partial
N_{j}}\right\}  =0.
\end{equation}
In the case of AGP states occupation number variations can be implemented by
the group $U_{occnum}$ giving
\begin{align}
&  \frac{\partial\mathcal{E}\left(  \mathbb{D}^{2},\mathbf{\varphi}\right)
}{\partial\eta_{j}}=\frac{1}{2\left(  N-1\right)  }\sum_{1\leq p\leq
2s}\left\{  \frac{\partial N_{p}\left(  \mathbf{\eta}\right)  }{\partial
\eta_{j}}\left\{  \mathcal{E}_{p}\left(  k\left(  \mathbf{\eta}\right)
D^{2N}k\left(  \mathbf{\eta}\right)  ,\mathbf{\varphi}\right)  +\left\langle
\left.  \varphi_{p}\right\vert h_{1}\varphi_{p}\right\rangle \right\}  \right.
\\
&  \left.  +\right.  \left.  N_{p}\left(  \mathbf{\eta}\right)  \frac
{\partial\mathcal{E}_{p}\left(  k\left(  \mathbf{\eta}\right)  D^{2N}k\left(
\mathbf{\eta}\right)  ,\mathbf{\varphi}\right)  }{\partial\eta_{j}}\right\}
=0
\end{align}%
\begin{align}
\partial_{\mathbf{\eta}_{j}}^{1}\mathcal{E}\left(  \mathbf{\eta}%
,\mathbf{\lambda},\mathbf{\theta}\right)   &  =\frac{1}{S_{N}\left(
\mathbf{\eta}\right)  }\left\{  \left\langle \left.  g\left(  \mathbf{\eta
},\mathbf{\lambda},\mathbf{\theta}\right)  ^{N}\right\vert \left[  H,\left(
a_{j}^{\dagger}a_{j}+a_{j+s}^{\dagger}a_{j+s}\right)  \right]  _{+}g\left(
\mathbf{\eta},\mathbf{\lambda},\mathbf{\theta}\right)  ^{N}\right\rangle
\right. \nonumber\\
&  -\left.  2\mathcal{E}\left(  \mathbf{\eta},\mathbf{\lambda},\mathbf{\theta
}\right)  \left\langle \left.  g\left(  \mathbf{\eta},\mathbf{\lambda
},\mathbf{\theta}\right)  ^{N}\right\vert \left(  a_{j}^{\dagger}a_{j}%
+a_{j+s}^{\dagger}a_{j+s}\right)  g\left(  \mathbf{\eta},\mathbf{\lambda
},\mathbf{\theta}\right)  ^{N}\right\rangle \right\}  =0;\quad1\leq j\leq s,
\end{align}%
\begin{align}
&  \frac{\partial N_{p}\left(  \mathbf{\eta}\right)  }{\partial\eta_{j}}%
=\frac{\partial}{\partial\eta_{j}}\left\{  \frac{1}{S_{N}\left(  \mathbf{\eta
}\right)  }\left\langle \left.  g\left(  \mathbf{\eta},\mathbf{\lambda
},\mathbf{\theta}\right)  ^{N}\right\vert a_{p}^{\dagger}a_{p}g\left(
\mathbf{\eta},\mathbf{\lambda},\mathbf{\theta}\right)  ^{N}\right\rangle
\right\} \\
&  =\frac{1}{S_{N}\left(  \mathbf{\eta}\right)  }\left\{  \left\langle \left.
g\left(  \mathbf{\eta},\mathbf{\lambda},\mathbf{\theta}\right)  ^{N}%
\right\vert \left\{  a_{j}^{\dagger}a_{j}a_{p}^{\dagger}a_{p}+a_{p}^{\dagger
}a_{p}a_{j}^{\dagger}a_{j}\right\}  g\left(  \mathbf{\eta},\mathbf{\lambda
},\mathbf{\theta}\right)  ^{N}\right\rangle \right. \\
&  \left.  -2N_{p}\left(  \mathbf{\eta}\right)  \left\langle \left.  g\left(
\mathbf{\eta},\mathbf{\lambda},\mathbf{\theta}\right)  ^{N}\right\vert
\left\{  a_{j}^{\dagger}a_{j}\right\}  g\left(  \mathbf{\eta},\mathbf{\lambda
},\mathbf{\theta}\right)  ^{N}\right\rangle \right\}
\end{align}%
\begin{align}
\frac{\partial\mathcal{E}_{p}\left(  k\left(  \mathbf{0}\right)
D^{2N}k\left(  \mathbf{0}\right)  ,\mathbf{\varphi}\right)  }{\partial\eta
_{j}}  &  =\frac{1}{N_{p}}\left\{  \left\langle \left.  g\left(  \mathbf{\eta
},\mathbf{\varphi}\right)  ^{N}\right\vert \left[  a_{p}^{\dagger}%
Ha_{p},\left(  a_{j}^{\dagger}a_{j}+a_{j+s}^{\dagger}a_{j+s}\right)  \right]
_{+}g\left(  \mathbf{\eta},\mathbf{\varphi}\right)  ^{N}\right\rangle \right.
\\
&  -\mathcal{E}_{p}\left(  \left(  \mathbf{\eta},\mathbf{\varphi}\right)
\right)  \left\langle \left.  g\left(  \mathbf{\eta},\mathbf{\varphi}\right)
^{N}\right\vert \left[  a_{p}^{\dagger}a_{p},\left(  a_{j}^{\dagger}%
a_{j}+a_{j+s}^{\dagger}a_{j+s}\right)  \right]  _{+}g\left(  \mathbf{\eta
},\mathbf{\varphi}\right)  ^{N}\right\rangle
\end{align}


\section{Aufbau Principle\label{aufbau}}

The energy of any AGP state can be expressed in terms of weighted ionization
potentials $\left\{  N_{p}\left(  \mathbf{\psi}\right)  I\left(  \psi
_{p}\right)  \right\}  $ and the one electron energy $\mathcal{H}_{1}\left(
\mathbf{n},\mathbf{\psi}\right)  $ as
\begin{equation}
\mathcal{E}\left(  \mathbf{n},\mathbf{\psi}\right)  =\frac{1}{2}\left\{
\sum_{1\leq p\leq2s}N_{p}\left(  \mathbf{\psi}\right)  I\left(  \psi
_{p}\right)  +\mathcal{H}_{1}\left(  \mathbf{n},\mathbf{\psi}\right)
\right\}  .
\end{equation}
Let $\left\vert g^{N}\left(  \mathbf{n},\mathbf{\psi}\right)  \right\rangle $
be a \emph{globally minimized }AGP state such that \emph{\ }
\begin{equation}
I\left(  \psi_{a}\right)  <I\left(  \psi_{b}\right)
\end{equation}
and
\begin{equation}
N_{b}>N_{a}.
\end{equation}
We will compare the energy $\mathcal{E}\left(  \mathbf{n},\mathbf{\psi
}\right)  $ with the energy $\mathcal{E}\left(  \mathbf{n}_{ba},\mathbf{\psi
}\right)  $ produced when the occupation numbers $n_{a}$ and $n_{b}$ are
interchanged, but everything else is left constant. Interchanging $n_{a}$ with
$n_{b}$ interchanges $N_{a}$ and $N_{b}$ as $\mathbf{n}$ is monotonically in
$1-1$ correspondence with $N$. Changing the occupations numbers
\begin{equation}
\left(  n_{a}\leftrightarrow n_{b}\Leftrightarrow N_{a}\leftrightarrow
N_{b}\right)
\end{equation}
is equivalent to changing the CGSO's
\begin{equation}
\psi_{a}\leftrightarrow\psi_{b}\text{ and }\psi_{a+s}\leftrightarrow\psi_{b+s}%
\end{equation}
and can be produced by either a GSO\ or an occupation number transformation
\begin{equation}
\left\vert g^{N}\left(  \mathbf{n},\mathbf{\psi}_{ba}\right)  \right\rangle
=\left\vert g^{N}\left(  \mathbf{n}_{ba},\mathbf{\psi}\right)  \right\rangle
=u\left(  \mathbf{\alpha}_{ba}\right)  \left\vert g^{N}\left(  \mathbf{n}%
,\mathbf{\psi}\right)  \right\rangle =k\left(  \mathbf{\eta}_{ba}\right)
\left\vert g^{N}\left(  \mathbf{n},\mathbf{\psi}\right)  \right\rangle .
\end{equation}
to produce the new energy
\begin{equation}
\mathcal{E}\left(  \mathbf{n}_{ba},\mathbf{\psi}\right)  =\frac{1}{2}%
%TCIMACRO{\tsum \limits_{1\leq p\leq2s}}%
%BeginExpansion
{\textstyle\sum\limits_{1\leq p\leq2s}}
%EndExpansion
\left\{  \frac{1}{S_{N}\left(  \mathbf{n}_{ba},\mathbf{\psi}\right)
}\left\langle \left.  g^{N}\left(  \mathbf{n}_{ba},\mathbf{\psi}\right)
\right\vert \left\{  a_{p}^{\dagger}a_{p}H-a_{p}^{\dagger}Ha_{p}\right\}
g^{N}\left(  \mathbf{n}_{ba},\mathbf{\psi}\right)  \right\rangle
+N_{p}\left\langle \left.  \psi_{p}\right\vert h_{1}\psi_{p}\right\rangle
\right\}  . \label{enba}%
\end{equation}
If $\left\vert g^{N}\left(  \mathbf{n},\mathbf{\psi}_{ba}\right)
\right\rangle $ is also a stationary state then eq. $\left(  \ref{enba}%
\right)  $ becomes
\begin{equation}
\mathcal{E}\left(  \mathbf{n}_{ba},\mathbf{\psi}\right)  =\frac{1}{2}%
\sum_{1\leq p\leq2s}N_{p}I_{ba}\left(  \psi_{p}\right)  +\mathcal{H}%
_{1}\left(  \mathbf{n}_{ba},\mathbf{\psi}\right)  =\frac{1}{2}\sum_{1\leq
p\leq2s}N_{p}I_{ba}\left(  \psi_{p}\right)  +\mathcal{H}_{1}\left(
\mathbf{n},\mathbf{\psi}\right)
\end{equation}
and
\begin{align}
\mathcal{E}\left(  \mathbf{n},\mathbf{\psi}\right)  -\mathcal{E}\left(
\mathbf{n}_{ba},\mathbf{\psi}\right)   &  =N_{a}I\left(  \psi_{a}\right)
+N_{b}I\left(  \psi_{b}\right)  -N_{b}I\left(  \psi_{a}\right)  -N_{a}I\left(
\psi_{b}\right) \nonumber\\
&  =\left(  N_{a}-N_{b}\right)  \left\{  I\left(  \psi_{a}\right)  -I\left(
\psi_{b}\right)  \right\}  >0.
\end{align}
Thus one concludes that
\begin{equation}
\mathcal{E}\left(  \mathbf{n},\mathbf{\psi}\right)  <\mathcal{E}\left(
\mathbf{n}_{ba},\mathbf{\psi}\right)
\end{equation}
which is a contradiction, hence
\begin{equation}
N_{a}>N_{b}.
\end{equation}%
\begin{align}
&  \mathcal{E}\left(  \mathbf{n}_{ba},\mathbf{\psi}\right) \nonumber\\
&  =\frac{1}{2}%
%TCIMACRO{\tsum \limits_{1\leq p\leq2s}}%
%BeginExpansion
{\textstyle\sum\limits_{1\leq p\leq2s}}
%EndExpansion
\left\{  \frac{1}{S_{N}\left(  \mathbf{n}_{ba},\mathbf{\psi}\right)
}\left\langle \left.  g^{N}\left(  \mathbf{n}_{ba},\mathbf{\psi}\right)
\right\vert \left\{  a_{p}^{\dagger}a_{p}H-a_{p}^{\dagger}Ha_{p}\right\}
g^{N}\left(  \mathbf{n}_{ba},\mathbf{\psi}\right)  \right\rangle \right. \\
&  +\left.  N_{p}\left\langle \left.  \psi_{p}\right\vert h_{1}\psi
_{p}\right\rangle \right\} \\
&  =\frac{1}{2}\left\{
%TCIMACRO{\tsum \limits_{1\leq p\leq2s}}%
%BeginExpansion
{\textstyle\sum\limits_{1\leq p\leq2s}}
%EndExpansion
\frac{1}{S_{N}\left(  \mathbf{n},\mathbf{\psi}\right)  }\left\langle \left.
g^{N}\left(  \mathbf{n}_{ba},\mathbf{\psi}\right)  \right\vert \left\{
a_{p}^{\dagger}a_{p}H-a_{p}^{\dagger}Ha_{p}\right\}  g^{N}\left(
\mathbf{n}_{ba},\mathbf{\psi}\right)  \right\rangle +N_{p}\left\langle \left.
\psi_{p}\right\vert h_{1}\psi_{p}\right\rangle \right\}
\end{align}%
\begin{align}
\mathcal{E}\left(  \mathbf{n}_{ba},\mathbf{\psi}\right)   &  =\frac{1}%
{2}\left\{  \frac{2N}{S_{N}\left(  \mathbf{n},\mathbf{\psi}\right)
}\left\langle \left.  g^{N}\left(  \mathbf{n}_{ba},\mathbf{\psi}\right)
\right\vert Hg^{N}\left(  \mathbf{n}_{ba},\mathbf{\psi}\right)  \right\rangle
-%
%TCIMACRO{\tsum \limits_{1\leq p\leq2s}}%
%BeginExpansion
{\textstyle\sum\limits_{1\leq p\leq2s}}
%EndExpansion
\left\langle \left.  g^{N}\left(  \mathbf{n}_{ba},\mathbf{\psi}\right)
\right\vert a_{p}^{\dagger}Ha_{p}g^{N}\left(  \mathbf{n}_{ba},\mathbf{\psi
}\right)  \right\rangle \right\} \\
&  \qquad\qquad\qquad\qquad\qquad\qquad\qquad\qquad\qquad\qquad\qquad
\qquad\qquad+\frac{1}{2}\mathcal{H}_{1}\left(  \mathbf{n},\mathbf{\psi
}\right)
\end{align}%
\begin{equation}
F\left(  \mathbb{D}^{2},\mathbf{\varphi}\right)  _{pp}=\frac{1}%
{\operatorname{Tr}\left\{  D^{2N}\right\}  }\operatorname{Tr}\left\{
\overset{}{\left\{  u^{\dagger}a_{p}^{\dagger}uu^{\dagger}a_{p}uu^{\dagger
}Hu-u^{\dagger}a_{p}uu^{\dagger}Huu^{\dagger}a_{p}^{\dagger}u\right\}  }%
D^{2N}\right\}
\end{equation}%
\begin{equation}
\operatorname{Tr}\left\{  \left(  \mathbb{H}_{1}\left(  \mathbf{\varphi
}\right)  \mathbb{D}^{1}\right)  _{pp}+\frac{1}{2}\left(  L_{2}^{1}\left(
\mathbb{V}_{2}\left(  \mathbf{\varphi}\right)  \mathbb{D}^{2}\right)  \right)
_{pp}\right\}
\end{equation}



\end{document}